% !TEX program = xelatex
% Build: xelatex main; biber main; xelatex main; xelatex main.
\documentclass[lang=cn,chinesefont=nofont,scheme=chinese,11pt,
 color=blue,mode=fancy,device=pad,bibend=biber]{elegantbook}
% Original ebook customization on the unmodified ElegantBook v4.7 class.
% Template: https://github.com/ElegantLaTeX/ElegantBook (LPPL 1.3c+).
\usepackage{mathtools,tabularx,longtable,array,ragged2e}
\usepackage{pgfplots}
\pgfplotsset{compat=1.18}
\usetikzlibrary{arrows.meta,positioning,calc,shapes.geometric,fit,decorations.pathreplacing}
\usepackage{bookmark}
\usepackage[nameinlink,noabbrev]{cleveref}
\geometry{paperwidth=190mm,paperheight=250mm,left=20mm,right=20mm,
 top=21mm,bottom=22mm,headheight=16pt,headsep=7mm,footskip=11mm}
\newif\ifbookwindowsfonts
\bookwindowsfontsfalse
\IfFontExistsTF{SimSun}{
 \IfFontExistsTF{Microsoft YaHei}{
  \IfFontExistsTF{KaiTi}{\bookwindowsfontstrue}{}
 }{}
}{}
\ifbookwindowsfonts
 \setCJKmainfont[BoldFont=Microsoft YaHei,ItalicFont=KaiTi]{SimSun}
 \setCJKsansfont{Microsoft YaHei}
 \setCJKmonofont{Microsoft YaHei}
\else
 \setCJKmainfont[BoldFont=FandolSong-Bold,ItalicFont=FandolKai-Regular]{FandolSong-Regular}
 \setCJKsansfont[BoldFont=FandolHei-Bold]{FandolHei-Regular}
 \setCJKmonofont{FandolFang-Regular}
\fi
\definecolor{ink}{HTML}{203C44}
\definecolor{teal}{HTML}{276A73}
\definecolor{gold}{HTML}{A87934}
\definecolor{paper}{HTML}{FAF8F2}
\colorlet{structurecolor}{ink}
\colorlet{main}{teal}
\colorlet{second}{gold}
\colorlet{third}{teal}
\hypersetup{pdftitle={扩散模型：从分布直觉到多模态生成},
 pdfauthor={扩散模型学习手册},pdfsubject={DDPM、score、SDE、概率流、Flow Matching 与多模态生成},
 pdfkeywords={Diffusion Models,DDPM,Flow Matching,DiT,中文教材},
 linkcolor=teal,citecolor=teal,urlcolor=teal,bookmarksopenlevel=1,
 pdfpagelayout=OneColumn,pdfstartview=FitH}
\setcounter{tocdepth}{1}
\setcounter{secnumdepth}{2}
\setstretch{1.32}
\setlength{\parskip}{3pt plus 1pt}
\setlength{\emergencystretch}{2em}
\raggedbottom
\allowdisplaybreaks[2]
\AtBeginDocument{\setlength{\abovedisplayskip}{9pt plus 2pt minus 2pt}
 \setlength{\belowdisplayskip}{9pt plus 2pt minus 2pt}}
\setlist[itemize]{leftmargin=2em,itemsep=3pt,topsep=5pt}
\setlist[enumerate]{leftmargin=2em,itemsep=3pt,topsep=5pt}
\renewcommand{\arraystretch}{1.2}
\newcolumntype{Y}{>{\RaggedRight\arraybackslash}X}
\fancyhf{}
\fancyhead[L]{\small\color{ink}扩散模型}
\fancyhead[R]{\small\color{ink}\nouppercase{\rightmark}}
\fancyfoot[C]{\small\color{ink}\thepage}
\renewcommand{\headrulewidth}{0.4pt}
\fancypagestyle{plain}{\fancyhf{}\fancyfoot[C]{\small\color{ink}\thepage}
 \renewcommand{\headrulewidth}{0pt}\renewcommand{\headrule}{}}
\tcbset{bookbox/.style={enhanced,breakable,sharp corners,
 boxrule=0pt,borderline west={1.4pt}{0pt}{teal},colback=teal!4,
 colframe=teal,coltitle=ink,colbacktitle=teal!7,
 fonttitle=\sffamily\bfseries,fontupper=\normalfont,
 left=10pt,right=10pt,top=7pt,bottom=7pt,before skip=10pt,after skip=10pt}}
\newtcolorbox{intuitionbox}[1]{bookbox,title={直觉\quad #1}}
\newtcolorbox{keybox}[1]{bookbox,title={要点\quad #1}}
\newtcolorbox{cautionbox}[1]{bookbox,borderline west={1.4pt}{0pt}{gold},colback=gold!5,title={辨析\quad #1}}
\newtcolorbox{examplebox}[2][]{bookbox,title={例子\quad #2},#1}
\newtcolorbox{derivationbox}[1]{bookbox,title={推导\quad #1}}
\newtcolorbox{exercisebox}[1]{bookbox,colback=black!2,title={练习\quad #1}}
\newcommand{\readingguide}[1]{\begin{keybox}{本章阅读路线}#1\end{keybox}}
\newcommand{\chaptertakeaway}[1]{\begin{keybox}{带走的理解}#1\end{keybox}}
\newcommand{\term}[2]{\hyperref[#2]{\textcolor{teal}{#1}}}
\newcommand{\E}{\mathbb{E}}
\newcommand{\R}{\mathbb{R}}
\newcommand{\N}{\mathcal{N}}
\DeclareMathOperator{\KL}{KL}
\newcommand{\dd}{\mathop{}\!\mathrm{d}}
\newcommand{\norm}[1]{\left\lVert#1\right\rVert}
\newcommand{\ip}[2]{\left\langle#1,#2\right\rangle}
\newcommand{\pdata}{p_{\mathrm{data}}}
\newcommand{\pnoise}{p_{\mathrm{noise}}}
\newcommand{\score}{s}
\newcommand{\eps}{\epsilon}
\newcommand{\vtheta}{v_\theta}
\newcommand{\utheta}{u_\theta}
\newcommand{\grad}{\nabla}
\DeclareMathOperator{\divg}{div}
\DeclareMathOperator{\tr}{tr}
\DeclareMathOperator{\Var}{Var}
\DeclareMathOperator{\Cov}{Cov}
\DeclareMathOperator{\SNR}{SNR}
\DeclareMathOperator{\argmin}{arg\,min}
\crefname{chapter}{章}{章}\crefname{section}{节}{节}
\crefname{equation}{式}{式}\crefname{figure}{图}{图}\crefname{table}{表}{表}
\ExecuteBibliographyOptions{sorting=none,maxbibnames=6,maxcitenames=2,doi=false,isbn=false,eprint=false}
\setlength{\bibitemsep}{7pt}
\AtEveryBibitem{\clearfield{urlday}\clearfield{urlmonth}\clearfield{urlyear}}

\addbibresource{references.bib}
\begin{document}
\hypersetup{pageanchor=false}
\begin{titlepage}
\pagecolor{paper}\color{ink}
\begin{tikzpicture}[remember picture,overlay]
 \fill[ink] (current page.north west) rectangle ($(current page.north east)+(0,-7mm)$);
 \draw[gold,line width=0.7pt] ($(current page.north west)+(20mm,-34mm)$) -- ($(current page.north east)+(-20mm,-34mm)$);
 \begin{scope}[shift={($(current page.center)+(0,-24mm)$)},scale=1.1]
  \foreach \i in {0,...,35}{
   \pgfmathsetmacro{\px}{-5+mod(\i*17,31)/31*3}
   \pgfmathsetmacro{\py}{-1.5+mod(\i*11,37)/37*3}
   \fill[teal!45] (\px,\py) circle (1.3pt);
  }
  \foreach \i in {0,...,20}{
   \pgfmathsetmacro{\ang}{\i*17}
   \draw[teal!25,thin] ({-2+0.07*\i},{-1.1+0.1*mod(\i*7,21)})
    .. controls (0,{-1+0.1*\i}) and (1.3,{1.3*sin(\ang)}) .. ({3.3+0.85*cos(\ang)},{1.1*sin(\ang)});
   \fill[teal] ({3.3+0.85*cos(\ang)},{1.1*sin(\ang)}) circle (1.5pt);
  }
  \node[font=\small,text=ink] at (-3.9,-2.1) {简单的噪声};
  \node[font=\small,text=ink] at (3.3,-2.1) {有结构的分布};
 \end{scope}
\end{tikzpicture}
\vspace*{38mm}
{\sffamily\fontsize{34}{42}\selectfont\bfseries 扩散模型\par}
\vspace{8mm}
{\sffamily\Large 从分布直觉到多模态生成\par}
\vspace{6mm}
{\large Diffusion Models\par}
\vfill
{\color{teal}\rule{28mm}{1pt}\par}
\vspace{5mm}
{\large 概念\quad 数学\quad 方法\quad 应用\par}
\vspace{3mm}
{\small 面向熟悉数学分析、线性代数与机器学习的读者\par}
\vspace{12mm}
{\small 2026 年 10 月 9 日\quad 第一版\par}
\end{titlepage}
\nopagecolor\color{black}
\hypersetup{pageanchor=true}
\frontmatter
\chapter*{写在前面}
\addcontentsline{toc}{chapter}{写在前面}
这本书想回答一个看似朴素的问题：怎样让一团随机数逐渐成为一幅有意义的图像？答案并不是把某张原图藏进噪声，再把它找出来。我们要学习一个分布，以及把简单分布变成复杂分布的规律。图像只是最直观的入口；相同思想还可以用于视频、音频、三维结构和其他数据。

本书假定你知道导数、积分、矩阵、条件概率和神经网络的基本训练方式。但“知道数学”与“熟悉某篇论文的语言”是两回事。我们会解释重要术语，交代每个公式为何出现，再把它们连接到可实现的方法。随机微积分不是阅读前提；第二章需要的内容会从有限步直觉开始，并链接到附录。

\begin{keybox}{四章是一条连续的理解路线}
第一章建立分布、几何和去噪的直觉。第二章沿着可计算的高斯过程，推导 DDPM、score、反向 SDE、概率流 ODE 和速度场。第三章区分概率路径、训练目标、求解器与网络，解释主要方法为何演进。第四章把数学接回真实系统，从 U-Net 与隐空间走到 DiT、视频和统一多模态生成。
\end{keybox}

你可以顺序阅读，也可以先读第一章、第三章的算法框架和第四章的网络图，再回到第二章补齐证明。若只想掌握核心数学，重点读第二章的前向边缘、ELBO、去噪 score matching、概率流和条件流匹配。附录中的术语表是一个可随时跳转的工具；正文的青绿色术语、章节编号、公式编号和参考文献均可点击。大多数 PDF 阅读器可以用“返回上一视图”回到跳转前的位置。

为了减少不必要的困难，全书把连续扩散的时间 $t$ 固定为“数据到噪声”，把流匹配的时间 $\tau$ 固定为“噪声到数据”。这个约定与部分论文不同。我们宁可反复标注时间方向，也不让一个负号把整个理解颠倒。第二章篇幅较长，是因为这里既要展示答案，也要展示答案为什么成立。

本书不是论文摘要的拼接。历史事实和特定方法归属引向原始论文；例子、推导的组织方式与教学图为本书撰写。我们会把等式所需的假设、理想极限与实际误差分开：理论上的精确 score 不等于训练出的神经网络，连续时间的同边缘分布不等于有限步采样完全相同，条件插值直线也不意味着学到的整体轨迹一定是直线。

第四章的文献核查截至 2026 年 10 月 9 日。它选取对理解技术路线有帮助的代表，而不声称覆盖所有当日新工作。论文、预印本、官方技术报告和产品披露具有不同的信息边界；未公开的架构不作猜测。未来方向是从已知瓶颈出发的判断，不能当成已经实现的能力。

版式基于 ElegantBook v4.7，保留其原始类文件和 LPPL 1.3c 或更新版本的许可证；本书另行定制了电子书页幅、墨青色章节、留白和可跨页的教学框。模板来源：\href{https://github.com/ElegantLaTeX/ElegantBook}{ElegantLaTeX/ElegantBook}。所有论文链接集中在参考文献，可从正文直接进入。
\tableofcontents
\chapter*{符号与时间速查}
\addcontentsline{toc}{chapter}{符号与时间速查}
\begin{longtable}{p{30mm}p{108mm}}
\toprule 符号 & 本书含义\\\midrule\endhead
$x_0, X$ & 数据样本；大写常用于随机变量或流匹配端点。\\
$x_k,\ K$ & 离散 DDPM 的第 $k$ 个含噪样本及总步数。\\
$\beta_k,\alpha_k$ & 单步噪声方差，$\alpha_k=1-\beta_k$。\\
$\bar\alpha_k$ & 累计信号保留量 $\prod_{j=1}^k\alpha_j$。\\
$t\in[0,1]$ & 连续扩散：$t=0$ 为数据，$t=1$ 为高噪声端；生成向时间减小的方向积分。\\
$a(t),b(t)$ & 表示 $x_t=a(t)x_0+b(t)\epsilon$ 的信号、噪声系数。\\
$\epsilon,Z$ & 标准高斯噪声；流匹配用 $Z$ 表示起点。\\
$p_t(x),s_t(x)$ & 时刻 $t$ 的边缘密度与 score，$s_t=\nabla_x\log p_t$。\\
$f(x,t),g(t)$ & SDE 的漂移与各向同性扩散强度。\\
$v_t^{\mathrm{PF}}$ & 扩散时间下的概率流 ODE 速度。\\
$v=a\epsilon-bx_0$ & VP 路径常用的网络预测目标，须与实际时间导数区分。\\
$\tau\in[0,1]$ & 流匹配：$\tau=0$ 为噪声，$\tau=1$ 为数据。\\
$Y_\tau,u_\tau$ & 流匹配的插值随机变量与边缘速度。\\
$c,\theta,d$ & 条件信息、神经网络参数与数据向量维数。\\
$q,p_\theta$ & DDPM 已知前向过程与学得的反向生成过程。\\
$\E,\KL,\norm{x}$ & 期望、KL 散度、欧氏范数。\\
\bottomrule
\end{longtable}
\mainmatter
\IfFileExists{chapters/01-concepts.tex}{\chapter{生成的本质：从图像到分布}
\label{ch:concepts}

\readingguide{本章先建立一个能够反复使用的直觉：图像生成要学会的，不是一张代表性图片，而是产生许多合理图片的规律。我们会依次讨论分布、高维几何、加噪与去噪、隐空间和条件生成，再把扩散放回生成模型的全景中。读完后，你应当能够解释为什么一个预测噪声的网络可以参与图像生成，也能识别这种解释的边界。涉及的概率术语可点击跳转到附录；完整推导将在第\ref{ch:math}章展开。}

\section{先问一个看似简单的问题：什么叫会生成？}
\label{sec:c-generation}

假设我们给模型一百万张猫的照片，然后要求它生成一张新照片。最容易想到的方案是输出这批照片的平均值。这个方案在某种意义下十分合理：如果使用像素平方误差，且模型只能输出一张固定图片，平均图像就是最优答案。但它往往是一团模糊的轮廓。猫的头可以向左，也可以向右；眼睛可以睁开，也可以闭上；背景可以是窗台，也可以是草地。这些可能性逐像素平均之后，不再构成其中任何一个合理场景。

问题出在任务的描述。我们希望模型生成一个可能发生的结果，而不是在所有可能结果之间折中。一张照片只能呈现一个视角、一种姿态和一种光照；一套生成规则却应当容纳许多视角、姿态和光照。描述这套规则最自然的数学对象，就是\term{概率分布}{gl:distribution}。

用 $X$ 表示随机图像，用 $x$ 表示它的一次具体取值。即使我们最终只想得到一张图片，也应当把目标写成
\[
  \widehat X\sim p_\theta,\qquad p_\theta\approx \pdata,
\]
其中 $\pdata$ 是我们想要建模的数据分布，$p_\theta$ 是参数为 $\theta$ 的模型产生的分布。符号 $\sim$ 读作“服从”；它把一次生成和背后的重复试验联系起来。我们真正关心的是：如果多次运行模型，出现的图片是否合理，是否有足够的变化，各类内容出现的比例是否符合目标。

\begin{examplebox}[unbreakable]{平均答案与随机答案}
一个数据集只有两个标量值 $-2$ 和 $2$，各占一半。输出 $0$ 的模型实现了最小的常数平方误差，却从不输出数据中可能出现的值。另一个模型先掷一枚公平硬币，再输出 $-2$ 或 $2$，才正确复现了这个分布。图像中的“不同姿势”与这个例子的两个端点相似，只是维数高得多。
\end{examplebox}

由此可以区分三个经常混在一起的目标。重建要求输出尽量接近某张给定图片；预测要求根据条件判断未来或缺失的内容；生成要求从一组可能性中抽取一个结果。它们可以共享网络，但评价标准不同。重建的答案经常可以与原图逐像素比较，无条件生成的答案没有这样一个唯一参照。把生成当作单一答案回归，模糊便成为一种可以预期的结果。

“复现分布”也不意味着逐张复制训练集。如果训练集为 $\{x^{(1)},\ldots,x^{(n)}\}$，它对应的经验分布给每张训练图分配质量 $1/n$。随机抽取并原样输出一张训练图，可以完美复现这个经验分布，却没有学会在未见情形下生成。实际学习希望利用有限样本、模型结构和训练目标，推断更丰富的总体规律。这就是生成模型里的泛化问题。训练样本本身有偏，学到的规律也可能偏；网络容量足够大时，它同样可能记忆样本。扩散并没有天然免除这些问题。

还应当明确“真实”的含义。互联网上收集的图片分布不是自然界全部可能图像的分布，更不是人类审美的统一标准。相机、网页选择、标注方式和筛选过程都在改变它。若数据集中人物通常居中，模型就可能把“居中”当作高概率规律。我们在全书使用 $\pdata$，指的是指定任务与数据获取过程下的目标分布；它不是一个脱离任务而客观存在的完美图库。

\begin{intuitionbox}{生成能力是重复运行的性质}
单张图片可以漂亮，但一个模型是否学得好，要看它产生的整批结果。只输出一张漂亮图片的模型，可能缺少覆盖；能产生很多不同图片的模型，也可能经常输出不合理内容。合理性、覆盖范围和概率分配是彼此相关、又不能完全互相替代的要求。
\end{intuitionbox}

\section{分布、密度和采样：三个不同的问题}
\label{sec:c-distribution}

\subsection{分布描述概率，密度描述概率如何铺开}

一张 $H\times W$ 的 RGB 图像可表示为 $d=3HW$ 维向量。对于八位像素，严格说来它属于一个巨大的离散集合；若把像素归一化并作为实数处理，便常写 $x\in\R^d$。这两种建模方式有不同的概率含义，不能在计算\term{似然}{gl:likelihood}时随意交换。连续模型处理离散图像时，还需要对量化区间或去量化作出约定。

分布 $P$ 为事件分配概率，例如“图像里有两只猫”或“某个像素位于某个区间”。当分布有\term{概率密度}{gl:density} $p(x)$ 时，
\[
  \mathbb P(X\in A)=\int_A p(x)\,\dd x.
\]
因此 $p(x)$ 本身不是“恰好生成图片 $x$ 的概率”。对通常的连续分布，单点概率为零；一个小区域的概率，才近似等于密度乘以区域体积。密度也可以大于一，因为它有“每单位体积”的含义。只要积分为一，就是合法的密度。

这一区分在高维空间尤其关键。密度比较的是单位体积中的集中程度，采样频率比较的是整个区域累积的概率质量。面积很小的高密度区域，可能比面积巨大的中等密度区域更少被抽到。若把生成理解为“寻找 $p(x)$ 最大的图片”，就会丢掉这种体积效应，也会把多样性缩成一个优化问题。

\begin{examplebox}{最密的点，未必是最典型的样本}
设 $Z\sim\N(0,I_d)$。密度在原点最大，但
\[
  \E\norm{Z}^2=d,\qquad
  \operatorname{Var}(\norm{Z}^2)=2d.
\]
由切比雪夫不等式，对任意 $\delta>0$，
\[
 \mathbb P\!\left(\left|\frac{\norm{Z}^2}{d}-1\right|>\delta\right)
 \leq \frac{2}{d\delta^2}.
\]
维数很大时，大部分样本的半径接近 $\sqrt d$，而不是接近零。原点密度高，原点附近足够小的球却几乎没有总体概率质量。
\end{examplebox}

这个例子不只是一个概率练习。扩散模型常从高维高斯噪声出发；所谓“随机噪声的常见样子”，就是许多坐标同时有中等幅度的随机波动，而不是所有坐标接近零。类似地，一张图片在模型下具有较高密度，未必就具有更好的视觉质量。密度还依赖坐标和量化约定，视觉质量则依赖结构、语义和感知。它们需要联系起来研究，不能当作同一个指标。

\subsection{知道密度和会采样并不等价}

假设有人给你函数 $p(x)$，你能计算任意图片的密度。你仍然不一定能高效地产生它的样本。高维空间里随机乱找，几乎不会碰到合理图片；寻找最大值，又只是在做优化。反过来，一个网络可以把随机向量直接变成图片，因而很容易采样，但未必能给这张图片计算归一化密度。生成模型的不同家族，往往就在这两个能力之间作不同安排。

对确定性生成器，可以写
\[
 Z\sim\pnoise,\qquad \widehat X=G_\theta(Z).
\]
生成器把输入分布“推送”为输出分布。这里的数学含义是
\[
 \mathbb P(\widehat X\in A)
 =\mathbb P\!\left(Z\in G_\theta^{-1}(A)\right).
\]
这里 $G_\theta^{-1}(A)$ 表示所有会被映射进 $A$ 的输入集合，称为原像；它不要求 $G_\theta$ 存在逆函数。即使 $G_\theta$ 的每一步完全确定，只要 $Z$ 是随机的，输出仍然是随机变量。一个固定函数可以产生无限多的结果，并不需要为每个结果单独存储一个模板。

这一观察能提前解开一个常见困惑：扩散采样一定要在每一步加入随机噪声吗？答案是否定的。某些采样方式只在初始状态随机，随后沿确定性轨迹运动；另一些方式在中途继续注入随机性。只要最终的分布正确，两种方式都可以进行生成。第\ref{ch:math}章将解释随机微分方程与概率流常微分方程为什么能在适当条件下拥有相同的时刻分布。

分布还可通过样本来认识，而不必先知道一个可计算的公式。训练集提供的就是这种信息：我们会从 $\pdata$ 取样，却不知道 $\pdata(x)$ 的解析表达式。有效的训练目标必须适应这个现实。扩散之所以吸引人，一个重要原因是：它能利用已知的加噪规则，把“只有干净样本”转化成大量有监督的预测任务。

\subsection{边缘、条件和路径分别在说什么}

假设 $X_t$ 是某个时间的含噪图像。$p_t(x)$ 描述这一时刻的一整群样本，这是时刻\term{边缘分布}{gl:distribution}。$q(x_t\mid x_0)$ 描述给定一张原图之后，加噪的结果怎样变化，这是\term{条件分布}{gl:conditional}。$q(x_0,x_1,\ldots,x_K)$ 则描述一条完整随机轨迹的联合规律。边缘、条件和联合分布不应混为一谈。

仅知道每一个时刻的边缘分布，不会唯一决定样本之间怎样连线。想象两张纸上各有一百个点，第一张代表初始噪声，第二张代表最终数据。你可以按不同规则把两张纸上的点配对，最终点云完全相同，运动路线却不同。第\ref{ch:methods}章中不同的采样过程和第\ref{sec:flowmatching}节中的路径设计，都将使用这一区分。

\section{为什么图像分布这么难学？}
\label{sec:c-geometry}

\subsection{维数高，更困难的是结构强}

一张 $256\times256$ 的 RGB 图像就有 $196\,608$ 个坐标。若每个坐标独立取一个像素值，所得图片几乎总是杂乱的噪声。自然图像的像素之间存在大量依赖：物体边缘连续，光照在空间上变化有规律，人的左右肢体受共同姿态约束，文字必须由一致的笔画组成。困难不只是坐标多，而是需要把这些约束同时满足。

“局部看起来合理”也不足以保证全局合理。每一小块都像皮毛，整张图却可能有三条尾巴；每个字母都有清楚的线条，一行文字却可能没有拼成任何词。反过来，全局构图正确，细部纹理也可能不可信。有效模型需要在不同尺度上协调信息。这为后文的 U-Net、多尺度特征以及注意力机制提供了问题背景，但概率原理并不规定必须使用某一种网络。

有限数据给学习增加了另一层困难。设想给每个坐标划十个格子，再统计每个格子组合的出现频率，便有 $10^d$ 种组合；现实数据远远无法覆盖。神经网络并非消除了这个指数级问题，而是用共享参数和结构偏好压缩了需要学习的规律。某种边缘检测模式可以在许多位置复用；某种“物体与阴影相协调”的关系，可以跨越许多具体图片迁移。

\subsection{流形是一幅有用的地图，但不是字面的全部真相}

一个常见想法是，自然图像集中在高维空间中的较低维结构附近。相机位置、物体姿态、光照和材料等因素的变化，可以在像素空间画出连续轨迹。我们用\term{数据流形}{gl:manifold}作为理解这种几何结构的语言：周围空间很大，而有意义的变化只占其中某些方向。

例如，固定一只手写数字的风格，只连续改变它的倾角，图片会沿一条曲线变化。改变倾角和笔画粗细，可能形成一个局部二维曲面。随机改变每一个像素，则大多沿着离开这个曲面的方向运动。对这个理想例子，去噪可以理解为把偏离曲面的观测拉回附近，同时保留曲面内部的有效变化。

不过，真实图像还包含传感器噪声、压缩误差、遮挡、离散类别和复杂组合。不同类别可能对应不同区域；某些变化会产生不连续的拓扑变化。不能因此断言所有图像恰好位于一张光滑、低维、无边界的流形上。“靠近结构化集合”往往比“严格落在一个流形上”更符合实际。后面有些数学公式需要完整空间中的平滑密度，正是加噪帮助我们获得了这种良好的对象。\cite{r_ncsn}

这里还能区分几何和概率。知道哪里有合理图片，只告诉了我们分布的支持或高概率区域；还要知道各区域应分到多少概率质量。两张同样合理的图片，可以在数据中拥有完全不同的出现频率。把样本投影到一个结构化集合，并不自动实现正确的生成分布。

\begin{figure}[htbp]
\centering
\begin{tikzpicture}[x=0.85cm,y=0.85cm]
  \draw[->,gray] (-0.4,0) -- (6.2,0) node[right] {$x^{(1)}$};
  \draw[->,gray] (0,-0.4) -- (0,3.6) node[above] {$x^{(2)}$};
  \draw[line width=1.4pt,blue!70!black]
    plot[smooth] coordinates {(0.45,0.7) (1.1,1.7) (2,2.15) (3,1.55) (4,1.45) (5.4,2.8)};
  \foreach \px/\py in {0.55/0.86,1.12/1.74,1.8/2.14,2.5/1.94,3.2/1.5,4.25/1.63,5.18/2.59}
    \fill[blue!70!black] (\px,\py) circle (2pt);
  \foreach \px/\py/\qx/\qy in {1.05/0.65/1.3/1.88,2.0/2.98/2.0/2.15,3.2/2.55/3.25/1.5,4.75/0.8/4.45/1.8}
    {\fill[orange!85!black] (\px,\py) circle (2.1pt);
     \draw[->,orange!85!black,line width=0.75pt] (\px,\py) -- (\qx,\qy);}
  \node[align=left,font=\small] at (4.5,3.25) {蓝线：结构化数据区域\\橙点：受到扰动的观测};
  \node[font=\small,align=center] at (3,-0.8) {箭头仅示意局部纠正方向；\\沿蓝线的概率分配仍然需要学习。};
\end{tikzpicture}
\caption{二维图示帮助理解高维几何。真实图像通常更复杂，箭头也不必是欧氏最近点投影。图为原创教学示意。}
\label{fig:c-manifold}
\end{figure}

\subsection{没有一条坐标轴天然叫“猫的姿势”}

像素坐标里，改变姿势会同时改变许多位置的颜色；改动一个像素通常不能完成一个清楚的语义变化。这种坐标与语义之间的不匹配，是高维学习的又一个难点。学习到的特征可以让某些规律更容易表达，但也不保证每个特征维度只控制一个独立概念。

因此我们应当谨慎对待“模型学到了语义空间”这句话。它可能表示模型的内部表示对语义有用，也可能只是表示相似内容在某种特征距离下比较接近。它并不自动意味着坐标可解释、变量互相独立、任意插值都合理。即使是一个效果很好的生成模型，仍可能把姿势、背景和风格纠缠在一起。

\section{加噪：先把复杂问题变成一串较温和的问题}
\label{sec:c-noising}

\subsection{扩散不是把一张图均匀地涂成灰色}

扩散模型的基本安排是：规定一个把数据逐渐扰乱的过程，然后学习怎样反向推进。这个想法在早期扩散概率模型中已经出现；DDPM 则给出了影响深远的训练与采样形式。\cite{c_sohl2015diffusion,r_ddpm}

先不讨论具体时间步，考虑
\begin{equation}
  X_t=a(t)X_0+b(t)\eps,\qquad
  X_0\sim\pdata,\quad \eps\sim\N(0,I_d),
  \label{eq:c-noising}
\end{equation}
其中 $X_0$ 与 $\eps$ 独立。$a(t)$ 控制保留多少原信号，$b(t)$ 控制注入多少噪声。在常见设置里，$t$ 从 $0$ 向 $1$ 增加，信号减弱、噪声增强。本书约定生成过程沿相反方向，从 $t=1$ 推进到 $t=0$。不同方法的系数和路径可以不同，不能只凭这一个写法判定整个过程。

对于固定原图 $x_0$，$X_t$ 的条件分布是
\[
 q_t(x\mid x_0)=\N\!\left(x;\,a(t)x_0,b(t)^2I_d\right).
\]
这句话很容易理解：每张原图周围都有一团高斯云。对所有原图取平均，便得到整个含噪数据的边缘分布
\begin{equation}
 p_t(x)=\int q_t(x\mid x_0)\,\pdata(x_0)\,\dd x_0,
 \label{eq:c-mixture}
\end{equation}
这里为便于书写，暂且假设 $\pdata$ 有密度；对离散或奇异分布，应把积分写为对概率测度的积分，含义相同。加噪改变的是整群样本的分布，而不是只对一张图片做一个视觉滤镜。

如果直接输出平均图像，不同对象会叠在同一张图里；高斯加噪则使每张原图产生不同的随机观测。前者把不确定性压成一个均值，后者用分布表达不确定性。图像最终看起来像噪声，是随机扰动遮蔽了结构，而不是所有原图都收缩成同一个灰点。

\subsection{高斯平滑为什么有帮助}

为看清加噪的作用，暂设 $a(t)=1$、$b(t)=\sigma$。此时 $p_\sigma$ 是原分布与高斯核的卷积。原来尖锐、分离或者没有普通密度的结构，经过正方差的高斯扰动后，会形成一个平滑、处处为正的密度。这让“在某个观测附近，概率怎样变化”有了良好定义。

小噪声保留细部差异，每一簇数据附近仍有自己的局部结构；大噪声使相邻簇重叠，细节逐渐难以辨认。对于两个分开的数据簇，先研究较宽、较平滑的概率地形，再逐步增加分辨率，通常比从远离数据的任意点直接寻找细小簇更有希望。这个多尺度安排正是理解扩散的重要入口，但它不是对有限模型必然成功的证明。

\begin{figure}[htbp]
\centering
\begin{minipage}{0.48\linewidth}
\centering
\begin{tikzpicture}
\begin{axis}[width=\linewidth,height=4.5cm,domain=-4.5:4.5,samples=120,
 xmin=-4.5,xmax=4.5,ymin=0,ymax=0.55,axis lines=left,
 xlabel={$x$},ylabel={$p_\sigma(x)$},title={较小噪声 $\sigma=0.4$},
 tick label style={font=\footnotesize},label style={font=\small},
 title style={font=\small}]
\addplot[blue!70!black,line width=1.2pt]
 {0.5/(sqrt(2*pi)*0.4)*(exp(-(x+2)^2/(2*0.4^2))+exp(-(x-2)^2/(2*0.4^2)))};
\end{axis}
\end{tikzpicture}
\end{minipage}\hfill
\begin{minipage}{0.48\linewidth}
\centering
\begin{tikzpicture}
\begin{axis}[width=\linewidth,height=4.5cm,domain=-6:6,samples=120,
 xmin=-6,xmax=6,ymin=0,ymax=0.55,axis lines=left,
 xlabel={$x$},ylabel={$p_\sigma(x)$},title={较大噪声 $\sigma=2.2$},
 tick label style={font=\footnotesize},label style={font=\small},
 title style={font=\small}]
\addplot[orange!85!black,line width=1.2pt]
 {0.5/(sqrt(2*pi)*2.2)*(exp(-(x+2)^2/(2*2.2^2))+exp(-(x-2)^2/(2*2.2^2)))};
\end{axis}
\end{tikzpicture}
\end{minipage}
\caption{对等概率的两个点 $-2$ 与 $2$ 添加不同强度的高斯噪声。右图的密度更平滑，两个来源更难区分。本图改变了噪声标准差，未进行信号缩放；因此它说明平滑作用，而不是完整的方差保持扩散过程。}
\label{fig:c-smoothing}
\end{figure}

回到式\eqref{eq:c-noising}。当 $a(t)$ 足够小、$b(t)$ 接近 $1$ 时，$X_t$ 的分布接近标准高斯。这里的“接近”有必要保留：有限时间的 DDPM 通常只有近似高斯的末端分布。若 $a(1)=0$ 且 $b(1)=1$，则末端才精确地独立于 $X_0$。生成前为初始随机状态指定的分布称为生成先验；选择一个容易采样、又尽量接近末端的先验，是生成链的第一环。

从矩的角度看，
\[
 \E X_t=a(t)\E X_0,\qquad
 \operatorname{Cov}(X_t)
 =a(t)^2\operatorname{Cov}(X_0)+b(t)^2I_d,
\]
假设相应矩存在。信号系数趋于零时，数据的均值与协方差影响逐渐减弱。仅仅矩接近还不等于分布完全接近，但这能解释为什么标准高斯是一个合理的终点设计。

\subsection{失去信息之后，反向怎么可能？}

如果加噪逐渐销毁信息，从纯噪声出发怎么可能还原图片？困惑来自“还原”的两个含义。恢复加噪前那一张特定原图，需要观测中保留足够信息；产生一个符合原分布的新样本，却只需要知道怎样对所有可能的原图进行选择。扩散生成通常做的是后者。

设观察到一张模糊而含噪的图。它可能来自窗台上的猫，也可能来自沙发上的猫。一个生成模型无须知道历史上到底是哪张图片被扰乱，而要知道哪些解释合理、各自有多少概率。每次反向采样可选出一个不同解释，重复运行后再现总体规律。信息没有从噪声中凭空出现；关于图像结构的信息来自训练数据和已经学到的参数，随机性负责选择具体结果。

这也解释了为何加噪过程可以是已知的，而反向过程需要学习。给原图加多少噪声，我们可以控制；给定噪声观测，哪种结构更合理，则需要了解数据分布。若把 $\pdata$ 换成手写数字、医学切片或者乐谱，正向高斯规则仍可相同，反向规律却完全不同。

在贝叶斯语言中，尚未考虑当前观测时的规律称为先验；结合观测之后，对未知原图的条件分布称为后验。这里 $\pdata(x_0)$ 提供原图先验，$q_t(x_t\mid x_0)$ 描述观测如何产生，二者通过贝叶斯公式得到后验 $p(x_0\mid x_t)$。后文讨论的后验均值和后验采样，都是针对这份经过更新的概率规律。它与生成链末端的噪声先验承担不同角色。

\section{去噪为什么会携带分布的信息？}
\label{sec:c-denoising}

\subsection{训练标签可知，不代表推理时答案可知}

训练时，我们先取一张干净图 $x_0$，再自行产生 $\eps$，计算 $x_t=a(t)x_0+b(t)\eps$。因此 $x_0$ 和 $\eps$ 都是现成标签。可用网络 $D_\theta(x_t,t)$ 预测干净图，也可用 $\eps_\theta(x_t,t)$ 预测噪声。表面上看，这只是一个普通回归任务。

但是，推理时只给网络 $x_t$ 和 $t$，不给它产生这张观测的 $x_0$ 和 $\eps$。同一个 $x_t$ 可以由不同的原图与噪声组合得到，所以网络不可能一般性地精确识别这一次所加的噪声。平方误差下，它学到的理想函数是条件均值：
\[
 D^*(x,t)=\E[X_0\mid X_t=x],\qquad
 \eps^*(x,t)=\E[\eps\mid X_t=x].
\]
这不是失败，而是回归目标的正确统计含义。给定二者之一，在 $a(t),b(t)$ 非零时，可以换算另一个：
\[
 D^*(x,t)=\frac{x-b(t)\eps^*(x,t)}{a(t)}.
\]

条件均值依赖 $\pdata$。若训练数据只有水平线，去噪器会倾向解释出水平线；若数据只有圆圈，同样的模糊输入会得到不同预测。因此回归任务并非只在学一种通用滤波器，它还在学习“这种观测通常由哪些数据产生”。

\subsection{一个能算清楚的两点世界}

继续使用 $X_0\in\{-h,h\}$、各占一半的例子，设
\[
 Y=X_0+\sigma\eps,\qquad \eps\sim\N(0,1),\quad h>0.
\]
其中 $X_0$ 与 $\eps$ 独立。观察 $Y=y$ 后，对两个来源的相对支持为
\[
 \frac{\mathbb P(X_0=h\mid Y=y)}
      {\mathbb P(X_0=-h\mid Y=y)}
 =\exp\!\left(\frac{2hy}{\sigma^2}\right).
\]
因此最优平方误差去噪器为
\begin{equation}
 D^*(y)=h\tanh\!\left(\frac{hy}{\sigma^2}\right).
 \label{eq:c-two-point-denoiser}
\end{equation}
正的观测支持右边端点，负的观测支持左边端点；证据越强，预测越接近对应端点。当 $\sigma$ 较大或 $y$ 接近零时，两种解释难分，均值就停留在它们之间。

这个例子同时说明了去噪器的能力和局限。$D^*(0)=0$，而原始数据从不取 $0$。即使一个去噪器统计上完全正确，一次输出的条件均值也不必是数据分布中的合法样本。高噪声时“预测的干净图很模糊”，并不直接说明模型没学好；可能只是后验不确定性很大。问题是如何利用这个预测构造正确的反向过程，而不是期待一个均值估计器一次完成全部生成。

对相同含噪观测，后验采样可以返回 $-h$ 或 $h$，后验均值则返回它们的加权平均。两种运算回答的是不同问题。扩散采样会通过一系列更新逐步降低不确定性，不能简单等同于对每一步都“直接输出预测原图”。具体更新式由时间路径、反向分布和采样器共同决定。

\subsection{去噪残差是一种概率方向}

令 $p_\sigma(y)$ 是 $Y=X_0+\sigma\eps$ 的密度。一个关键恒等式是
\begin{equation}
 \grad_y\log p_\sigma(y)
 =\frac{D^*(y)-y}{\sigma^2}.
 \label{eq:c-tweedie-intuition}
\end{equation}
左边称为\term{分数函数}{gl:score}，通常简称 score。这个词与“给图片打一个分数”无关：它是对数密度关于观测坐标的梯度，是一个与图片同维度的向量。式\eqref{eq:c-tweedie-intuition}告诉我们，去噪器建议的修正方向包含了含噪分布的局部信息。去噪与 score matching 的联系是扩散的重要思想来源。\cite{c_vincent2011connection}

读者熟悉梯度就可以理解其局部含义。在一个很小的位移 $\Delta y$ 下，
\[
 \log p_\sigma(y+\Delta y)-\log p_\sigma(y)
 \approx \ip{\grad\log p_\sigma(y)}{\Delta y}.
\]
score 指向密度局部增加最快的方向，同时给出变化的大小。它告诉运动机制在每个位置应该怎样响应，而不需要直接输出一个巨大的密度函数。完整恒等式及不同噪声系数下的版本将在第\ref{sec:score}节推导。

仍以两点世界为例，代入式\eqref{eq:c-two-point-denoiser}就有
\[
 \score_\sigma(y)
 =\frac{h\tanh(hy/\sigma^2)-y}{\sigma^2}.
\]
小噪声时，在两个端点附近，箭头把偏离簇中心的点拉回；在两簇之间，它又根据观测位置区分两种可能性。噪声尺度改变，整个向量场也改变。网络接收 $t$，就是为了知道现在应当使用哪一张“概率地图”。

\begin{cautionbox}{沿梯度上山不是分布采样}
若只反复沿 score 做梯度上升，样本通常会向密度极大值集中，不能保证保持正确的概率比例。生成需要一个使整群样本按照目标分布演化的机制。随机项、扩散系数、时间相关的漂移，以及确定性流中的体积变化，都会参与这个机制；它不是普通的图片质量优化。
\end{cautionbox}

这里也能解释“梯度不知道归一化常数，为何还足够”。若 $p(x)=\widetilde p(x)/C$，其中 $C$ 与 $x$ 无关，则 $\grad\log p(x)=\grad\log\widetilde p(x)$。局部梯度不必计算这个常数。不过，只知道梯度并不等于可以随意选择采样器；在断开的支持上，相对概率质量也有额外问题。正方差的高斯扰动把分布平滑到整个空间，有助于避开这些几何障碍。

\section{局部回归怎样组织成全局生成？}
\label{sec:c-global}

\subsection{很多简单任务，共用一个网络}

训练时，我们可以随机选噪声级别 $t$，直接由式\eqref{eq:c-noising}获得观测。无需每次都真的从 $0$ 一步一步加噪到 $t$。一个时间条件网络在大量样本和噪声级别上学习，得到一组相互关联的预测规则。小噪声任务偏向恢复细部，大噪声任务更依赖总体结构，但这只是常见现象，不是每个采样步都具有固定语义职责的定理。

许多时间任务共享参数，既节省资源，也使模型能利用尺度之间的共同规律。可以把它想象成一位修复师：轻微损伤时，看清的细节较多，需要微调；严重损伤时，必须更多依靠对对象的总体知识。但真正的网络未必像人一样先认出一个离散语义再动笔，它执行的是连续函数计算。

最简单的噪声预测训练目标具有形式
\[
 \E_{t,X_0,\eps}
 \left[\norm{\eps-\eps_\theta(a(t)X_0+b(t)\eps,t)}^2\right].
\]
这条式子很好写，也容易构造训练数据，但它的作用不能只凭直觉判断。为何这种目标与概率建模有关，不同时间是否应有不同权重，以及怎样换算成反向均值，都是第\ref{ch:math}章要完整回答的问题。形式简单，不代表背后的对象简单。

\subsection{训练和采样走在不同方向上}

\begin{table}[htbp]
\centering
\caption{把训练与采样分开理解。二者复用同一个网络，但输入来源、已知信息与任务方向不同。}
\label{tab:c-training-sampling}
\small
\begin{tabularx}{\linewidth}{@{}p{0.17\linewidth}XX@{}}
\toprule
 & 训练 & 生成采样\\
\midrule
起点 & 真实数据 $x_0$ 与自行抽取的噪声 & 一份从简单先验抽取的随机状态\\
已知信息 & 原图、所加噪声和噪声级别 & 当前状态、时间和可能的外部条件\\
网络任务 & 学习条件均值、score 或相关预测量 & 为采样器提供当前状态的更新信息\\
时间使用 & 通常随机选择时间，无须依次遍历 & 按设定的时间序列或连续积分推进\\
最终目的 & 调整参数，使许多局部预测准确 & 将初始样本群搬运为目标分布\\
\bottomrule
\end{tabularx}
\end{table}

若真实反向转移分布已知，并且从正确的末端分布出发，逐步使用这些转移，就会返回正确的数据分布。这本质上是联合分布的条件概率分解，而不是“噪声会自动变成图片”的神奇规律。学习模型用有限网络近似反向规律，采样器再用有限步数近似所需过程。训练误差、末端近似和求解误差都会改变最终输出。

因此扩散有效的直觉可以分成三层。第一，正向过程人为规定，训练样本容易构造。第二，邻近噪声级别之间的变化比较小，可以把困难生成拆成较温和的预测。第三，概率论给出这些局部规则应怎样组织成全局过程。真正把三层接起来的是数学；直觉负责让公式的用途变得可见。

请特别注意“邻近”这个词。跨越很大噪声间隔时，反向条件分布可能高度多峰，一个简单高斯未必表达得好。小时间步下的局部近似更可信，但代价是采样步骤增多。这正是后来快速采样、高阶求解、蒸馏和路径设计反复面对的张力。扩散不是免费消除复杂性，而是把复杂性分配到一条有结构的路径上。

\subsection{随机性不是误差，它承担选择的工作}

同一句提示词，多次运行模型得到不同构图，通常是我们想要的行为。未被提示词规定的姿势、背景细节、纹理等，都需要在合理范围内随机选择。初始噪声提供这种随机性；在随机采样过程里，中途噪声也参与选择。在确定性采样过程里，初始随机向量已经足以产生结果差异。

随机种子只是产生伪随机序列的一个索引，不是模型内部的“图片编号”，也不是可解释的语义编码。改变种子会改变采样轨迹，不能保证只改变一个指定对象。固定种子有助于比较参数修改的效果，但这种控制依赖于相同的模型、时间序列和实现；换了采样设置，同一个种子也未必对应同一构图。

另外，反向随机过程不必逐样本重走当初那条正向加噪路径。正确性通常关心各时刻和最终的分布，而不是每个噪声粒子能否找回“自己的原图”。这一区分将把 DDPM、DDIM 和确定性概率流统一到同一个理解框架中。

\section{隐空间：哪些东西被隐藏了？}
\label{sec:c-latent}

\subsection{同一个词，三种常见用法}

\term{隐变量或潜变量}{gl:latent}是模型中未被直接观测、却参与描述或产生数据的变量。“隐空间”则是这些变量所处的空间。它并不专指一种网络，也不意味着人类看不懂的地方一定叫隐空间。讨论扩散时，至少应区分三种用法。

第一种是生成器的随机输入 $Z$。例如 GAN 用它决定输出的变化；$Z$ 未作为训练图像的直接标签被观测，因此可称为隐变量。它的维数、分布和语义由设计与训练共同决定。高斯输入不保证各坐标变成独立语义因素。

第二种是概率模型中解释观测的未观测变量。VAE 用 $z$ 参与 $p_\theta(x\mid z)$；DDPM 的 $x_1,\ldots,x_K$ 对观测 $x_0$ 而言也是一串辅助隐变量。于是，像素空间扩散仍可被称为隐变量模型。这里“隐”说的是概率角色，不是在说它已经压缩了图像。

第三种是潜空间扩散中的压缩表示。一个编码器 $E$ 把图片变成 $z_0=E(x_0)$，扩散在这个表示空间里学习；一个解码器 $D$ 把生成的 $z_0$ 转回图片。它把表示学习与分布学习拆成两个环节，是 latent diffusion 的主要安排。\cite{c_rombach2022ldm}

这三种用法可以同时出现在同一个系统里。最初随机抽取的 $z_K$ 是随机状态；去噪中的 $z_k$ 是轨迹上的隐变量；最终 $z_0$ 是需要解码的压缩表示。符号常都写作 $z$，但它们的统计角色和时间位置不同。阅读论文时先问“这个 $z$ 是编码结果、先验样本还是含噪状态”，会减少许多误解。

\subsection{压缩表示为什么值得做}

很多像素变化对理解画面并不同等重要。相邻像素的细微差别、某些纹理和压缩噪声，可能花费大量维数，却不改变对象与构图。如果编码器能把重要结构放进较小的空间，扩散网络便可以在较少的位置上运算，最后由解码器补回可见细节。

作一个维数比较：一张 $512\times512\times3$ 图片有 $786\,432$ 个标量；若表示为 $64\times64\times4$ 的潜张量，则有 $16\,384$ 个标量，数量为原来的 $1/48$。这说明表示大小可以显著减少，但实际运行时间不必减少 $48$ 倍，因为通道宽度、网络深度、注意力、编码解码成本和硬件行为都在影响计算。

这里的 $64\times64\times4$ 只是一个示例，不是所有潜空间扩散的固定格式。潜变量也可以有不同空间压缩率、通道数、离散编码或时间维度。对视频，压缩时间轴可以降低开销，同时也可能损失快速运动细节。不同表示的选择，决定了后面的生成器容易保留什么、难以恢复什么。

\begin{figure}[htbp]
\centering
\begin{tikzpicture}[x=1cm,y=1cm]
 \node[draw,rounded corners,fill=blue!6,minimum width=2.3cm,minimum height=1.0cm,align=center] (img) at (0,0) {像素图像\\$x_0$};
 \node[draw,rounded corners,fill=blue!6,minimum width=2.3cm,minimum height=1.0cm,align=center] (lat) at (3.6,0) {压缩表示\\$z_0=E(x_0)$};
 \node[draw,rounded corners,fill=orange!8,minimum width=2.3cm,minimum height=1.0cm,align=center] (noise) at (7.2,0) {含噪表示\\$z_k$};
 \draw[->,thick] (img) -- node[above,font=\small] {编码} (lat);
 \draw[->,thick] (lat) -- node[above,font=\small] {加噪} (noise);
 \draw[->,thick] ([yshift=-0.25cm]noise.west) to[bend left=28] node[below,font=\small] {学习反向生成} ([yshift=-0.25cm]lat.east);
 \node[draw,rounded corners,fill=green!7,minimum width=2.3cm,minimum height=1cm,align=center] (out) at (3.6,-2.0) {输出图像\\$\widehat x_0=D(\widehat z_0)$};
 \draw[->,thick] (lat) -- node[right,font=\small] {解码} (out);
 \node[font=\small,align=center] at (0,-2.0) {表示决定信息瓶颈\\扩散学习潜变量分布};
\end{tikzpicture}
\caption{潜空间扩散的两个学习层次。图中编码与解码说明表示的角色；生成时从潜噪声开始，并不先输入一张真实图片。}
\label{fig:c-latent}
\end{figure}

\subsection{潜空间不是免费的完美语义空间}

压缩意味着取舍。若 $D(E(x))$ 已经把小字、精细线条或某种医学特征损坏，后面的扩散模型就不能普遍保证把它们准确恢复。它可能生成一个看起来合理的细节，但合理并不等于与原始事实一致。潜空间扩散的输出能力受到表示的重建能力和解码器行为共同限制。

从概率角度看，扩散首先学习的是训练图像经编码后形成的潜变量分布。设 $Z_0=E(X_0)$，理想生成结果满足 $\widehat Z_0$ 与 $Z_0$ 同分布，那么解码后复现的是 $D(Z_0)$ 的分布。除非编码解码在目标任务下足够保真，它并不精确等于原始 $X_0$ 的分布。这个推送关系比“潜空间里学好了，图像就自然学好了”更准确。

同样，给每张图像的编码加一个高斯正则约束，也不意味着整批编码形成的分布精确是标准高斯。单个样本的近似后验、所有样本的聚合分布和生成先验是不同对象。我们可以专门在潜空间里训练扩散，恰恰是为了学习那个复杂的聚合分布。第\ref{ch:applications}章会进一步讨论表示模型与生成主干的合作。

\section{条件生成：把“可能”缩小到“符合要求”}
\label{sec:c-conditional}

\subsection{一句描述通常对应许多图片}

无条件模型学习 $\pdata(x)$；条件模型学习 $p_\theta(x\mid c)$。条件 $c$ 可以是类别、文字、草图、深度图、已有视频片段或测量结果。一句“窗边的一只橘猫”约束了部分内容，却没有规定相机角度、猫的动作、窗户形状以及每一根毛的状态。于是条件生成仍然是分布问题，不能一般性地压成唯一正确图像。

训练时若有配对样本 $(x_0,c)$，就可以给网络额外输入 $c$，学习
\[
 \eps_\theta(x_t,t,c)\approx \E[\eps\mid X_t=x_t,C=c].
\]
从 score 的角度看，它在估计 $\grad_x\log p_t(x\mid c)$。条件越有信息，允许的解释一般越受限制；但“信息多”不保证数据中的条件可靠，也不保证网络准确使用了每个条件。

如果提示词遗漏某种属性，模型会依据训练分布补充；如果提示词要求一个训练中很少见的组合，模型可能只完成其中一部分。这说明条件不只是用户界面的输入，还决定了概率建模的可识别性与泛化难度。词语组合能够表达的空间通常远大于有限配对数据实际覆盖的空间。

\subsection{条件是怎样进入图像网络的}

条件可以转换成特征，再通过拼接、调制或者注意力与图像表示发生交互。文字编码器会把文本转换为一组\term{词元}{gl:token}表示；交叉\term{注意力}{gl:attention}让图像位置根据当前状态读取有关的文本信息。此处重点是信息通道：概率目标说明应学习什么关系，具体结构说明怎样把条件送入计算。

文本编码器强，不意味着生成结果一定能逐字遵循指令。生成主干仍要把语言关系映射到空间关系；多个对象的属性绑定、数量、位置和细小文字都可能出错。条件表示、训练配对和网络交互方式分别影响不同环节。第\ref{ch:applications}章会从具体模型中看到这些问题怎样被处理。

另一个常见工具是 guidance，中文通常称“引导”。它在采样时加强条件的影响。第\ref{sec:guidance}节将从概率角度说明其由来，第三章再比较具体实现。现在只需记住：增强提示词符合度通常会改变输出分布，可能同时改变多样性、饱和度与细节。调高引导强度不等于无条件地提高模型的全部能力。

\subsection{图像修复也是条件概率问题}

设真实图像 $X$ 经测量算子 $A$ 得到 $Y=AX+\eta$，例如 $A$ 只保留部分像素，或执行低分辨率采样。恢复任务应研究 $p(x\mid y)$。给定缺损图像，缺失区域可能有多种合理补全；给定低分辨率图像，高频纹理也可能无法唯一确定。生成式修复正是用数据先验补充测量没有提供的信息。

对于拍照滤镜，我们可能欢迎丰富而美观的补充；对于科学测量，则要区分观测支持的结构与先验补出的结构。数学上的条件分布让这一区分可以被清楚表达：结果的某些部分被数据约束，另一些部分仍然存在不确定性。理解扩散的能力，也需要理解它对不确定性的使用方式。

\section{把扩散放回生成模型的全景}
\label{sec:c-families}

扩散不是第一个学习图像分布的模型，也不是把所有其他方法简单替换掉的一条单线进化。不同家族选择不同的概率表达、学习信号和采样计算；许多现代系统会混合多个家族。下面按“如何组织分布”理解它们，而不是只记住几个缩写。

\subsection{自回归：把联合概率拆成条件概率}

\term{自回归模型}{gl:ar}利用概率链式法则，
\[
 p_\theta(x_1,\ldots,x_d)
 =\prod_{i=1}^{d}p_\theta(x_i\mid x_1,\ldots,x_{i-1}).
\]
先决定顺序，再预测下一个位置。对图像，原始 PixelRNN 系列采用像素级的条件建模；对于现代视觉系统，序列元素也可以是压缩后的离散词元。\cite{c_oord2016pixelrnn}这种方式使似然分解直接，训练可以通过已知前缀提供监督，采样则依赖顺序生成。

它的关键选择是顺序。图片本来是二维整体，却要被排成一个序列；之前生成的内容成为之后生成的条件。扩散则通常在每一步同时更新所有位置，通过多次整体处理协调它们。这里的“同时”指状态更新，不意味着每个位置互相独立。

不应简单地说自回归一定慢、扩散一定快。逐像素采样可能有很长序列，压缩词元会缩短它；扩散每一步能并行处理位置，却常需要多次调用网络。实际速度由序列长度、采样步数、模型大小和硬件共同决定。概率分解的不同，让二者的计算瓶颈也不同。

\subsection{VAE：通过潜变量解释数据}

\term{变分自编码器}{gl:vae}首先规定
\[
 p_\theta(x)=\int p(z)p_\theta(x\mid z)\,\dd z,
\]
然后用编码分布 $q_\phi(z\mid x)$ 近似潜变量后验。这样既有生成路径，也有从观测推断潜变量的路径。由于积分常不可直接计算，训练使用\term{证据下界}{gl:elbo}；重参数化使其中的重要期望可以通过梯度进行优化。\cite{c_kingma2014vae}

普通自编码器只学编码和重建，并不自动保证“随机抽一个潜变量就能得到合理图片”。VAE 增加了概率模型和分布约束，但这些约束同样存在表达能力与重建精度的取舍。某些 VAE 的模糊输出与简单解码分布及平方误差有关，不能把“所有 VAE 都模糊”当成家族定义。

DDPM 也能从含有许多辅助隐变量的模型出发写出变分下界。因此扩散与 VAE 不是完全无关的两个数学世界。不过，标准 VAE 常用较直接的潜变量到图像解码，扩散把生成分配到许多噪声尺度；潜空间扩散又可以用一个 VAE 类表示模型承接图像与潜变量。关系在概率建模层面与系统组成层面各有一层。

\subsection{GAN：通过对抗反馈学习一个采样器}

\term{生成对抗网络}{gl:gan}把随机输入经过生成器 $G_\theta$ 变成样本，再让判别器区分真实与生成样本。生成器根据判别器提供的反馈更新，从而学习一个能产生目标分布样本的映射。\cite{c_goodfellow2014gan}它通常无须计算输出的归一化密度，采样也可以只需一次主生成器调用。

这种安排能让学习信号强调真假可辨的细部，因此在图像生成史上取得了重要成果；同时，两个网络互相影响会造成优化上的困难，生成器还可能忽略一些数据模式。所谓“模式”是某类具有独特内容或结构的可能性，例如一类姿态或一个数字类别，而不只是纹理风格。模式缺失就是输出分布没有充分覆盖它们。

扩散把训练组织成已知扰动下的预测，通常不需要判别器参与核心目标。这有助于避免某些对抗训练难题，但不保证不会发生覆盖不足、训练失衡或记忆。后来的快速扩散模型也可以结合对抗损失。因此“使用去噪”与“使用判别器”并非互斥的系统属性。

\subsection{归一化流：通过可逆变换管理密度}

\term{归一化流}{gl:flow}取可逆映射 $x=f_\theta(z)$，用变量替换公式写出
\[
 \log p_\theta(x)
 =\log p_Z(f_\theta^{-1}(x))
 +\log\left|\det J_{f_\theta^{-1}}(x)\right|.
\]
这里 $J$ 是\term{雅可比矩阵}{gl:jacobian}。它描述变换怎样拉伸或压缩体积，因此密度需要随体积改变。Real NVP 等模型通过特殊结构使可逆计算和行列式计算可行。\cite{c_dinh2017realnvp}

归一化流把采样和密度计算联系得十分直接，但计算方便的可逆结构也约束了设计。与之相比，常见去噪网络不必自身可逆；整个生成过程的数学规律负责组织分布变换。连续的确定性流与概率流 ODE 又会把这些思路重新联系起来。

还有一个术语陷阱：“normalizing flow”中的 flow，与“flow matching”不是同一个方法名。前者强调用可逆变换和体积变化表达密度；后者强调用回归目标学习一个连续速度场。二者可以相交，例如由满足适当条件的速度场积分得到连续归一化流，但不能仅看到 flow 就假设模型使用同一种目标或同一种采样器。

\subsection{能量模型：先描述相对偏好，再想办法采样}

能量模型写成
\[
 p_\theta(x)=\frac{\exp(-E_\theta(x))}{Z_\theta},\qquad
 Z_\theta=\int \exp(-E_\theta(x))\,\dd x.
\]
低能量表示相对偏好；配分函数 $Z_\theta$ 使密度归一化。一个任意的能量函数未必拥有有限配分函数，所以数学上还需要可积性。对高维神经能量模型，归一化以及训练中的采样往往很难。

能量梯度与 score 有直接联系：$\grad_x\log p_\theta(x)=-\grad_xE_\theta(x)$。模型可以通过这类局部方向参与随机采样。现代神经能量模型提供了这种思路的具体实例。\cite{c_du2019ebm}扩散 score 模型则对一系列噪声分布分别学习方向，通常不需要在每次训练更新中先运行完整的生成链。

但 score 网络也不一定显式表示为某个标量能量的精确梯度。理想目标是一个梯度场，有限网络近似后可能不完全具有对应的保守场结构。这不会阻止我们把它放入采样器，却提醒我们：一个有数学来源的目标，与有限学习结果满足全部数学性质，是两个不同层次。

\begin{table}[htbp]
\centering
\caption{生成模型家族的核心安排。比较针对典型形式，实际系统可混合多个方法。}
\label{tab:c-family-overview}
\small
\begin{tabularx}{\linewidth}{@{}p{0.13\linewidth}XXX@{}}
\toprule
家族 & 分布如何表达 & 主要学习信号 & 典型采样计算\\
\midrule
自回归 & 按顺序分解条件概率 & 已知前缀下的似然 & 依次产生序列元素\\
VAE & 潜变量与条件解码分布 & 变分下界与重建 & 抽潜变量，再解码\\
GAN & 随机输入的生成映射 & 真假判别的反馈 & 执行生成映射\\
归一化流 & 可逆变换与体积因子 & 可计算的对数密度 & 从简单分布经可逆映射\\
能量模型 & 能量与归一化常数 & 似然或替代目标 & 常需迭代随机采样\\
扩散 / score & 多噪声级别的反向规律 & 去噪、score 与变分目标 & 多步随机或确定性更新\\
\bottomrule
\end{tabularx}
\end{table}

\section{建立四层坐标，避免把名字混成一团}
\label{sec:c-four-layers}

读到新的生成模型时，建议先把它放到四层坐标上。第一层是表示：状态是像素、压缩张量还是离散词元？第二层是概率路径：数据怎样与简单分布联系起来，时间中间是什么分布？第三层是学习目标：网络预测噪声、干净数据、score 还是速度？第四层是网络与求解：怎样表达这个函数，怎样按有限步数生成？

例如，“潜空间”“DiT”“噪声预测”“DDIM”分别主要回答上述不同层的问题。它们不一定相互竞争，反而可以组合成一个系统。一个 Transformer 主干可以做噪声预测，也可以做速度回归；一个潜表示可以服务扩散，也可以服务自回归生成；换一个采样器通常不等于重新学习全部表示。

这里还需把“结构”一词说清楚。U-Net 和 DiT 主要指神经网络结构；DDPM 指一组概率建模、训练和采样安排；DDIM 主要提供与训练边缘相联系的另一类生成轨迹；flow matching 是学习连续速度场的一类方法。第三章将沿数学思想梳理它们的发展关系，而不会把这些名字放在同一张架构排名表上。

\begin{keybox}{进入第二章前，先保留这五句话}
生成的对象是分布；随机性选择具体样本。加噪给复杂分布建立了一串较平滑的中间分布。去噪的条件均值可以转换为局部概率方向。局部方向需要正确的时间机制才能产生全局采样。表示、概率路径、预测目标和网络结构是相互关联但可分别讨论的层次。
\end{keybox}

我们现在已经获得扩散可能有效的直觉，却还欠着严格答案：离散反向均值为何可以用噪声预测来表示？训练目标为何来自变分下界？score 怎样进入反向 SDE？概率流 ODE 如何让随机与确定性生成相遇？速度场又如何统一描述概率搬运？第\ref{ch:math}章将把这些欠下的推导逐一补齐。

\section{练习与解答思路}
\label{sec:c-exercises}

\begin{exercisebox}{练习一：平方误差为何偏爱平均答案？}
设 $X$ 是二阶矩有限的随机向量，模型必须输出固定向量 $m$。证明 $\E\norm{X-m}^2$ 在 $m=\E X$ 时最小。再解释：为什么这个结论不意味着生成模型应当总输出均值？
\end{exercisebox}

\begin{exercisebox}{练习二：一个有偏的两点分布}
设 $\mathbb P(X_0=h)=\pi$、$\mathbb P(X_0=-h)=1-\pi$，其中 $0<\pi<1$，并观察 $Y=X_0+\sigma\eps$。写出后验对数几率和最优去噪器。若 $\pi>1/2$，在 $Y=0$ 时去噪结果是否仍为零？这说明了什么？
\end{exercisebox}

\begin{exercisebox}{练习三：确定性函数能产生随机分布吗？}
设 $U$ 在 $(0,1)$ 上均匀分布，$X=\mathbf 1_{\{U>\frac34\}}$。求 $X$ 的分布。随后设 $G$ 是固定网络，讨论“采样过程没有中途噪声，所以不可能产生多样性”这句话哪里有误。
\end{exercisebox}

\begin{exercisebox}{练习四：理想潜空间生成能保证什么？}
设编码器为 $E$、确定性解码器为 $D$，潜空间生成器精确复现 $Z=E(X)$ 的分布。证明输出分布等于 $D(E(X))$ 的分布。给出一个简单例子，说明它不一定等于 $X$ 的分布；再给出一个充分条件，使二者相等。
\end{exercisebox}

\begin{exercisebox}{练习五：把一个系统放到四层坐标上}
某图像系统先将图片编码为连续潜张量，以高斯扰动构造训练数据，用 Transformer 预测噪声，再用确定性多步更新采样。分别指出它的表示、概率路径、学习目标与网络/求解选择。说明“它是 Transformer，所以不属于扩散”为什么不成立。
\end{exercisebox}

\subsection*{解答思路}

\paragraph{练习一。}
令 $\mu=\E X$，展开得到
\[
 \E\norm{X-m}^2
 =\E\norm{X-\mu}^2+\norm{\mu-m}^2,
\]
交叉项因 $\E(X-\mu)=0$ 而消失。均值是给定限制下的最优估计，但生成要求产生符合整个分布的随机结果，而不是在固定向量集合里挑选一个最优代表。对多峰分布，均值甚至可能落在低概率区域。

\paragraph{练习二。}
贝叶斯公式给出后验对数几率
\[
 \ell(y)=\log\frac{\pi}{1-\pi}+\frac{2hy}{\sigma^2}.
\]
设后验正端点概率为 $r(y)=1/(1+\exp(-\ell(y)))$，则
$D^*(y)=h(2r(y)-1)=h\tanh(\ell(y)/2)$。
当 $y=0$ 时，$D^*(0)=h(2\pi-1)$。观测无明显倾向时，数据先验仍会参与估计；去噪器确实携带分布的概率权重，而不只是做几何投影。

\paragraph{练习三。}
$\mathbb P(X=1)=1/4$，$\mathbb P(X=0)=3/4$。映射是确定的，但输入是随机的。相同道理适用于 $G(Z)$ 和从随机初值出发的确定性积分。多样性的质量取决于映射是否把输入概率正确组织起来，不取决于是否每一步都掷硬币。

\paragraph{练习四。}
若 $\widehat Z$ 与 $E(X)$ 同分布，则对任意可测集合 $A$，
\[
 \mathbb P(D(\widehat Z)\in A)
 =\mathbb P(E(X)\in D^{-1}(A))
 =\mathbb P(D(E(X))\in A).
\]
若 $X$ 等概率取 $-1,1$，$E$ 把两者都编码为 $0$，且 $D(0)=0$，输出就只取 $0$。一个充分条件是 $D(E(X))=X$ 几乎处处成立；“重建视觉上相近”只是近似条件，不能推出分布精确相等。

\paragraph{练习五。}
表示是连续压缩潜张量；训练路径是所指定的高斯扰动路径，还需知道具体系数才能进一步分类；预测目标是噪声；网络是 Transformer，求解是确定性多步更新。主干结构决定函数如何计算，扩散的概率安排决定它为什么这样训练和采样，二者属于不同层次。

\chaptertakeaway{本章把“生成一张图”改写为“从一个学到的分布取样”。图像的高维结构解释了学习困难，平滑与多尺度解释了加噪的作用，条件回归与 score 的联系解释了去噪为何携带概率信息。我们同时保留了三个边界：条件均值不是后验样本，密度上升不是正确采样，潜空间学得好也受编码解码的信息瓶颈限制。带着这些区分进入第二章，公式就会成为解决明确问题的工具。}
}{}
\IfFileExists{chapters/02-math.tex}{\chapter{数学主线：从去噪回归到概率输运}\label{ch:math}
\readingguide{这一章把“从噪声生成数据”变成可以证明和计算的命题。先推导离散 DDPM 的边缘分布、后验和训练目标；再说明噪声回归为何给出密度的 score；最后进入 SDE、概率流 ODE 与条件流匹配。每条主线都围绕同一个问题：我们没有真实密度的公式，怎样仅凭数据样本学到足以生成的规律？}

\section{先辨认数学对象，再开始推导}
给定训练样本，我们想让最终生成随机变量的分布接近 $\pdata$。这句话包含三个不同对象。第一个是单个样本的位置 $x$；第二个是描述大量样本如何分布的密度 $p(x)$；第三个是让样本移动的规则，例如转移概率、漂移或速度。生成模型的任务通常是学习第三个对象，从而控制第二个对象。不能把一张图像的变化轨迹当成整个分布的变化规律。

为了把复杂问题拆开，我们人为选定一条容易制造样本的概率路径。DDPM 选定逐步加噪的马尔可夫链；连续扩散选定 SDE；流匹配可以直接选定两个端点之间的插值。这些路径是我们设计的训练工具，不是数据天然附带的物理过程。训练后，我们从路径的简单端出发生成，要求恢复的是数据的分布，而不是逐张恢复训练集。

\begin{keybox}{本章的理想化条件}
我们把图像表示成 $\R^d$ 中的连续向量，先讨论无条件生成。需要微分密度时，假设噪声端以外的相关密度足够光滑且正，积分、求导和条件期望可以交换；涉及积分分部时要求边界项消失。平方投影要求监督具有有限二阶矩；跨时间的 DSM 还要求加权风险有限，可以截断零噪声端点，或选择合适权重（例如 $\lambda=b^2$）。ODE 与 SDE 的存在唯一性还需要相应的正则性与增长条件。真实离散像素、低维流形和零噪声端点如何处理，见附录\ref{app:probability}与\ref{app:stochastic}。
\end{keybox}

基本概率符号也值得重申：$q(x_k\mid x_0)$ 是给定某个数据样本后的条件分布；$q(x_k)=\int q(x_k\mid x_0)\pdata(x_0)\dd x_0$ 是将所有数据样本混合后的边缘分布。前者可以是简单高斯，后者仍然可能具有多峰和复杂几何。扩散学习最巧妙的地方，正是利用简单条件分布提供监督，学会处理复杂边缘分布。

\section{DDPM 前向过程：可计算的高斯链}\label{sec:ddpm-forward}
\subsection{单步加噪与累计信号}
离散时间用 $k=0,1,\ldots,K$。定义
\begin{align}
 q(x_{1:K}\mid x_0)&=\prod_{k=1}^Kq(x_k\mid x_{k-1}),\\
 q(x_k\mid x_{k-1})&=\N\bigl(\sqrt{\alpha_k}x_{k-1},\beta_kI\bigr),\qquad
 \alpha_k=1-\beta_k,\quad 0<\beta_k<1.
\end{align}
这里 $\beta_k$ 是单步引入的噪声方差，$\alpha_k$ 是单步信号平方尺度。$\sqrt{\alpha_k}$ 轻微缩小输入，避免反复加噪导致总方差无限增长。若输入已经是标准高斯，那么缩放后的方差 $\alpha_kI$ 加上噪声方差 $\beta_kI$，仍是 $I$。因此标准高斯是这条链的不变分布。

这不是说任意输入一步就成为高斯。一步得到的分布是“缩小后的数据分布与高斯的卷积”。要使数据的信息逐渐退场，需要累计信号衰减。

记 $\bar\alpha_k=\prod_{j=1}^k\alpha_j$。写出前两步就能看出规律：
\begin{align}
 x_1&=\sqrt{\alpha_1}x_0+\sqrt{\beta_1}z_1,\\
 x_2&=\sqrt{\alpha_2\alpha_1}x_0
 +\sqrt{\alpha_2\beta_1}z_1+\sqrt{\beta_2}z_2,
\end{align}
其中 $z_1,z_2$ 相互独立，均为标准高斯。后两项之和是高斯，方差为
$\alpha_2\beta_1+\beta_2=1-\alpha_1\alpha_2$。一般地，若第 $k-1$ 步的条件方差是 $1-\bar\alpha_{k-1}$，下一步为
\begin{equation}
 \alpha_k(1-\bar\alpha_{k-1})+\beta_k=1-\bar\alpha_k.
\end{equation}
数学归纳得到
\begin{equation}\boxed{
 q(x_k\mid x_0)=\N\bigl(\sqrt{\bar\alpha_k}x_0,(1-\bar\alpha_k)I\bigr).
 }\label{eq:forward-marginal}\end{equation}
因此训练时不必真的走过 $k$ 次加噪：抽取一个 $\epsilon\sim\N(0,I)$，直接令
\begin{equation}
 x_k=\sqrt{\bar\alpha_k}x_0+\sqrt{1-\bar\alpha_k}\epsilon.
\end{equation}
这是一种\term{重参数化}{gl:reparameterization}：把带参数的随机变量表示成确定性变换加上标准噪声。一次直接构造只匹配这个时刻的边缘；若多个时刻独立重采样 $\epsilon$，它们的联合分布不再是原来那条马尔可夫链。训练噪声回归只需要单时刻监督，故这一区别不构成问题。

\subsection{为什么终点只是近似标准高斯}
当 $\bar\alpha_K$ 很小时，$x_K$ 的数据项接近零，条件分布接近标准高斯。但只要 $\bar\alpha_K>0$，在有限步数下一般不会完全等于标准高斯。若数据的均值为 $m$、协方差为 $C$，则
\begin{align}
 \E[x_K]&=\sqrt{\bar\alpha_K}m,\\
 \Cov(x_K)&=\bar\alpha_KC+(1-\bar\alpha_K)I.
\end{align}
即使这些一、二阶量接近高斯，高阶统计量也可能残留。把生成起点设成 $\N(0,I)$ 是有意采用的近似；噪声日程、终点信噪比与训练、采样的一致性都影响误差。

\term{信噪比}{gl:snr}在 VP 表示下通常定义为
\begin{equation}
 \SNR(k)=\frac{\bar\alpha_k}{1-\bar\alpha_k}.
\end{equation}
它比较信号与注入噪声的平方尺度；如果数据各方向方差不同，真实方向上的信噪比还要乘以对应数据方差。因此它是日程的指标，并不是所有图像和所有频率具有完全相同的可辨识度。

\section{反向后验：有原图时为什么能算}\label{sec:ddpm-posterior}
我们希望生成时使用 $q(x_{k-1}\mid x_k)$，但它混合了所有可能的 $x_0$，通常难以直接计算。训练时我们却知道抽到的原图，因而可以计算更强的条件后验 $q(x_{k-1}\mid x_k,x_0)$。这是“可计算教师”和“可部署学生”之间的关键区别。

利用贝叶斯公式和前向边缘，对 $k\ge2$ 有
\begin{equation}
 q(x_{k-1}\mid x_k,x_0)\propto
 q(x_k\mid x_{k-1})q(x_{k-1}\mid x_0).
\end{equation}
令 $y=x_{k-1}$，对数密度中与 $y$ 有关的部分是
\begin{equation}
 -\frac{\norm{x_k-\sqrt{\alpha_k}y}^2}{2\beta_k}
 -\frac{\norm{y-\sqrt{\bar\alpha_{k-1}}x_0}^2}{2(1-\bar\alpha_{k-1})}.
\end{equation}
把平方展开，$\norm{y}^2$ 的精度系数与线性项分别为
\begin{align}
 A&=\frac{\alpha_k}{\beta_k}+\frac{1}{1-\bar\alpha_{k-1}}
   =\frac{1-\bar\alpha_k}{\beta_k(1-\bar\alpha_{k-1})},\\
 h&=\frac{\sqrt{\alpha_k}}{\beta_k}x_k+
 \frac{\sqrt{\bar\alpha_{k-1}}}{1-\bar\alpha_{k-1}}x_0.
\end{align}
配方 $-A\norm{y}^2/2+\ip{h}{y}=-A\norm{y-A^{-1}h}^2/2+\text{常数}$，于是
\begin{align}
 q(x_{k-1}\mid x_k,x_0)&=\N(\widetilde\mu_k,\widetilde\beta_k I),\label{eq:posterior}\\
 \widetilde\beta_k&=\frac{1-\bar\alpha_{k-1}}{1-\bar\alpha_k}\beta_k,\\
 \widetilde\mu_k&=\frac{\sqrt{\bar\alpha_{k-1}}\beta_k}{1-\bar\alpha_k}x_0
 +\frac{\sqrt{\alpha_k}(1-\bar\alpha_{k-1})}{1-\bar\alpha_k}x_k.
\end{align}
把 $x_0=(x_k-\sqrt{1-\bar\alpha_k}\epsilon)/\sqrt{\bar\alpha_k}$ 代回，可将均值写为
\begin{equation}
 \widetilde\mu_k=\frac{1}{\sqrt{\alpha_k}}
 \left(x_k-\frac{\beta_k}{\sqrt{1-\bar\alpha_k}}\epsilon\right).
 \label{eq:posterior-noise}
\end{equation}
注意这一项里的 $\epsilon$ 是此次训练构造中实际使用的噪声。生成时它不可观测；网络会根据 $x_k$ 和 $k$ 估计它。

\begin{cautionbox}{高斯后验与真实逆过程}
“已知 $x_0$ 的后验是高斯”不意味着“只知道 $x_k$ 的真实反向条件也是同一个高斯”。后者对 $x_0$ 的后验作混合，可能多峰。DDPM 使用单个高斯建模局部反向转移，在足够小的噪声步长下与反向扩散相协调；有限步的大跨度跳跃不自动保留这一近似精度。
\end{cautionbox}

\subsection{后验混合的均值与方差：误差藏在哪里}
把后验均值写成 $\widetilde\mu_k=A_kx_0+B_kx_k$，其中 $A_k=\sqrt{\bar\alpha_{k-1}}\beta_k/(1-\bar\alpha_k)$。对 $x_0\mid x_k$ 作平均，有
\begin{equation}
 \E[x_{k-1}\mid x_k]=A_k\E[x_0\mid x_k]+B_kx_k.
\end{equation}
因此一个达到理想噪声回归最优值的网络，可以得到真实逆条件的均值。这是比“高斯近似可能合理”更精确的结论。

但是方差还含有原图后验的不确定性。由全方差公式，
\begin{align}
 \Cov(x_{k-1}\mid x_k)
 &=\E[\Cov(x_{k-1}\mid x_k,x_0)\mid x_k]
 +\Cov(\E[x_{k-1}\mid x_k,x_0]\mid x_k)\\
 &=\widetilde\beta_k I+A_k^2\Cov(x_0\mid x_k).
\end{align}
第二项可以具有方向性，甚至在某些观测附近很大。仅把逆转移方差设为 $\widetilde\beta_kI$，不能把整个混合分布的协方差全部表示出来；采用另一个固定方差或学习方差，是处理近似的不同选择。

这也解释了为何一步从几乎纯噪声跳回干净图像困难。高噪声下许多原图都合理，后验均值可能缺少清晰结构；局部小步则不必一次解决所有歧义。小步并不是“让网络重复猜同一个答案”，而是在不断变化的边缘分布中逐步重建信息。其代价是多次网络调用，第三章的加速工作就从这里出发。

\section{ELBO 如何变成噪声均方误差}\label{sec:ddpm-elbo}
\subsection{先写完整的生成模型}
反向生成模型是
\begin{equation}
 p_\theta(x_{0:K})=p(x_K)\prod_{k=1}^Kp_\theta(x_{k-1}\mid x_k),\qquad
 p(x_K)=\N(0,I).
\end{equation}
$p_\theta(x_0)$ 必须把全部隐藏状态积分掉；直接求这个积分很困难。我们使用已知的 $q(x_{1:K}\mid x_0)$ 作为辅助分布。由 Jensen 不等式，
\begin{align}
 \log p_\theta(x_0)
 &=\log\E_q\left[\frac{p_\theta(x_{0:K})}{q(x_{1:K}\mid x_0)}\right]\\
 &\ge\E_q\left[\log p_\theta(x_{0:K})-\log q(x_{1:K}\mid x_0)\right].
\end{align}
右边是\term{证据下界}{gl:elbo}。下界与真实对数似然的差恰是辅助后验与模型后验之间的 KL；这个恒等式及 Jensen 的适用条件见附录\ref{app:probability}。

把负下界记作 $L_{\mathrm{VLB}}$。前向链也可沿反向条件分解为
\begin{equation}
 q(x_{1:K}\mid x_0)=q(x_K\mid x_0)
 \prod_{k=2}^Kq(x_{k-1}\mid x_k,x_0).
\end{equation}
它与生成分解逐项相减，得到
\begin{align}
 L_{\mathrm{VLB}}(x_0)
 =&\ \KL\bigl(q(x_K\mid x_0)\Vert p(x_K)\bigr)\nonumber\\
 &+\sum_{k=2}^K\E_{q(x_k\mid x_0)}
 \KL\bigl(q(x_{k-1}\mid x_k,x_0)\Vert p_\theta(x_{k-1}\mid x_k)\bigr)\nonumber\\
 &-\E_{q(x_1\mid x_0)}\log p_\theta(x_0\mid x_1).
 \label{eq:elbo}
\end{align}
三部分依次是终点匹配、每一步逆转匹配和最后一步重建。若前向日程固定，终点匹配不依赖 $\theta$。最后一步仍然是有效似然项；量化图像可以使用离散化高斯解码似然，连续教材模型也可使用适当连续解码分布。它不能凭空消失。

\subsection{两高斯之间的 KL 给出回归问题}
设 $p_\theta(x_{k-1}\mid x_k)=\N(\mu_\theta(x_k,k),\sigma_k^2 I)$，暂时把 $\sigma_k$ 固定。由高斯 KL 公式，
\begin{align}
 &\KL\bigl(\N(\widetilde\mu_k,\widetilde\beta_k I)
 \Vert\N(\mu_\theta,\sigma_k^2I)\bigr)\nonumber\\
 &=\frac{\norm{\widetilde\mu_k-\mu_\theta}^2}{2\sigma_k^2}
 +\frac d2\left(\frac{\widetilde\beta_k}{\sigma_k^2}-1+
 \log\frac{\sigma_k^2}{\widetilde\beta_k}\right).
\end{align}
对均值网络而言，第二项是常数。选取与式\eqref{eq:posterior-noise}匹配的参数化
\begin{equation}
 \mu_\theta(x_k,k)=\frac{1}{\sqrt{\alpha_k}}
 \left(x_k-\frac{\beta_k}{\sqrt{1-\bar\alpha_k}}\epsilon_\theta(x_k,k)\right),
\end{equation}
相减并平方后，每一步与参数有关的 KL 就是
\begin{equation}
 \frac{\beta_k^2}{2\sigma_k^2\alpha_k(1-\bar\alpha_k)}
 \E\norm{\epsilon-\epsilon_\theta(x_k,k)}^2.
 \label{eq:weighted-noise}
\end{equation}
这是一个有严格来源的加权均方误差。但常用 DDPM 训练目标会简化权重：
\begin{equation}\boxed{
 L_{\mathrm{simple}}=\E_{\substack{k\sim\mathrm{Unif}\{1,\ldots,K\}\\
 x_0\sim\pdata,\ \epsilon\sim\N(0,I)}}
 \norm{\epsilon-\epsilon_\theta(x_k,k)}^2.
 }\label{eq:simple-loss}\end{equation}
这一简化在 DDPM 中有重要作用\cite{r_ddpm}。它重新分配不同噪声尺度的学习重点，不能把它称为“原封不动的最大似然”或“完整 ELBO”。网络容量无限且每个时刻的权重严格为正时，各时刻最优回归函数不变；共享有限网络下，权重会改变优化、资源分配和最终误差。

如果同时学习方差，KL 中的方差项也依赖参数；只保留均值误差就不再是完整目标。为什么很多实现混合简单损失与变分损失，第三章会继续解释。

\subsection{回归最优值不是训练时那份噪声}
平方损失存在一个值得单独证明的性质。令输入为 $U$、监督为 $V$，记 $m(U)=\E[V\mid U]$。则
\begin{align}
 \E\norm{h(U)-V}^2
 &=\E\norm{h(U)-m(U)}^2+\E\norm{V-m(U)}^2.
 \label{eq:conditional-projection}
\end{align}
交叉项为零，因为 $\E[V-m(U)\mid U]=0$。因此
\begin{equation}
 \epsilon_*(x_k,k)=\E[\epsilon\mid x_k,k].
\end{equation}
这不是精确认出某次加噪的随机数。多个原图和多个噪声都能产生相同的 $x_k$；网络输出是后验平均。你可能担心“平均不是会造成模糊吗？”如果直接把这个均值作为整张最终图，确实可能如此。但扩散把平均作为局部转移的组成部分，配合不同尺度的迭代和必要的随机性来生成整个分布。接下来我们会看到，这个后验平均竟然恰好编码了边缘密度的梯度。

\section{Score：把条件去噪连接到边缘密度}\label{sec:score}
\subsection{为什么对密度取对数，再求梯度}
给定光滑正密度，定义
\begin{equation}
 s_t(x)=\grad_x\log p_t(x)=\frac{\grad_xp_t(x)}{p_t(x)}.
\end{equation}
这里的\term{score}{gl:score}是对状态 $x$ 求导，而不是经典统计里对参数求导的 score。它是向量场，给出当前位置上对数密度增加最快的方向和速率。若未归一化密度写为 $p(x)=e^{-E(x)}/Z$，那么 $s(x)=-\grad E(x)$；归一化常数 $Z$ 对 $x$ 的导数为零。这解释了为何学习 score 有机会避开昂贵的归一化积分。

但 score 不是“图像质量梯度”，也不是全局语义坐标。一个真实分布的两个对称峰之间，score 可以为零，而该位置的密度很低；仅做梯度上升可能停在那里或集中到少数峰。如何从 score 得到正确分布，需要动力学而不只是方向直觉。

\subsection{条件 score 的平均就是边缘 score}
采用连续表示 $x_t=a(t)x_0+b(t)\epsilon$，缩写 $a=a(t),b=b(t)>0$。已知条件核
\begin{equation}
 q_t(x\mid x_0)=\N(ax_0,b^2I),\qquad
 \grad_x\log q_t(x\mid x_0)=-\frac{x-ax_0}{b^2}=-\frac\epsilon b.
\end{equation}
我们不能直接计算 $p_t(x)$，但可以作形式上的求导：
\begin{align}
 \grad_xp_t(x)
 &=\int\pdata(x_0)\grad_xq_t(x\mid x_0)\dd x_0\\
 &=\int\pdata(x_0)q_t(x\mid x_0)
       \grad_x\log q_t(x\mid x_0)\dd x_0.
\end{align}
除以 $p_t(x)$，由贝叶斯公式识别出 $q_t(x_0\mid x)$，得到
\begin{equation}\boxed{
 s_t(x)=\E\left[\grad_x\log q_t(x\mid x_0)\mid x_t=x\right]
       =-\frac{\E[\epsilon\mid x_t=x]}{b(t)}.
 }\label{eq:dsm}\end{equation}
结合平方损失投影性质，噪声预测网络的最优值恰给出
$s_\theta(x,t)=-\epsilon_\theta(x,t)/b(t)$。这并不需要知道真实边缘密度的归一化常数。

称下面的可训练目标为\term{去噪 score matching}{gl:score}：
\begin{equation}
 L_{\mathrm{DSM}}=\E_{t,x_0,\epsilon}\lambda(t)
 \norm{s_\theta(x_t,t)+\epsilon/b(t)}^2.
\end{equation}
令 $H=-\epsilon/b(t)$，在每个 $t$ 上使用式\eqref{eq:conditional-projection}，便有
\begin{equation}
 L_{\mathrm{DSM}}=\E_{t,x_t}\lambda(t)\norm{s_\theta(x_t,t)-s_t(x_t)}^2+C,
\end{equation}
其中 $C$ 与 $\theta$ 无关。这说明条件监督和难以直接计算的边缘监督具有相同优化方向。它们的损失数值一般不相同；条件监督仍包含不可消除的方差。早期 score-based 生成及其连续统一框架分别见\cite{r_ncsn,r_sde}，更一般的 score matching 背景见\cite{r_scorematching}。

\subsection{Tweedie 关系：从 score 找到后验均值}
把式\eqref{eq:dsm}写成 $s_t(x)=-(x-a\E[x_0\mid x_t=x])/b^2$，可得
\begin{equation}\boxed{
 \E[x_0\mid x_t=x]=\frac{x+b(t)^2s_t(x)}{a(t)},\quad a(t)\ne0.
 }\label{eq:tweedie}\end{equation}
它常称为高斯去噪的 Tweedie 关系。加性噪声 $a=1$ 时最容易理解：后验均值等于当前点加上由 score 指示的修正。噪声大时偏离数据规律的可能性大，修正尺度 $b^2$ 也大。不过真实 score 自身会随噪声变化，不能只看这个系数推断每一步移动距离。

\begin{examplebox}{单峰高斯可以把关系全部算出来}
设 $x_0\sim\N(m,cI)$，则 $x_t\sim\N(am,(a^2c+b^2)I)$，
\begin{equation}
 s_t(x)=-\frac{x-am}{a^2c+b^2}.
\end{equation}
代入式\eqref{eq:tweedie}，得到
\begin{equation}
 \E[x_0\mid x_t=x]
 =m+\frac{ac}{a^2c+b^2}(x-am).
\end{equation}
当观测信号很弱时，估计向先验均值收缩；当噪声很小时，观测本身占主导。这里的“收缩”来自贝叶斯条件平均，并不是网络偶然学到的图像滤镜。
\end{examplebox}

\subsection{一个双峰例子：为什么需要多个尺度}
考虑一维数据为 $\tfrac12\N(-m,c)+\tfrac12\N(m,c)$。在加性噪声 $x=x_0+b\epsilon$ 下令 $h=c+b^2$，则
\begin{equation}
 p_b(x)\propto\exp\left(-\frac{x^2+m^2}{2h}\right)
 \cosh\left(\frac{mx}{h}\right),
\end{equation}
从而
\begin{equation}
 s_b(x)=-\frac{x}{h}+\frac m h\tanh\left(\frac{mx}{h}\right).
\end{equation}
在高噪声下，两个峰合并，score 更平滑；在低噪声下，中间附近的 score 将左右两侧分别引向两个峰。对称中心始终满足 $s_b(0)=0$，但当 $h<m^2$ 时它是低密度局部极小值。实际随机起点不会全部精确落在中心；随着路径演化，整团概率分配到不同峰。这个例子提醒我们，score 的零点与高概率点不能等同。

\begin{figure}[htbp]\centering
\begin{tikzpicture}
\begin{axis}[width=.94\linewidth,height=57mm,xmin=-4,xmax=4,
 xlabel={$x$},ylabel={密度},ymin=0,legend style={font=\small,draw=none},
 legend pos=north east,axis lines=left,samples=150,domain=-4:4]
\addplot[teal,thick]{0.5/sqrt(2*pi*0.3)*(exp(-(x-1.6)^2/(2*0.3))+exp(-(x+1.6)^2/(2*0.3)))};
\addlegendentry{低噪声：$h=0.3$}
\addplot[gold,thick,dashed]{0.5/sqrt(2*pi*3)*(exp(-(x-1.6)^2/6)+exp(-(x+1.6)^2/6))};
\addlegendentry{高噪声：$h=3$}
\end{axis}
\end{tikzpicture}
\caption{原创解析例子：高斯平滑把两个清晰峰变成更容易学习的宽分布。曲线由上式直接计算，$m=1.6$；它展示边缘密度而非某条去噪轨迹。}\label{fig:math-mixture}
\end{figure}

\subsection{Langevin 动力学：为什么沿 score 走还要加噪声}
在固定噪声尺度上，如果只用 $\dot x=s(x)$，每个样本都在爬升对数密度。最终可能集中到若干局部模式，而不是保持目标密度的峰宽与相对概率。要抽取整个分布，另一种经典动力学是
\begin{equation}
 \dd X_r=s(X_r)\dd r+\sqrt2\dd W_r.
 \label{eq:langevin}
\end{equation}
这里 $r$ 是在固定目标密度下运行的采样时钟，区别于改变噪声尺度的扩散时间 $t$。它的 Fokker--Planck 方程是
\begin{equation}
 \partial_r\rho_r=-\divg(\rho_rs)+\Delta\rho_r.
\end{equation}
如果 $s=\grad\log p$，代入候选稳态 $\rho=p$，因为 $ps=\grad p$，右边等于 $-\Delta p+\Delta p=0$。因此 $p$ 是不变密度。合适的遍历性、尾部和正则性条件才进一步保证从别的初始分布趋近它；仅仅验证右边为零并不证明收敛速度，也不证明任意离散步长下都能无偏抽样。

这个抵消提供了很具体的直觉：score 将粒子聚向高密度区域，随机扩散将它们铺开，两种作用在正确的密度下平衡。如果把随机项删掉，平衡就被破坏。与反向扩散不同，Langevin 在此追踪一份不随 $r$ 变化的目标密度；实际 annealed Langevin 依次更换平滑尺度，希望从较简单的密度走向较复杂的密度。有限次更新可能尚未达到每个尺度的平衡。

其离散更新为 $x\leftarrow x+\eta s(x)+\sqrt{2\eta}z$。它常用于 predictor--corrector 中的 corrector。加噪在这里是为了修正分布探索，不是因为网络没预测出应当消去的所有噪声。第三章会把它与用于跨时刻的 predictor 放在一起比较。

\begin{examplebox}{温度改变了目标，而不只是改变随机种子}
若固定漂移为 $s=\grad\log p$，将噪声强度改为 $\sqrt{2T}$，则在可归一化的条件下，候选不变密度为 $\rho_T\propto p^{1/T}$。因为 $\grad\rho_T=\rho_Ts/T$，代入 $-\divg(\rho_Ts)+T\Delta\rho_T$ 就得到零。$T<1$ 使密度更集中，$T>1$ 使它更分散。这个结论是固定密度的稳态分析，不能直接当成时间变化的 CFG 终点分布公式。
\end{examplebox}

\section{预测目标的等价与不等价}\label{sec:parameterization}
在 $x_t=ax_0+b\epsilon$ 下，预测噪声、预测原图与预测 score 之间可以作代数转换：
\begin{equation}
 \widehat x_0=\frac{x_t-b\widehat\epsilon}{a},\qquad
 \widehat s=-\frac{\widehat\epsilon}{b}.
\end{equation}
这些公式要求分母非零。靠近纯噪声端，$a$ 很小，噪声预测的一点误差转换成 $x_0$ 误差会被放大；靠近数据端，$b$ 很小，噪声误差转换成 score 误差也会被放大。不同目标不仅是改名字，还影响数值条件。

若 $a^2+b^2=1$，定义常用的 $v$ 目标\cite{r_vpred}：
\begin{equation}
 v=a\epsilon-bx_0.
\end{equation}
把两个关系合并，就是二维系数平面中的旋转：
\begin{equation}
 \begin{bmatrix}x_t\\v\end{bmatrix}
 =\begin{bmatrix}a&b\\-b&a\end{bmatrix}
 \begin{bmatrix}x_0\\\epsilon\end{bmatrix}.
\end{equation}
矩阵是正交矩阵，逆变换不需要除以接近零的 $a$ 或 $b$：
\begin{equation}\boxed{
 \widehat x_0=ax_t-b\widehat v,\qquad
 \widehat\epsilon=bx_t+a\widehat v.
 }\label{eq:vconversion}\end{equation}
这给出选择 $v$ 参数化的一个直观理由。但在同一有限网络和固定损失权重下，不同参数化的训练效果没有数学上自动相同的保证。

若写 $a=\cos\phi(t),b=\sin\phi(t)$，在固定 $x_0,\epsilon$ 的插值曲线上，
\begin{equation}
 \frac{\dd x_t}{\dd t}=\phi'(t)(a\epsilon-bx_0)=\phi'(t)v.
\end{equation}
$v$ 是对角度 $\phi$ 的导数目标；只有相应角度时间尺度下才直接等于这条条件曲线的速度。更进一步，网络的最优 $v$ 是条件均值，边缘 ODE 速度需要对隐含端点平均。SDE 的实际样本路径则几乎处处没有普通时间导数。把所有这些对象都称作“速度”会掩盖重要差异。

\subsection{损失权重与时间抽样也属于模型设计}
假设 $x_0$ 预测误差为 $e_x$，经代数转换的噪声误差是 $e_\epsilon=-(a/b)e_x$。因此
\begin{equation}
 \norm{e_\epsilon}^2=\frac{a^2}{b^2}\norm{e_x}^2.
\end{equation}
均匀权重的噪声损失等价于以 SNR 加权的原图损失，而不是均匀的原图损失。$v$ 损失经式\eqref{eq:vconversion}转换，也有自己的尺度权重。

同样，“均匀抽时间”只有在确定了时钟后才有意义。令 $\lambda(t)=\log(a(t)^2/b(t)^2)$ 为对数信噪比；从均匀 $t$ 改为在某个有限截断区间上均匀 $\lambda$，改变的是概率质量分配给各噪声尺度的方式。完整无穷 log-SNR 区间上不存在均匀概率分布。若目标积分为 $\int w(t)\ell(t)\dd t$ 而抽样密度为 $r(t)$，无偏估计应使用 $w(t)\ell(t)/r(t)$。实际方法也可能有意改变目标权重；要区分“重要性抽样加速同一目标”和“重新定义目标”。

\section{从离散链走向连续 SDE}\label{sec:sde}
\subsection{微小高斯增量意味着什么}
将时间区间分为小段 $\Delta t$，令单步噪声方差 $\beta_k=\beta(t_k)\Delta t$。一阶展开
\begin{equation}
 \sqrt{1-\beta(t)\Delta t}=1-\tfrac12\beta(t)\Delta t+O(\Delta t^2).
\end{equation}
因此 DDPM 的增量近似为
\begin{equation}
 x_{k+1}-x_k=-\tfrac12\beta(t_k)x_k\Delta t
 +\sqrt{\beta(t_k)}\sqrt{\Delta t}\,z_k.
\end{equation}
第一项尺度是 $\Delta t$，第二项是 $\sqrt{\Delta t}$；这两个尺度的差别决定了随机微积分。把极限写成\term{随机微分方程}{gl:sde}：
\begin{equation}
 \dd x_t=-\tfrac12\beta(t)x_t\dd t+\sqrt{\beta(t)}\dd W_t.
 \label{eq:vp-sde}
\end{equation}
$W_t$ 是标准\term{布朗运动}{gl:brownian}，独立增量 $W_{t+h}-W_t\sim\N(0,hI)$。$\dd W_t$ 不是一个普通函数的微分；它表示随机积分里的增量。给 SDE 写导数 $\dot x_t$ 通常没有意义。

更一般地，我们先限制到各向同性、仅依赖时间的扩散系数：
\begin{equation}
 \dd x_t=f(x_t,t)\dd t+g(t)\dd W_t.
 \label{eq:general-sde}
\end{equation}
$f$ 称为漂移，描述可预测的平均运动；$g$ 称为扩散强度，$g^2\Delta t$ 是小步噪声协方差。这里 $g$ 与图像生成模型的名字“Diffusion”相关，但不要把“diffusion coefficient”误译成网络预测的噪声值。

式\eqref{eq:vp-sde}称为方差保持（VP）SDE，因为标准高斯是其不变分布。这不意味着任何初始数据的方差在所有时刻都恒定。设
\begin{equation}
 B(t)=\int_0^t\beta(r)\dd r,\quad
 a(t)=e^{-B(t)/2},\quad b(t)^2=1-e^{-B(t)}.
\end{equation}
线性 SDE 的积分因子解给出
\begin{equation}
 x_t=a(t)x_0+a(t)\int_0^t a(r)^{-1}\sqrt{\beta(r)}\dd W_r.
\end{equation}
后面的随机积分为均值零高斯，其方差由 It\^o 等距得到
\begin{equation}
 a(t)^2\int_0^t a(r)^{-2}\beta(r)\dd r=1-a(t)^2.
\end{equation}
故条件边缘仍是 $\N(a(t)x_0,b(t)^2I)$。这恰连接到 DDPM 的累计信号，而不是一个与离散模型毫无关系的新对象。

另一类方差增长（VE）SDE 不缩小数据，而让噪声方差增加。若 $\sigma(0)=0$ 且 $\sigma^2$ 可微单调增长，可写
\begin{equation}
 \dd x_t=\sqrt{\frac{\dd\sigma(t)^2}{\dd t}}\dd W_t,
 \qquad x_t\mid x_0\sim\N(x_0,\sigma(t)^2I).
\end{equation}
实际方法常从小的正噪声 $\sigma_{\min}$ 开始；这时核的方差取决于初始分布是否已作平滑，必须明确。VP 通过信号衰减得到简单端，VE 通过噪声压过数据尺度得到近似简单端。两者是不同的路径设计，不是不同的网络层。

\subsection{从粒子方程推导密度方程}
取一个二阶可微的标量测试函数 $\varphi(x)$。It\^o 公式给出
\begin{equation}
 \dd\varphi(x_t)=\left(f\cdot\grad\varphi+\tfrac12g^2\Delta\varphi\right)\dd t
 +g\grad\varphi\cdot\dd W_t.
\end{equation}
取期望，随机积分均值为零，得到
\begin{equation}
 \frac{\dd}{\dd t}\int\varphi(x)p_t(x)\dd x
 =\int\left(f\cdot\grad\varphi+\tfrac12g^2\Delta\varphi\right)p_t(x)\dd x.
\end{equation}
对右边积分分部一次和两次，在边界项消失时，有
\begin{equation}\boxed{
 \partial_tp_t=-\divg(f p_t)+\tfrac12g(t)^2\Delta p_t.
 }\label{eq:fokker-planck}\end{equation}
这就是\term{Fokker--Planck 方程}{gl:fp}。SDE 描述一个随机粒子；Fokker--Planck 描述大量粒子的密度。拉普拉斯项让密度向四周扩散；漂移项把密度随向量场搬运。$1/2$ 来自 It\^o 二阶项，而不是任意选的缩放。

这个推导也解释了为什么随机微积分必须保留二阶项：布朗增量的平方量级是 $\Delta t$，与普通漂移增量一样重要。附录\ref{app:stochastic}会用有限步 Taylor 展开补齐这一点。

\section{反向 SDE：倒转随机过程需要 score}\label{sec:reverse-sde}
\subsection{先把时钟说清楚}
继续使用前向时间 $t$，但生成时令 $t$ 从 $1$ 走向 $0$。在前述正则性条件下，反向 SDE 为\cite{r_anderson,r_sde}
\begin{equation}\boxed{
 \dd x_t=\left[f(x_t,t)-g(t)^2s_t(x_t)\right]\dd t
 +g(t)\dd\overline W_t,\qquad \dd t<0.
 }\label{eq:reverse-sde}\end{equation}
$\overline W_t$ 是按反向时间计的布朗运动；它的小步方差是 $|\dd t|$，不是负数。反向 SDE 与前向 SDE 一般不共用同一条噪声实现。我们反转的是随机过程的分布规律，而不是逐次撤销同一组随机数。

若不喜欢负时间增量，设反向时钟 $r=1-t$，$y_r=x_{1-r}$。用正向增长的 $r$ 写就是
\begin{equation}
 \dd y_r=\left[-f(y_r,1-r)+g(1-r)^2s_{1-r}(y_r)\right]\dd r
 +g(1-r)\dd\widetilde W_r.
 \label{eq:reverse-clock}
\end{equation}
两式含义相同。看到论文中的 $f-g^2s$ 和 $-f+g^2s$ 时，先查时钟，再判断符号。生成小步 $h>0$ 的 Euler--Maruyama 形式为
\begin{equation}
 x_{t-h}=x_t-h f(x_t,t)+h g(t)^2s_\theta(x_t,t)
 +g(t)\sqrt h\,z.
 \label{eq:reverse-euler}
\end{equation}

\subsection{score 修正的来源}
为什么不能仅仅给漂移加一个负号？因为给定当前 $x$，过去的样本并不是在所有方向等概率。贝叶斯公式把前一时刻更常出现的位置赋予更大后验权重。

以纯加性的小步噪声为例，$X_t=X_{t-h}+g\sqrt h Z$。对 $X_t=x$，高斯去噪关系给出
\begin{equation}
 \E[X_{t-h}\mid X_t=x]=x+h g^2\grad_x\log p_t(x).
\end{equation}
再加入前向漂移的反向抵消，局部反向均值就是
$x-hf(x,t)+hg^2s_t(x)+o(h)$。局部协方差主导项仍是 $hg^2I$。这便揭示了反向漂移中 score 项的来源。一般情形的小步密度推导与矩阵扩散修正见附录\ref{app:proofs}。

仅检查一条密度方程，并不足以唯一确定“哪个 SDE 才是原过程的时间反转”，因为不同概率流可以给出相同的密度变化。反向 SDE 的严格结论还依赖逆向转移概率及其小步极限。这个区分会在概率流 ODE 中再次出现。

\subsection{与离散 DDPM 的局部联系}
对 VP SDE，$f=-\beta(t)x/2,g^2=\beta(t)$。反向小步均值为
\begin{equation}
 x+h\beta(t)\left(\tfrac12x+s_t(x)\right).
\end{equation}
DDPM 均值参数化用 $s_\theta=-\epsilon_\theta/\sqrt{1-\bar\alpha_k}$，于是
\begin{equation}
 \mu_\theta=\frac{x_k+\beta_k s_\theta(x_k,k)}{\sqrt{1-\beta_k}}
 =x_k+\beta_k\left(\tfrac12x_k+s_\theta\right)+O(\beta_k^2).
\end{equation}
两者一阶对应。这说明 DDPM 的高斯均值与连续反向扩散共享局部结构。但有限步方差选取、端点解码和时间离散化会使实际采样器不同；不能从一阶展开推出所有离散算法完全等价。

\section{概率流 ODE：把扩散项写成输运}\label{sec:pfode}
\subsection{概率流中的一半系数}
由于 $\grad p_t=p_ts_t$ 且 $g$ 不依赖状态，式\eqref{eq:fokker-planck}可以重写为
\begin{align}
 \partial_tp_t
 &=-\divg(f p_t)+\tfrac12g^2\divg(p_ts_t)\\
 &=-\divg\left[\left(f-\tfrac12g^2s_t\right)p_t\right].
\end{align}
而确定性 ODE $\dot x_t=v_t(x_t)$ 所诱导的密度满足\term{连续性方程}{gl:continuity}
\begin{equation}
 \partial_tp_t+\divg(p_tv_t)=0.
\end{equation}
因此选择
\begin{equation}\boxed{
 \frac{\dd x_t}{\dd t}=v_t^{\mathrm{PF}}(x_t)
 =f(x_t,t)-\tfrac12g(t)^2s_t(x_t)
 }\label{eq:pfode}\end{equation}
即可得到与原 SDE 相同的单时刻边缘密度，条件是精确 score、相同的初始边缘与适当的解唯一性。生成时将此 ODE 从 $t=1$ 积分到 $t=0$。这就是概率流（probability flow）ODE\cite{r_sde}。

反向 SDE 的 score 系数是 $g^2$，概率流 ODE 是 $g^2/2$。如果把反向 SDE 的随机项删掉，仍保留完整的 $g^2$，所得 ODE 一般不再具有所需边缘分布。删掉噪声必须同时调整漂移，因为随机扩散本身已经在搬运概率。

\begin{examplebox}{一条随机游走的过程与一条静止的过程}
设数据本来就是 $\N(0,I)$，并使用 VP SDE。每个时刻都保持 $\N(0,I)$，所以 $s_t(x)=-x$。概率流速度为
\begin{equation}
 -\tfrac12\beta x-\tfrac12\beta(-x)=0.
\end{equation}
ODE 样本一直停在最初随机抽到的位置。SDE 样本却不断随机运动。两者每个时刻的密度完全一致，联合轨迹分布显然不同。这是检验“同边缘不等于同路径”最简洁的例子。
\end{examplebox}

\subsection{连续性方程的几何意义}
对一个小区域 $A$，其概率质量随时间的变化是穿过边界的净流量：
\begin{equation}
 \frac{\dd}{\dd t}\int_Ap_t(x)\dd x
 =-\int_{\partial A}p_tv_t\cdot n\dd S.
\end{equation}
$p_tv_t$ 是概率流密度，$n$ 为边界外法向。散度表示单位体积的净流出量：$\divg(p_tv_t)>0$ 时该位置的密度下降。不能只看 $\divg v_t$ 判断固定位置密度，因为密度还会随位置搬运；沿轨迹的对数密度变化才由 $-\divg v_t$ 控制。

给定一条密度路径，速度未必唯一。在 $d\ge2$ 的区域，如果向量场 $w_t$ 满足 $\divg(p_tw_t)=0$，那么 $v_t+w_t$ 产生相同密度变化。例如径向对称密度下的绕中心旋转，可以改变粒子轨迹而不改变密度。概率流 ODE 提供一个由前向 SDE 自然导出的选择，并不是唯一可能的确定性生成机制。

\subsection{速度场与 score 为什么不是同一个概念}
Score 是密度对数的梯度；光滑情况下其旋度为零。速度场只需满足概率输运的要求，并不必是某个标量势函数的梯度。概率流速度包含前向漂移 $f$，若 $f$ 自身具有旋转部分，整体速度就可能不是梯度场。流匹配也不要求网络速度是保守场。

在不同噪声尺度或路径上把 score 转换成速度，需要知道日程系数；式\eqref{eq:pfode}和式\eqref{eq:gaussian-path-velocity}正是这样的转换。一个张量的形状相同，不能证明它们具有相同意义。$\epsilon_\theta$、$s_\theta$、$v_\theta$ 和 $u_\theta$ 都可以是 $d$ 维向量，区别在于监督、单位与放入生成方程的方式。

如果密度路径具有对称性，某些速度的旋转部分可能不改变任何单时刻密度，但会改变起点到终点的配对。这说明生成建模既有分布层面，也有输运几何层面。讨论编辑、反演与保持身份时，后一层往往变得特别重要；仅用“终点分布接近”描述这些任务并不充分。

\subsection{确定性为何仍然可以有多样性}
ODE 在给定初始点时通常是确定性的。但初始点从 $\N(0,I)$ 随机抽取，不同初始点被流映射到不同的最终点。整体分布是一个随机起点分布的\term{推前}{gl:transport}。随机性可以集中在起点，也可以沿 SDE 路径持续注入；两者都可以生成多样本。

若 ODE 在整个闭区间满足足够强的可逆流条件，光滑正密度不会在有限时间内精确坍缩到严格低维流形。因此现实生成中所谓“回到数据流形”通常是近似或正噪声平滑意义上的说法。对奇异数据分布，要么保留一个小噪声端点，要么允许端点速度失去正则性，不能同时宣称处处正则可逆与精确奇异终点。

\subsection{沿 ODE 计算似然}
将连续性方程展开为
\begin{equation}
 \partial_t p_t+v_t\cdot\grad p_t+p_t\divg v_t=0.
\end{equation}
沿满足 $\dot x_t=v_t(x_t)$ 的轨迹使用链式法则，得到
\begin{equation}
 \frac{\dd}{\dd t}\log p_t(x_t)=-\divg v_t(x_t).
\end{equation}
把数据点从 $0$ 积分到简单端 $1$，因此
\begin{equation}
 \log p_0(x_0)=\log p_1(x_1)+\int_0^1\divg v_t(x_t)\dd t.
 \label{eq:cnf-likelihood}
\end{equation}
这个正号可以通过“扩张会降低沿轨迹密度”核查。对神经网络速度场，散度是 Jacobian 的迹。高维下可以用满足 $\E[\xi\xi^\top]=I$ 的随机向量估计：
\begin{equation}
 \E_\xi\left[\xi^\top(\partial_xv_t)\xi\right]=\tr(\partial_xv_t).
\end{equation}
无需构造完整 $d\times d$ Jacobian，向量积可通过自动微分计算。有限次迹估计有随机误差，数值积分有离散误差；“连续流有精确变化量公式”与“实际计算完全无误差”要分开。

还有一处边界：代入近似 $s_\theta$ 之后，神经 ODE 自身定义了一个生成密度 $p_\theta^{\mathrm{ODE}}$，可以用该变化量公式计算它的似然。但该 ODE 与使用同一个近似 score 的 SDE 不必仍然共享精确边缘，也不必等于真实数据密度。理想等价关系不能替代训练误差分析。

\section{Flow Matching：直接学习搬运概率的速度}\label{sec:flowmatching}
\subsection{换一个问题：不先规定随机加噪链}
现在使用 $\tau\in[0,1]$，方向为噪声到数据。抽取 $Z\sim\pnoise$ 和 $X\sim\pdata$，先用独立\term{耦合}{gl:coupling} $\pi(z,x)=\pnoise(z)\pdata(x)$。定义线性插值
\begin{equation}
 Y_\tau=(1-\tau)Z+\tau X,\qquad \dot Y_\tau=X-Z.
 \label{eq:linear-interpolant}
\end{equation}
我们可以直接抽到 $Y_\tau$ 和对应的条件速度标签 $X-Z$，不需要数值求解一条当前网络轨迹。能否用这样的局部监督学到真正生成整个分布的 ODE？Flow Matching 的核心答案是可以\cite{r_fm}。

首先，不同 $(Z,X)$ 对的直线会相交。在相同位置 $y$ 和时刻 $\tau$，不同配对可能要求不同速度。单个网络 $u_\theta(y,\tau)$ 只能输出一个向量。因此它必须学到一个对相交路径进行平均的边缘速度：
\begin{equation}
 u_\tau(y)=\E[X-Z\mid Y_\tau=y].
 \label{eq:marginal-velocity}
\end{equation}
这个平均为何正确，不能只凭“看起来合理”接受。

\subsection{用测试函数证明边缘速度搬运了所需密度}
对紧支撑光滑测试函数 $\varphi$，在可交换微分与期望的条件下，
\begin{align}
 \frac{\dd}{\dd\tau}\E[\varphi(Y_\tau)]
 &=\E[\grad\varphi(Y_\tau)\cdot(X-Z)]\\
 &=\E[\grad\varphi(Y_\tau)\cdot u_\tau(Y_\tau)]\\
 &=\int\grad\varphi(y)\cdot u_\tau(y)p_\tau(y)\dd y\\
 &=-\int\varphi(y)\divg(p_\tau(y)u_\tau(y))\dd y.
\end{align}
所以 $p_\tau$ 以弱形式满足
\begin{equation}
 \partial_\tau p_\tau+\divg(p_\tau u_\tau)=0.
\end{equation}
若满足适当的正则性与唯一性，从 $Y_0\sim\pnoise$ 解 ODE $\dot y=u_\tau(y)$，就得到所需边缘路径，终点分布为 $\pdata$。端点可能需要取极限；弱形式允许我们先讨论随机变量的分布，而不强迫每条确定性条件路径拥有普通密度。

这个证明有一个很有力量的地方：ODE 的生成轨迹不必沿着训练时某个配对的直线。它只需要让所有粒子的分布在每个时刻正确。条件标签可能互相冲突，条件均值恰好化解了冲突，并保留整体概率守恒。

\subsection{为什么条件流匹配能训练边缘速度}
给网络使用可计算目标
\begin{equation}\boxed{
 L_{\mathrm{CFM}}=\E_{\tau,(Z,X)\sim\pi}
 \norm{u_\theta(Y_\tau,\tau)-(X-Z)}^2.
 }\label{eq:cfm}\end{equation}
假设 $\E\norm{X-Z}^2<\infty$，令 $V=X-Z$，再次使用条件投影恒等式：
\begin{align}
 L_{\mathrm{CFM}}
 &=\E_{\tau,Y_\tau}\norm{u_\theta(Y_\tau,\tau)-u_\tau(Y_\tau)}^2\\
 &\quad+\E_{\tau,Z,X}\norm{V-u_\tau(Y_\tau)}^2.
\end{align}
第二项不依赖参数。于是可观测的条件速度回归，与不可直接计算的边缘速度回归，具有相同最优解和梯度（在可交换微分与期望的条件下）。这与去噪 score matching 的数学骨架完全相同：用条件目标的样本监督，通过条件均值得到边缘对象。

\begin{intuitionbox}{两个无法直接获得的教师}
扩散无法直接告诉我们边缘密度的 score，但能提供条件核的 score；流匹配无法直接告诉我们整体边缘速度，但能提供配对插值的速度。平方损失使神经网络学会后验平均，而这一平均分别满足正确的密度梯度关系与连续性方程。这是“只凭样本训练生成规则”的共同桥梁。
\end{intuitionbox}

\subsection{耦合为什么改变难度，却不改变端点}
只要 $\pi$ 的两个边缘仍是 $\pnoise$ 和 $\pdata$，插值的端点就不变。但中间密度、路径交叉程度、条件标签方差和 ODE 轨迹都依赖 $\pi$。独立配对容易抽样，却可能使许多相反方向在同一位置竞争。

若选取或近似一个较好的输运耦合，可以减少不必要的长途搬运或冲突。但必须区分几种说法：“端点之间每条条件路径是直线”，不等于“边缘 ODE 轨迹是直线”，更不等于“已经找到全局最优输运”。最优输运还需要明确代价函数，并对所有具有指定边缘的耦合优化。Flow Matching 中某些条件路径的最优性质，不能无条件提升为全局边缘耦合最优性质。

Rectified Flow 与 reflow 会利用模型产生的噪声、数据配对重新训练速度场。它们怎样使路径趋向易于积分的形状，第三章会解释；这里要带走的是：耦合是一种训练自由度，与网络结构和数值求解器属于不同层面。

\subsection{高斯条件路径的通用速度}
回到数据到噪声的时钟 $t$，固定数据 $x_0$，采用 $x_t=a(t)x_0+b(t)Z$。条件曲线导数为 $a'x_0+b'Z$。在条件密度 $q_t(x\mid x_0)$ 中，$Z=(x-a x_0)/b$，故条件速度可写为
\begin{equation}
 w_t(x\mid x_0)=a'x_0+\frac{b'}b(x-a x_0),\quad b>0.
\end{equation}
边缘速度是 $w_t$ 的条件平均。利用式\eqref{eq:tweedie}和 $\E[Z\mid x_t=x]=-bs_t(x)$，得到
\begin{equation}\boxed{
 w_t(x)=\frac{a'}a x+\left(\frac{a'}a b^2-bb'\right)s_t(x),\quad a,b>0.
 }\label{eq:gaussian-path-velocity}\end{equation}
这个公式连接了 Gaussian Flow Matching 与 diffusion。若前向线性 SDE 的漂移为 $f(x,t)=\gamma(t)x$，其中 $\gamma=a'/a$，条件方差满足
\begin{equation}
 (b^2)'=2\gamma b^2+g^2,
\end{equation}
那么 $\gamma b^2-bb'=-g^2/2$，式\eqref{eq:gaussian-path-velocity}正是概率流 ODE 的速度。相同的概率路径可以通过 score 训练，也可以通过条件速度训练；训练参数化和损失权重仍需转换。

反过来，不是所有任意选择的 $a,b$ 都对应一个可用的前向加噪 SDE：计算出的 $g^2=(b^2)'-2(a'/a)b^2$ 必须非负，并满足相应正则性。流匹配不要求先找到这个非负扩散系数，因此可以使用更宽的路径家族。随机插值框架还把概率输运与扩散进一步联系起来\cite{r_stochasticinterpolants}。

\section{条件生成与引导：怎样改变概率方向}\label{sec:guidance}
\subsection{条件 score 的贝叶斯分解}
给定条件 $c$，例如类别或文本，直接训练条件模型学 $s_t(x\mid c)$ 是最自然的做法。若另有噪声状态上的分类器 $p_t(c\mid x)$，贝叶斯公式给出
\begin{equation}
 \grad_x\log p_t(x\mid c)=s_t(x)+\grad_x\log p_t(c\mid x).
 \label{eq:classifier-guidance}
\end{equation}
$\log p(c)$ 对 $x$ 的梯度为零。分类器必须在相应噪声分布上有意义；把只见过干净图像的分类器直接用于高噪声输入，不满足这个要求。

条件 score 已经包含某个条件的真实统计规律。为什么还要额外放大条件方向？真实数据里的图文配对可能不够紧，模型也有误差；实际用户往往希望图像更强地表达条件。于是引导允许生成过程偏离普通条件抽样，换取更明显的条件遵循。

\subsection{Classifier-free guidance 的数学形式}
训练时以一定概率丢弃条件，用同一网络学习条件预测和空条件预测。采用本书约定，令引导尺度 $w$ 满足
\begin{equation}\boxed{
 s^{\mathrm{CFG}}(x,t,c)
 =s_\theta(x,t,\varnothing)+w\left[s_\theta(x,t,c)-s_\theta(x,t,\varnothing)\right].
 }\label{eq:cfg}\end{equation}
$w=0$ 是无条件；$w=1$ 是普通条件；$w>1$ 外推条件方向\cite{r_cfg}。有的论文把 $s_c+\omega(s_c-s_u)$ 写成引导，这时 $w=1+\omega$。比较“引导值等于多少”之前，要先比较公式。

若两个 score 都精确，对于某个固定时刻，式\eqref{eq:cfg}形式上等于以下未归一化密度的 score：
\begin{equation}
 \widetilde p_t(x\mid c)\propto p_t(x\mid c)^w p_t(x)^{1-w}.
\end{equation}
利用贝叶斯也可写成 $p_t(x)p_t(c\mid x)^w$ 乘以与 $x$ 无关的因子。这给出“增强条件相关位置”的直觉，但需要右边可积；更重要的是，逐个时刻定义出的 $\widetilde p_t$ 一般不自动构成原前向 SDE 的一致边缘路径。因此不能据此声称最终生成密度严格等于 $p_0(x\mid c)^w/p_0(x)^{w-1}$。引导系统的全程分布还受动力学与数值积分影响。

在相同线性参数化下，噪声、$v$ 或流速度预测也可作同样的条件、空条件线性组合。不过对一般流匹配速度，不能直接将它解释为某个 log density 的梯度；该组合是改变输运场，不等于自动获得新的解析目标密度。

\begin{cautionbox}{更强条件与更好生成不是同一个量}
大引导可能增强文本吻合，也可能放大预测误差、削弱多样性或导致过饱和。理论上的条件梯度方向给出机制解释，却不保证任意尺度都稳定。第\ref{ch:applications}章还会区分语义条件、空间控制、编辑约束及其不同信息通道。
\end{cautionbox}

\section{把理想方程接回真实模型}
\subsection{把概率路径与样本路径再分开一次}
现在我们拥有三种很容易写成同一形式的对象：SDE 的随机样本路径、为了回归而构造的 $a(t)x_0+b(t)Z$ 条件曲线，以及生成 ODE 的积分轨迹。它们可以拥有相同的单时刻边缘密度，却有不同的联合分布和不同的端点耦合。

前向 SDE 的噪声在时间上持续注入；用同一份 $Z$ 构造条件曲线时，所有时刻的噪声完全相关；概率流 ODE 只在初始条件中使用随机性，之后按一个固定场移动。选择它们中的哪一种，不仅影响采样的外观，也影响你是否能把某个中间状态与某个原始端点唯一对应。

\begin{examplebox}{同一个标准高斯边缘也可以有不同耦合}
设 $Z,E$ 是独立标准高斯。过程 $X_t=Z$ 在所有时刻不动；过程 $Y_t=\cos(t)Z+\sin(t)E$ 在每个时刻也都是标准高斯。二者的边缘相同，但
$\E[X_sX_t]=1$，而 $\E[Y_sY_t]=\cos(t-s)$。如果只保存每个时刻的样本直方图，你看不出这种差别；若保存轨迹或配对，差别就清晰了。
\end{examplebox}

这个例子不是新的扩散采样器，只是说明边缘信息不足以确定路径。训练只要能够获得所需的单时刻监督，就可能使用不同的人工构造；生成是否正确则需要相应的动力学证明，不能只因边缘形式相像就把算法互换。

\subsection{终点与中间点的生成要求并不一样}
如果只要求最终分布正确，原则上允许很多不同中间路径。若还要求每个中间状态都具有指定噪声语义、能与一个既有 score 网络兼容，那么路径自由度就减少。一个预先训练的 DDPM 网络不能不经转换便在任意 Flow Matching 路径上使用，因为它学习的是自己训练日程里的条件统计。

同样，在同一条路径上更换求解器，通常可以复用网络，但仍需保持时间嵌入、输入缩放、预测目标和终点处理一致。将“训练规则允许不模拟当前网络轨迹”理解为“采样时也不需要积分”是另一种常见混淆。CFM 在训练时不需要反向穿过 ODE 求解器；生成时仍须求解学到的场，除非再训练一个一步或少步生成机制。

\subsection{一个最小的数学实现清单}
在真正写代码前，可以用六句话明确模型。数据在什么空间、采用什么尺度？概率路径的起点与终点是什么？随机训练样本怎样构造？网络预测的是哪个可观测条件标签？预测值如何转换成生成动力学？时间区间与端点怎样离散化？如果这些问题没有明确答案，增加复杂网络也只会把错误隐藏得更深。

以直线 CFM 为例：在选定的连续表示里，抽独立 $Z,X$；抽 $\tau$；形成 $Y_\tau$；回归 $X-Z$；从 $Z$ 解 $\dot Y=u_\theta(Y,\tau)$；采用合适时间网格到 $\tau=1$。以 VP score 生成而言：形成 $x_t=ax_0+b\epsilon$；回归 $-\epsilon/b$ 或等价噪声标签；明确 $f,g$；将预测代入反向 SDE 或概率流 ODE；从简单端向较小的 $t$ 积分。所有公式里相同的符号必须对应相同的日程与坐标尺度。

\subsection{四层误差分别来自哪里}
第一层是数据与表示：有限训练集、筛选偏差、量化像素、隐空间压缩决定了可学的对象。第二层是终点与路径：若噪声端仍残留数据结构，而生成起点采用标准高斯，就存在先验失配。第三层是统计学习：有限容量和优化使 $s_\theta$ 或 $u_\theta$ 偏离理想场。第四层是数值求解：有限步 ODE/SDE 近似产生额外偏差。

它们不会因为同一个公式推导正确而消失。模型提高容量未必修复一个错误日程；求解器增加阶数未必修复错误速度场；引导增强语义也可能加重积分难度。第三章的方法比较将沿这些轴线组织。

\subsection{为什么小的局部回归误差仍可能重要}
设真实 ODE 速度 $u$ 在状态上具有 Lipschitz 常数 $L$，估计场为 $u_\theta$。从相同起点出发，令轨迹误差 $e(\tau)=y_\theta(\tau)-y(\tau)$，有
\begin{align}
 \frac{\dd}{\dd\tau}\norm{e(\tau)}
 &\le L\norm{e(\tau)}+
 \norm{u_\theta(y_\theta(\tau),\tau)-u(y_\theta(\tau),\tau)}.
\end{align}
形式上应用 Gr\"onwall 不等式，得到
\begin{equation}
 \norm{e(1)}\le\int_0^1e^{L(1-r)}
 \norm{u_\theta(y_\theta(r),r)-u(y_\theta(r),r)}\dd r.
\end{equation}
它说明误差会沿路径累积，并可能被场的敏感性放大。但训练损失是在目标边缘上取平均，右边却沿模型轨迹计算；没有额外的覆盖、稳定性或一致误差条件，不能直接把一个小训练损失宣称为严格的小终点误差。这就是分布偏移在生成过程中的一种体现。

\subsection{一个统一图景}
\begin{figure}[htbp]\centering
\begin{tikzpicture}[>=Latex,
 box/.style={draw=teal!60,fill=teal!4,rounded corners=1pt,text width=61mm,align=center,inner sep=7pt,font=\small}]
\node[box] (samples) at (0,0) {数据样本与人为选定的概率路径};
\node[box] (score) at (-3.6,-2.5) {可计算条件 score\\平方回归得到边缘 score};
\node[box] (velocity) at (3.6,-2.5) {可计算条件速度\\平方回归得到边缘速度};
\node[box] (sde) at (-3.6,-5) {反向 SDE\\随机粒子动力学};
\node[box] (ode) at (3.6,-5) {生成 ODE\\确定性粒子动力学};
\node[box] (density) at (0,-7.5) {概率守恒方程\\目标边缘分布};
\draw[->,teal] (samples)--(score);\draw[->,teal](samples)--(velocity);
\draw[->,teal](score)--(sde);\draw[->,teal](velocity)--(ode);
\draw[->,gold](score.east)--node[above,font=\scriptsize]{概率流}(ode.north west);
\draw[->,teal](sde)--(density);\draw[->,teal](ode)--(density);
\end{tikzpicture}
\caption{本章的逻辑关系。score 与条件速度给出不同可训练入口；是否共享一条边缘路径，取决于路径和系数的匹配条件。}\label{fig:math-unified}
\end{figure}

\begin{table}[htbp]\centering\small
\begin{tabularx}{\linewidth}{p{24mm}YYY}\toprule
入口 & 可观测监督 & 理想边缘对象 & 生成规则\\\midrule
DDPM & 注入噪声 $\epsilon$ & $\E[\epsilon\mid x_k]$ & 学得的高斯逆转移\\
Score matching & 条件核梯度 & $\grad\log p_t$ & 反向 SDE 或概率流 ODE\\
Flow Matching & 条件路径导数 & $\E[\dot Y_\tau\mid Y_\tau]$ & 速度场 ODE\\\bottomrule
\end{tabularx}
\caption{同一个条件投影原理的三种入口。路径、目标与求解器的关系需要逐项检查。}
\end{table}

\section{练习与解答思路}
\begin{exercisebox}{2.1\quad 终点与信噪比}
设一维数据均值为 $2$、方差为 $4$，且 $\bar\alpha_K=0.01$。求 $x_K$ 的均值和方差；仅凭这两个量能否断定它是高斯？
\end{exercisebox}
\textbf{解答思路。}均值为 $0.2$，方差为 $0.01\times4+0.99=1.03$。不能断定高斯；高斯混合与单高斯可以拥有相同均值、方差。终点近似还需要考虑高阶结构与所用误差尺度。

\begin{exercisebox}{2.2\quad 条件回归的不可消除误差}
给定 $U$，监督 $V$ 以相同概率等于两个向量 $v_1,v_2$。求平方回归最优输出及最小条件损失。解释它与流匹配相交直线的关系。
\end{exercisebox}
\textbf{解答思路。}最优输出为 $(v_1+v_2)/2$，最小损失为 $\norm{v_1-v_2}^2/4$。边缘速度可以正确，条件标签损失却不为零；交叉的训练直线无需由网络逐条复现。

\begin{exercisebox}{2.3\quad 检查概率流系数}
数据为标准高斯，前向 SDE 为 $\dd x=-\beta x\dd t/2+\sqrt\beta\dd W$。若生成时直接删掉反向 SDE 的噪声，保留其漂移，得到的 ODE 是否保持标准高斯？
\end{exercisebox}
\textbf{解答思路。}真实 score 为 $-x$。反向时钟下删噪后的漂移为 $-\beta x/2$，样本向零收缩，方差不恒定。概率流需要 $g^2/2$ 的系数，才得到零速度。

\begin{exercisebox}{2.4\quad 直线流匹配的高斯解析解}
设 $Z\sim\N(0,I)$，$X\sim\N(m,cI)$，两者独立。求 $Y_\tau$ 的均值、方差，以及边缘速度 $\E[X-Z\mid Y_\tau=y]$。
\end{exercisebox}
\textbf{解答思路。}均值 $\tau m$，方差 $h(\tau)I$，其中 $h=(1-\tau)^2+\tau^2c$。联合高斯条件公式给出
\begin{equation}
 u_\tau(y)=m+\frac{\tau c-(1-\tau)}{h(\tau)}(y-\tau m)
 =m+\frac{h'(\tau)}{2h(\tau)}(y-\tau m).
\end{equation}
每条 ODE 轨迹满足 $y_\tau=\tau m+\sqrt{h(\tau)}z$。它通常不是训练配对的直线，却具有完全正确的边缘。这是可以直接代入验证的解析例子。

\begin{exercisebox}{2.5\quad 参数化改变了什么}
当 $a^2+b^2=1$ 时，证明从 $v$ 误差得到的原图误差为 $-b$ 倍，噪声误差为 $a$ 倍。讨论 $a\to0$ 与 $b\to0$ 时的数值意义。
\end{exercisebox}
\textbf{解答思路。}由式\eqref{eq:vconversion}立即得到。两个转换系数有界，而噪声、原图之间直接相互转换含 $a/b$ 或 $b/a$。这只说明转换的条件性更好，不构成任意训练设置下 $v$ 一定最佳的证明。

\chaptertakeaway{扩散模型的数学并非只有“加噪，再减噪”。高斯条件核提供可计算监督；条件期望把监督转化成边缘 score 或边缘速度；Fokker--Planck 与连续性方程解释这些场如何搬运整个概率分布。反向 SDE 和概率流 ODE 共享理想边缘而不共享轨迹；流匹配把相同的条件投影思想推广到直接设计的输运路径。下一章将用这张关系网理解具体算法，而不是记忆互不相连的名字。}
}{}
\IfFileExists{chapters/03-methods.tex}{\chapter{从数学到算法：扩散与流的设计谱系}
\label{ch:methods}

\readingguide{第二章给出了分布演化的语言，本章把它变成可以训练、可以采样、可以比较的算法。阅读时请始终追问四件事：选了什么概率路径，网络预测什么，怎样走完生成路径，网络用什么结构。前两件决定学习任务，第三件决定数值误差与速度，第四件决定表达能力和计算成本。把四件事分开，许多看似互相取代的方法就会成为可以组合的部件。}

\section{先画地图：方法名称为什么容易让人混淆}
\label{sec:m-method-map}

当我们说“用 DDPM 生成图片”，可能同时指一个逐步加噪的概率模型、一种噪声回归损失、一个高斯反向采样器，甚至默认了一套 U-Net。这些部件在最初的实现中一起出现，却没有数学理由要求它们永远绑在一起。同一个噪声预测网络可以用 DDPM、DDIM 或合适的 ODE 求解器采样；同一种 Flow Matching 损失也可以配合卷积网络或者 Transformer。要理解发展谱系，先拆掉名称造成的捆绑。

\begin{table}[htbp]
\centering
\caption{阅读生成方法时的四个独立问题。例子表示选择，不表示互斥的版本。}
\label{tab:m-four-axes}
\small
\begin{tabularx}{\linewidth}{p{0.13\linewidth}XX}
\toprule
设计维度 & 真正需要回答的问题 & 常见选择及影响\\
\midrule
概率路径 & 中间分布如何连接数据与噪声？ & VP、VE、高斯插值、直线插值；影响信噪比与运输几何\\
监督目标 & 网络输出哪个数学量？ & 噪声、数据、score、$v$、速度、端点；影响数值尺度与损失权重\\
采样算法 & 从起点怎样到达终点？ & 祖先采样、DDIM、Euler、Heun、DPM-Solver；影响离散误差与网络调用数\\
网络结构 & 用什么函数族表示输出？ & U-Net、DiT、多模态 Transformer；影响容量、归纳偏置与成本\\
\bottomrule
\end{tabularx}
\end{table}

第二章的统一线索是：已知的条件扰动产生训练样本，条件回归的最优解给出边缘分布所需的场，采样器再利用这个场移动样本。算法的差异来自在哪一环改变了设计。DDPM 把链式概率模型和去噪回归连起来；DDIM 改变跨时刻的联合构造；Score SDE 用连续时间统一随机与确定性生成；EDM 整理尺度和数值选择；Flow Matching 直接学习运输速度；蒸馏则把一段昂贵的运动压缩成较少的网络调用。

\begin{figure}[htbp]
\centering
\begin{tikzpicture}[font=\small,
  every node/.style={align=center},
  component/.style={draw=black!45,rounded corners,fill=black!3,
    minimum width=3.1cm,minimum height=0.85cm},
  >=stealth]
\node[component] (path) at (0,0) {概率路径\\$p_t$ 或 $p_\tau$};
\node[component] (target) at (4.4,0) {条件回归\\$\eps_\theta,s_\theta,u_\theta$};
\node[component] (sample) at (8.8,0) {生成算法\\SDE、ODE、端点映射};
\node[component] (net) at (4.4,-1.7) {函数表示\\U-Net、Transformer};
\draw[->] (path) -- node[above=5mm]{定义样本与标签} (target);
\draw[->] (target) -- node[above=5mm]{提供场或映射} (sample);
\draw[->] (net) -- node[right]{承载参数 $\theta$} (target);
\node at (0,-1.7) {DDIM：改变联合构造\\EDM：整理尺度与调度};
\node at (8.8,-1.7) {求解器：减少数值误差\\蒸馏：学习有限步结果};
\end{tikzpicture}
\caption{本章的原创方法地图。网络结构位于函数表示层，不能与采样器排在同一条替代链上。}
\label{fig:m-map}
\end{figure}

下面采用全书约定：离散时刻 $k$ 从数据 $x_0$ 走向噪声 $x_K$；连续扩散时间 $t$ 也从数据走向噪声，生成时让 $t$ 下降；流时间 $\tau$ 从噪声走向数据，生成时让 $\tau$ 上升。改变时钟并不自动改变路径，改变路径也不等于只把时钟反过来。

\section{DDPM：把反向高斯链变成监督学习}
\label{sec:m-ddpm}

\subsection{训练并不需要模拟整条加噪链}

DDPM 的实用突破，是把复杂的反向分布学习转换为一个容易造标签的任务：给网络一张含噪样本及其噪声级别，请它预测加入的噪声。原论文以变分模型为起点，并展示了去噪目标与高质量图像生成的联系。\cite{r_ddpm}

回顾式~\eqref{eq:forward-marginal}，令
\[
 a_k=\sqrt{\bar\alpha_k},\qquad b_k=\sqrt{1-\bar\alpha_k},
 \qquad x_k=a_kx_0+b_k\eps,\quad \eps\sim\N(0,I).
\]
这里 $a_k,b_k$ 是便于本章讨论的缩写，不改变第二章符号。由于边缘分布可以直接采样，一次训练只需随机选择一个 $k$，无需先计算 $x_1,x_2,\ldots,x_{k-1}$。网络输入的时间必须与 $a_k,b_k$ 对应，否则同一个含噪向量无法让网络判断“应去除多少噪声”。

\begin{examplebox}{DDPM 的一轮训练}
\begin{enumerate}[leftmargin=*,itemsep=2pt]
\item 从数据集中抽取 $x_0$；从 $\{1,\ldots,K\}$ 抽取 $k$；抽取独立的 $\eps\sim\N(0,I)$。
\item 直接构造 $x_k=a_kx_0+b_k\eps$。
\item 计算预测 $e=\eps_\theta(x_k,k,c)$；无条件模型省略条件 $c$。
\item 计算 $\ell=\norm{e-\eps}^2$，对 $\theta$ 求梯度并更新。
\end{enumerate}
批处理时，每个样本可以拥有自己的 $k$。网络共享参数，时间嵌入把不同噪声级别的任务区分开。
\end{examplebox}

随机的噪声标签并不意味着学习没有方向。固定 $x_k=x$ 后，平方损失的最优预测为
\begin{equation}
\eps^*(x,k)=\E[\eps\mid x_k=x].
\label{eq:m-noise-conditional}
\end{equation}
这正是第二章条件回归原理的应用。网络不会知道此次加噪使用的真实 $x_0$；它学到的是所有可能解释中由数据分布决定的条件平均。这个平均噪声可转换成 score，并指示哪些移动能把当前样本推向更高的含噪数据密度。不要把训练标签的可见性误解为生成阶段也能获得原图。

\subsection{预测噪声怎样变成一步生成}

由噪声预测可先计算一个数据估计
\begin{equation}
\widehat x_0(x_k,k)=\frac{x_k-b_k\eps_\theta(x_k,k)}{a_k}.
\label{eq:m-x0-est}
\end{equation}
将它代入式~\eqref{eq:posterior} 的后验均值，就得到常用反向均值
\begin{equation}
\mu_\theta(x_k,k)=\frac{1}{\sqrt{\alpha_k}}
\left(x_k-\frac{\beta_k}{b_k}\eps_\theta(x_k,k)\right).
\label{eq:m-ddpm-mean}
\end{equation}
采样器选择反向方差 $\sigma_k^2$，执行
\begin{equation}
x_{k-1}=\mu_\theta(x_k,k)+\sigma_k z_k,\qquad
z_k\sim\N(0,I).
\label{eq:m-ddpm-step}
\end{equation}
一般展示图像时，最后一步可使用均值输出，令 $z_1=0$；若讨论离散像素的精确似然，则要另外指定末端观测分布，不能以这条展示规则替代。

式~\eqref{eq:m-ddpm-step} 叫\emph{祖先采样}，意思是先从最上层噪声变量采样，再按照模型条件分布逐层采样后代。这一名称与祖先的时间先后没有关系。开始时取 $x_K\sim\N(0,I)$，要求 $A_K=\bar\alpha_K$ 足够小，使真实终端边缘接近这一先验；然后按 $k=K,\ldots,1$ 循环调用网络。输入已经没有原图，但每一步仍可以计算 $\mu_\theta$，这就是训练任务与生成任务接上的地方。

\begin{derivationbox}{后验均值为什么能用数据估计替换}
对 $k>1$，后验均值可写为
\[
 \widetilde\mu_k(x_k,x_0)
 =\frac{a_{k-1}\beta_k}{1-\bar\alpha_k}x_0
 +\frac{\sqrt{\alpha_k}(1-\bar\alpha_{k-1})}{1-\bar\alpha_k}x_k.
\]
它对 $x_0$ 是线性的，所以取条件期望时
\[
 \E[\widetilde\mu_k(x_k,x_0)\mid x_k]
 =\widetilde\mu_k(x_k,\E[x_0\mid x_k]).
\]
因此理想回归给出了真实反向条件分布的均值。但真实
$q(x_{k-1}\mid x_k)$ 还混合了不同 $x_0$ 的不确定性，一般并非一个精确高斯。均值正确不等于整个有限步反向分布都正确。
\end{derivationbox}

这也解释了“小步”的作用。用许多小噪声步，单步反向分布更容易用高斯近似。将 1000 步链随意删成 10 步，再沿用原来的单步 $\beta_k$，会使均值和方差不再对应新的跨度。正确的跨步祖先采样必须重新计算有效噪声参数，例如从 $k$ 跨到 $j<k$ 时，前向缩放比是 $a_k/a_j$，有效方差是 $1-a_k^2/a_j^2$。DDIM 与 ODE 求解器则提供了另一类跨步规则。

\subsection{实现前必须核对的尺度}

“噪声预测的损失接近 1”不是普遍的健康标准。损失是按维度求和还是求平均，数据值域是 $[0,1]$ 还是 $[-1,1]$，潜变量是否乘了缩放因子，都会改变数值。可操作的核对是：先计算恒为零的预测器的平均噪声均方误差，再比较训练模型能否显著降低它；按 $k$ 分组观察，避免容易的时刻掩盖困难的时刻。

另一个值得注意的现象是，式~\eqref{eq:m-x0-est} 在 $a_k$ 很小时会放大噪声预测误差。设预测误差为 $\delta e$，就有
\[
 \delta\widehat x_0=-\frac{b_k}{a_k}\delta e.
\]
高噪声区从 $\eps$ 转换到数据会病态；低噪声区从数据转换到 score 也可能病态。模型的输出参数化需要考虑这些尺度，而不是只比较输出名字。

\section{DDIM：保留单时刻训练，重新安排跨时刻运动}
\label{sec:m-ddim}

\subsection{相同边缘分布可以有不同联合分布}

DDPM 的去噪损失只需要 $(x_0,x_k)$ 的分布，它没有直接询问 $x_k$ 和 $x_j$ 在同一条随机路径中怎样相关。DDIM 利用这个自由度构造一族非马尔可夫的前向联合分布，保留 $q(x_k\mid x_0)$，并据此设计更短的生成过程。\cite{m_ddim}

用 $A_k=\bar\alpha_k$ 暂时简化公式，选定两个时刻 $j<k$，规定
\begin{align}
q_\sigma(x_j\mid x_k,x_0)
 &=\N\bigl(m_{j\mid k}(x_k,x_0),\,\sigma_{k\to j}^2I\bigr),\\
m_{j\mid k}(x_k,x_0)
 &=\sqrt{A_j}x_0
 +\sqrt{1-A_j-\sigma_{k\to j}^2}\,
 \frac{x_k-\sqrt{A_k}x_0}{\sqrt{1-A_k}}.
\label{eq:m-ddim-conditional}
\end{align}
要求 $0\leq\sigma_{k\to j}^2\leq1-A_j$。给定 $x_0$，分式中的残差服从标准正态；再加独立噪声后，$x_j$ 的总条件方差为
\[
 (1-A_j-\sigma_{k\to j}^2)+\sigma_{k\to j}^2=1-A_j,
\]
均值仍是 $\sqrt{A_j}x_0$。从较大时刻向较小时刻递归，就能保持所需的每个边缘分布。

保持这些边缘，意味着可以继续使用相同的去噪样本和标签，不意味着不同联合构造的完整变分目标数值完全相等。相应变分目标会得到不同的时刻权重；若允许每个时刻拥有充分表达能力，正权重不改变各自的理想条件回归解。实际网络共享参数且容量有限，权重仍会影响优化。DDIM 复用 DDPM 的常用去噪代理目标，应在这个意义下理解“相同训练”。

这里的“非马尔可夫”指所设计的前向推断联合分布：它依赖 $x_0$，不再仅由上一时刻决定。实际生成时用模型估计替换 $x_0$ 后，更新仍可以只依赖当前状态与时刻。两句话讨论不同的联合分布，不能据“非马尔可夫”断言生成代码必须保存所有历史图像。

\subsection{三项更新各自做什么}

令 $e_k=\eps_\theta(x_k,k)$，用式~\eqref{eq:m-x0-est} 得到 $\widehat x_0$，式~\eqref{eq:m-ddim-conditional} 给出
\begin{equation}
x_j=\sqrt{A_j}\widehat x_0
 +\sqrt{1-A_j-\sigma_{k\to j}^2}\,e_k
 +\sigma_{k\to j}z,\qquad z\sim\N(0,I).
\label{eq:m-ddim-step}
\end{equation}
第一项保留当前预测的数据成分，第二项保留一部分已估计的噪声方向，第三项引入新的随机性。这样看，采样器是在重新分配数据、已有噪声和新噪声的比例，而不是简单“少做几次去噪”。

常用选择是
\begin{equation}
\sigma_{k\to j}=
\eta\sqrt{\frac{1-A_j}{1-A_k}}
\sqrt{1-\frac{A_k}{A_j}},\qquad 0\leq\eta\leq1.
\label{eq:m-ddim-eta}
\end{equation}
$\eta=0$ 时，给定起点和条件，整个采样过程是确定性的；$\eta=1$ 对应所选跨度的后验方差。在相邻步及相应均值、方差选择下，它恢复 DDPM 的一种祖先采样形式；不能因此说它恢复所有采用不同反向方差的 DDPM 实现。对于 $j=0$，$A_0=1$，更新退化为输出数据估计。

确定性不妨碍生成多样性。每次从不同 $x_K$ 开始，就会得到不同输出。它只是去掉了路径中的额外随机数，因此同一随机起点更容易在不同步数的采样结果中保持相关语义。这个性质适合编辑与比较实验，但不保证任意改动采样网格后图像完全一致。

\begin{examplebox}{跨步 DDIM 的采样流程}
选定严格递减网格 $k_M=K>k_{M-1}>\cdots>k_0=0$，从
$x_{k_M}\sim\N(0,I)$ 开始。对 $i=M,\ldots,1$：
\begin{enumerate}[leftmargin=*,itemsep=2pt]
\item 令 $k=k_i,\ j=k_{i-1}$，在原训练时刻 $k$ 计算 $e_k$。
\item 从原训练噪声表取得 $A_k,A_j$，计算 $\widehat x_0$。
\item 用式~\eqref{eq:m-ddim-eta} 计算跨度方差，再用式~\eqref{eq:m-ddim-step} 更新。
\end{enumerate}
训练时刻仍使用原来的索引；不能把第 $i$ 次网络调用直接当作第 $i$ 个训练时刻。
\end{examplebox}

\subsection{DDIM 与概率流 ODE：关系需要一个条件}

确定性 DDIM 确实与概率流 ODE 有紧密关系，但“DDIM 就是任意一步的精确概率流解”会造成误解。看清它的数值坐标即可消除这个混淆。令
\[
 r(t)=\frac{b(t)}{a(t)},\qquad y(t)=\frac{x_t}{a(t)}.
\]
在相应高斯扩散路径和噪声预测转换下，概率流 ODE 可以写成
\begin{equation}
\frac{\dd y}{\dd r}=\eps_\theta(a(t)y,t).
\label{eq:m-ddim-rescaled-ode}
\end{equation}
确定性 DDIM 的更新恰好是
\[
 y_j=y_k+(r_j-r_k)\eps_\theta(x_k,k),
\]
即在缩放状态 $y$、噪声比 $r$ 坐标中冻结当前噪声预测的一阶更新。连续极限对应同一类概率流方程；有限步时，沿轨迹的网络输出会变化，冻结它就产生离散误差。原始时间 $t$ 中的朴素 Euler 更新与这个坐标变换后的更新也一般不同。

\begin{cautionbox}{边缘一致性不等于有限步精确生成}
DDIM 构造保证的是已知真实 $x_0$ 时前向条件边缘的正确性。生成时用条件平均预测替换 $x_0$，再跨越有限时间，既有函数近似误差，也有有限步运动误差。即使回归网络理想，粗网格也不能自动获得任意数据分布的精确样本。概率流 ODE 的同边缘结论还需要真实 score 与准确积分。
\end{cautionbox}

一个简单反例是全步跳到数据端：$j=0$ 时输出 $\E[x_0\mid x_K]$。如果 $x_K$ 几乎与数据独立，这个条件平均趋近数据均值；它当然不是完整的数据分布。理解这一点，也就理解了为什么“标签为噪声，误差很小”仍不足以支持一步生成。

\section{Score SDE：同一个 score 支持随机和确定性采样}
\label{sec:m-score-sde}

Score SDE 的贡献是把离散去噪链与多噪声级别 score 模型放在连续随机过程框架中，并给出反向 SDE、预测器--校正器和概率流 ODE 的统一关系。\cite{r_sde} 数学关系已在第~\ref{ch:math}~章推导，本节只讨论如何落实。

以标量扩散系数的前向过程
\[
 \dd x_t=f(x_t,t)\dd t+g(t)\dd W_t
\]
为例。设 $h>0$，从当前扩散时刻 $t$ 生成到 $t-h$。反向 SDE 的 Euler--Maruyama 预测步为
\begin{equation}
x^- = x_t-\bigl[f(x_t,t)-g(t)^2s_\theta(x_t,t)\bigr]h
      +g(t)\sqrt h\,z,\qquad z\sim\N(0,I).
\label{eq:m-em}
\end{equation}
公式中的负号来自时间下降，噪声方差仍为 $g(t)^2h$。直接把 $h$ 换成负数放进平方根，是常见的反向时间错误。\term{布朗运动}{gl:brownian}的增量依赖经过的正时间长度，而不是依赖时刻标签的符号。

到达新时刻后，可以固定这个噪声级别，做若干 Langevin 校正：
\begin{equation}
x\leftarrow x+\delta\,s_\theta(x,t-h)+\sqrt{2\delta}\,z'.
\label{eq:m-langevin}
\end{equation}
预测器负责沿时间演化，校正器尝试改善当前时刻的分布匹配。若 score 真实、步长适当、校正充分，Langevin 动力学以该时刻分布为平稳分布。有限次大步校正并不具有这个精确保证；额外调用网络也会增加代价。

因此，PC 采样适合从“演化”和“局部分布修正”两个角度理解。它不是对图像反复锐化，也不是每步都必须有校正。概率流 ODE 则用
$f-\frac12g^2s_\theta$ 作确定性漂移，在相同 score 理想条件下具有相同的单时刻边缘。SDE 保留路径随机性，ODE 可做确定性运输和连续流似然计算，二者的有限步效果取决于模型误差与数值方法，不能只凭“随机”或“确定”宣布优劣。

VP 和 VE 的区别也应放在路径轴上理解。VP 让数据系数下降、噪声系数增加，标准形式满足 $a(t)^2+b(t)^2=1$；VE 保持数据系数为 1，增加噪声标准差，因此样本方差变大。VP 适合直接连接标准化数据与近似标准高斯，VE 的加性噪声坐标便于讨论去噪器和噪声尺度。换路径后必须重新检查网络输入尺度与 score 转换，不能只替换采样器名称。

\section{Improved DDPM 与 EDM：让学习和数值尺度配合}
\label{sec:m-design}

\subsection{Improved DDPM：噪声日程与反向方差也是模型设计}

Improved DDPM 展示了噪声日程、学习反向方差和目标权重对生成质量、似然及采样效率的作用。\cite{m_improved} 这些改动没有要求抛弃 DDPM 的基本回归思想。

首先，逐步线性增加 $\beta_k$ 不代表信息逐步线性减少。噪声日程最好在累积量 $A_k$ 或\term{信噪比}{gl:snr}上检查，因为信号强度是多个缩放系数的乘积。一个常见余弦日程写成
\[
 A_k=\frac{F(k/K)}{F(0)},\qquad
 F(u)=\cos^2\left(\frac{u+s}{1+s}\frac{\pi}{2}\right),
\]
再由 $\beta_k=1-A_k/A_{k-1}$ 得到单步噪声。实际末端通常还要裁剪 $\beta_k$，避免奇异值。其意义是把信息变化安排得更平缓，让较多训练时刻保持有用信号；余弦函数并不是所有数据和所有路径上的普遍最优定律。

其次，反向方差可以由网络学习。对 $k>1$，已知后验方差
\[
 \widetilde\beta_k=
 \beta_k\frac{1-A_{k-1}}{1-A_k}.
\]
可将网络给出的插值量 $\rho_\theta$ 放入对数方差
\[
 \log \Sigma_\theta=
 \rho_\theta\log\beta_k+
 (1-\rho_\theta)\log\widetilde\beta_k,
\]
逐维应用以形成对角协方差。末端 $\widetilde\beta_1=0$ 必须专门处理。学习方差的动机是补充均值之外的反向不确定性，但仅靠噪声均方误差无法训练这一输出，需要变分目标或其混合形式。

混合目标可抽象写为
$L_{\mathrm{simple}}+\lambda L_{\mathrm{VLB}}$。这里 $L_{\mathrm{simple}}$
强调去噪回归，变分项补充方差与似然信息。某些实现对变分项中的均值分支停止梯度，以减少两个目标对均值学习的竞争；停止梯度是优化选择，不能说它仍在对完整原目标做同样的梯度下降。

若随机抽时刻的分布为 $q_k$，想无偏估计
$\frac1K\sum_k L_k$，需用 $L_k/(Kq_k)$。改变抽样概率和改变目标权重是两件事。前者配合正确的重要性权重可以降低估计方差；若不补偿，训练的目标本身已经变化。工程中二者都可能有用，但报告实验时应说明采用哪一种。

\subsection{EDM：把噪声标准差当作清楚的物理量}

EDM 系统拆分采样日程、数值积分、网络预条件和训练噪声选择，强调不同设计可以分别改进和组合。\cite{m_edm} 本节用其常用加性噪声坐标解释：设
\[
 x_\sigma=X+\sigma\eps,\qquad
 D_\theta(x,\sigma)\approx\E[X\mid x_\sigma=x].
\]
$\sigma$ 表示每个维度的噪声标准差，生成时从较大的 $\sigma$ 走向 0。此处 $\sigma$ 是噪声坐标，连续扩散 $t$ 的方向约定仍然不变。由去噪与 score 的关系
\[
 s_\theta(x,\sigma)=\frac{D_\theta(x,\sigma)-x}{\sigma^2},
\]
相应确定性 ODE 在这个坐标中为
\begin{equation}
\frac{\dd x}{\dd\sigma}
=\frac{x-D_\theta(x,\sigma)}{\sigma}.
\label{eq:m-edm-ode}
\end{equation}
可以把它理解为：含噪样本与当前数据估计的差值，按照当前噪声尺度转换成运动方向。降 $\sigma$ 时步长为负，因此会朝去噪估计移动。

问题在于 $\sigma$ 可能跨越多个数量级。若直接把所有 $x_\sigma$ 输入同一个网络，高噪声输入幅度大、低噪声残差小，网络要同时处理完全不同的尺度。预条件把已知的尺度变化放到解析系数中：
\begin{equation}
D_\theta(x,\sigma)=c_{\mathrm{skip}}(\sigma)x+
c_{\mathrm{out}}(\sigma)
F_\theta(c_{\mathrm{in}}(\sigma)x,c_{\mathrm{noise}}(\sigma)).
\label{eq:m-edm-precond}
\end{equation}
设数据标准差的代表尺度为 $\sigma_{\mathrm{data}}$，一种常用选择为
\begin{align}
c_{\mathrm{in}}&=(\sigma^2+\sigma_{\mathrm{data}}^2)^{-1/2},
&
c_{\mathrm{skip}}&=\frac{\sigma_{\mathrm{data}}^2}
{\sigma^2+\sigma_{\mathrm{data}}^2},\\
c_{\mathrm{out}}&=\frac{\sigma\sigma_{\mathrm{data}}}
{\sqrt{\sigma^2+\sigma_{\mathrm{data}}^2}},
&
c_{\mathrm{noise}}&=\tfrac14\log\sigma.
\end{align}
输入归一化让网络看到的幅度更稳定；跳连在低噪声区保留输入；输出系数让网络主要拟合尺度适当的残差。这些解析系数并不让一般数据分布的去噪器变成线性函数，而是减轻网络不必要的尺度负担。

\begin{intuitionbox}{预条件像换一把合适的尺子}
同一个温度预测问题，分别用摄氏度、毫摄氏度和绝对温标表示，并不会改变物理对象，却会改变数值优化。扩散的不同噪声时刻也有类似问题。预条件利用已知的噪声尺度调整输入和输出，让网络把容量花在数据结构上。
\end{intuitionbox}

相应损失常写为
\[
 \E_{\sigma,X,\eps}
 \left[w(\sigma)\norm{D_\theta(X+\sigma\eps,\sigma)-X}^2\right],
 \qquad
 w(\sigma)=\frac{\sigma^2+\sigma_{\mathrm{data}}^2}
 {(\sigma\sigma_{\mathrm{data}})^2}.
\]
这个权重等于 $c_{\mathrm{out}}^{-2}$，配合重写后的网络残差目标，使不同噪声级别更容易取得相近尺度。噪声抽样常在 $\log\sigma$ 空间安排，以覆盖多个数量级；具体分布仍需和数据、网络以及采样区间匹配。

数值方面，给出 $N$ 个正噪声节点
\[
 \sigma_i=\left[
 \sigma_{\max}^{1/\rho}
 +\frac{i}{N-1}\left(\sigma_{\min}^{1/\rho}
 -\sigma_{\max}^{1/\rho}\right)\right]^\rho,
 \qquad i=0,\ldots,N-1,
\]
再附加 $\sigma_N=0$。它在变换坐标中均匀取点，以把节点集中到需要精细分辨的噪声区。常见的 $\rho$ 数值是经验选择，不应脱离模型一律使用。生成通常从 $x_{\sigma_{\max}}=\sigma_{\max}z$、$z\sim\N(0,I)$ 开始；这是大噪声端的近似先验，其有效性也依赖噪声尺度足以盖过数据尺度。

记 $d_i=(x_i-D_\theta(x_i,\sigma_i))/\sigma_i$，
一步 Euler 先得到 $x'=x_i+(\sigma_{i+1}-\sigma_i)d_i$。
若 $\sigma_{i+1}>0$，再计算
$d'=(x'-D_\theta(x',\sigma_{i+1}))/\sigma_{i+1}$，用 Heun 更新
\begin{equation}
x_{i+1}=x_i+\frac{\sigma_{i+1}-\sigma_i}{2}(d_i+d').
\label{eq:m-heun}
\end{equation}
最后到 0 的一步可只用 Euler，避免除以 0。EDM 还讨论可控地增加噪声再继续生成的随机采样选择；解析 ODE、网络预条件和随机修正是不同部件，使用其中一个不代表必须同时使用全部。

\section{输出参数化：噪声、数据、score 和 \texorpdfstring{$v$}{v}}
\label{sec:m-parameterization}

第二章已经给出了参数化转换。本节强调它们在训练中的差异。固定 $t$，设 $x_t=aX+b\eps$，$a,b>0$。如果一个估计器给出 $\widehat\eps$，就可以得到
$\widehat X=(x_t-b\widehat\eps)/a$ 和
$\widehat s=-\widehat\eps/b$。反向也成立。对于无限容量、精确回归与相应权重，很多表示可描述同一个理想场；对于有限网络、有限计算与未调整的损失，它们不再等价。

\begin{table}[htbp]
\centering
\caption{相同含噪样本的几种输出。表中的误差关系按兼容的线性转换计算。}
\label{tab:m-targets}
\small
\begin{tabularx}{\linewidth}{p{0.12\linewidth}XX}
\toprule
输出 & 训练标签及理想含义 & 尺度提醒\\
\midrule
噪声 & $\eps$；最优为 $\E[\eps\mid x_t]$ & 数据误差为 $-(b/a)\delta\eps$，高噪声区放大\\
数据 & $X$；最优为 $\E[X\mid x_t]$ & 噪声误差为 $-(a/b)\delta X$，低噪声区放大\\
score & 条件标签 $-\eps/b$；最优为 $s_t$ & 小 $b$ 时标签幅度大，常需权重或尺度化\\
$v$ & $a\eps-bX$，限 $a^2+b^2=1$ & 旋转坐标稳定转换；与 ODE 速度需另作区别\\
流速度 & $\dot\gamma_0 Z+\dot\gamma_1 X$，系数按流路径定义 & 时间重参数化直接改变速度和权重\\
\bottomrule
\end{tabularx}
\end{table}

表中的 $\gamma_0,\gamma_1$ 是流路径的噪声系数和数据系数，将在第~\ref{sec:m-flow}~节定义；它们不自动等于连续扩散的 $a,b$。

例如，兼容估计器满足
\[
\norm{\widehat\eps-\eps}^2
=\frac{a^2}{b^2}\norm{\widehat X-X}^2.
\]
所以“均匀时刻下的无权重噪声损失”在数据预测坐标中，是按信噪比
$a^2/b^2$ 加权的损失。换目标却保持表面上同一个均方误差，实际上改变了训练关注的噪声区域。

在 VP 路径中，$v$ 的几何尤其清楚。写
$a=\cos\phi,\ b=\sin\phi$，就有
\begin{equation}
\begin{pmatrix}x_t\\v\end{pmatrix}
=
\begin{pmatrix}a&b\\-b&a\end{pmatrix}
\begin{pmatrix}X\\\eps\end{pmatrix},
\qquad
X=ax_t-bv,\quad\eps=bx_t+av.
\label{eq:m-v-rotation}
\end{equation}
这是正交旋转，逆变换没有 $1/a$ 或 $1/b$，因而在两端更稳定。
$v$ 预测在渐进蒸馏等工作中发挥了作用。\cite{r_vpred}
若固定一对 $(X,\eps)$，其插值轨迹对角度的导数正是
$\dd x/\dd\phi=v$。但对扩散时间的导数为
$\dd x/\dd t=\phi'(t)v$；网络给出条件平均时，对应的是该路径的边缘速度。只有连同路径、时钟和条件平均一起说明，才可以把 $v$ 与速度场联系起来。不同系数的 Flow Matching 速度不能直接套用式~\eqref{eq:m-v-rotation}。

对同一 $x_t$，兼容的 $v$ 与数据误差还有
$\delta X=-b\,\delta v$，所以无权重 $v$ 损失对应数据误差权重
$1/b^2$。这使它与噪声预测的信噪比权重不同。参数化是改善数值和优化的工具，不是凭命名保证更好生成的魔法。

\section{ODE 加速：把已知线性部分解析地走完}
\label{sec:m-solvers}

\subsection{为何更多阶数未必意味着更快}

使用第~\ref{sec:pfode}~节的概率流 ODE 后，采样变成数值积分问题。网络每次给出向量场的估计，求解器将其转成一段有限运动。Euler 用起点斜率，Heun 用预测终点的斜率修正，其他 Runge--Kutta 方法再利用中间点。高阶方法在足够光滑、步长足够小时可以更快降低误差，但每步往往需要更多网络调用。

因此评价时应区分采样步数与\emph{网络函数评估次数}（NFE）。一个 20 步 Heun 求解器可能接近调用网络 40 次；一个多步法可复用历史输出，每个新步只增加一次调用。文本条件生成中的\term{无分类器引导}{gl:cfg}还可能需要有条件和无条件两路网络计算：即使合并到一个批次，实际计算量也不凭批次数自动减半。延迟、吞吐量、显存和 NFE 各回答不同问题。

同样重要的是模型误差。记真实场为 $w$，网络场为 $w_\theta$，它们之差是学习误差；求解器对 $w_\theta$ 的误差是离散误差。无限增加步数最多精确求解网络给出的方程，不能自动修复网络场。反过来，一个更好的网络若只用极粗网格，也可能因数值误差失去优势。

\subsection{DPM-Solver 的基本积分恒等式}

DPM-Solver 利用扩散 ODE 的半线性结构：先解析处理已知缩放项，再对网络相关的积分做高阶近似。\cite{m_dpmsolver} 这与把整条方程当作黑箱的通用求解器是不同的利用信息方式。

对高斯路径 $x_t=a(t)X+b(t)\eps$，设前向漂移
$f(t)x$，其中 $f=a'/a$，扩散方差率
$g^2=(b^2)'-2fb^2$。噪声预测形式的概率流 ODE 为
\begin{equation}
\frac{\dd x}{\dd t}
=f(t)x+\frac{g(t)^2}{2b(t)}\eps_\theta(x,t)
=\frac{a'}{a}x+\left(b'-\frac{a'}{a}b\right)\eps_\theta(x,t).
\label{eq:m-semilinear}
\end{equation}
除以 $a$，已知线性项恰好抵消：
\[
 \frac{\dd(x/a)}{\dd t}
 =\frac{\dd(b/a)}{\dd t}\eps_\theta(x,t).
\]
这已经给出式~\eqref{eq:m-ddim-rescaled-ode}。再令
\[
 \lambda(t)=\log\frac{a(t)}{b(t)}
 =\tfrac12\log\operatorname{SNR}(t),
\]
注意本节 $\lambda$ 是半对数信噪比，文献有时将完整对数信噪比也记为同一字母。因为 $b/a=e^{-\lambda}$，从较高噪声的时刻 $t$ 生成到较低噪声的 $s<t$，有
\begin{equation}
x_s=\frac{a(s)}{a(t)}x_t
-a(s)\int_{\lambda(t)}^{\lambda(s)}
e^{-\lambda}\eps_\theta(x(\lambda),t(\lambda))\dd\lambda.
\label{eq:m-dpm-exact}
\end{equation}
这是学习 ODE 的精确解恒等式，不是假定网络不变。难点只剩被积函数沿轨迹如何变化。

令 $h=\lambda(s)-\lambda(t)>0$，如果把积分中的网络输出冻结为当前值 $e_t$，就得到
\begin{equation}
x_s\approx\frac{a(s)}{a(t)}x_t
-b(s)(e^h-1)e_t.
\label{eq:m-dpm-first}
\end{equation}
这与相应确定性 DDIM 更新代数等价。它给出了 DDIM 的一阶积分视角，也说明“解析处理线性项”不等于“整个 ODE 已经解析解出”。更高阶求解器通过额外时刻的预测或历史输出估计网络变化，再把指数权重下的多项式项积分掉。

\begin{derivationbox}{数据预测形式的积分也可以解析整理}
令 $D_\theta=(x-b\eps_\theta)/a$，式~\eqref{eq:m-semilinear} 可重写为
\[
 \frac{\dd x}{\dd t}
 =\frac{b'}b x+a\lambda'(t)D_\theta(x,t).
\]
于是
\[
 \frac{\dd(x/b)}{\dd\lambda}
 =e^\lambda D_\theta(x(\lambda),t(\lambda)),
\]
从而
\[
 x_s=\frac{b(s)}{b(t)}x_t+
 b(s)\int_{\lambda(t)}^{\lambda(s)}
 e^\lambda D_\theta(x(\lambda),t(\lambda))\dd\lambda.
\]
冻结数据预测得到
$x_s\approx\frac{b(s)}{b(t)}x_t+
a(s)(1-e^{-h})D_\theta(x_t,t)$。
一阶时它与式~\eqref{eq:m-dpm-first} 等价；高阶时，对噪声预测和数据预测分别近似，可能具有不同的稳定性。
\end{derivationbox}

DPM-Solver++ 用数据预测形式组织高阶积分，并针对引导采样提出稳定性改进与多步变体。\cite{m_dpmsolverpp} 其价值可从上面的两种积分理解：两者在连续方程上等价，却把非线性变化交给不同的被积函数近似。大引导会放大场的变化，更需要注意数值稳定性。

选择求解器时，至少要检查模型输出如何转换、时间与噪声日程是否一致、起终点是否落在可用区间。理论表达中的 $a=0$、$b=0$ 会让 $\lambda$ 发散，实际通常在有限端点积分，再做终端去噪。若潜变量没有天然的 $[-1,1]$ 边界，随意套用像素空间的数据裁剪还会改变所解的场。低 NFE 场景应实测质量与稳定性，不能只按求解阶数选名称。

\section{Flow Matching：直接学会把概率质量运过去}
\label{sec:m-flow}

\subsection{不必先制造随机扩散，才能训练生成 ODE}

Flow Matching 的入口是连续运输。先选择一族容易抽样的条件概率路径，再回归其条件速度；条件回归得到边缘速度，生成时解 ODE。它兼容扩散型高斯路径，也允许其他路径。\cite{r_fm} 与第~\ref{sec:flowmatching}~节一样，本节用流时钟 $\tau$，从噪声端 0 走向数据端 1。

最直观的训练方案是抽取 $Z\sim\pnoise$、$X\sim\pdata$，选择一对端点的\term{耦合}{gl:coupling}，构造
\begin{equation}
Y_\tau=(1-\tau)Z+\tau X,\qquad
V=X-Z.
\label{eq:m-linear-flow}
\end{equation}
然后最小化
\begin{equation}
L_{\mathrm{lin}}
=\E_{\tau,Z,X}\norm{u_\theta(Y_\tau,\tau)-(X-Z)}^2.
\label{eq:m-flow-loss}
\end{equation}
若端点独立，这叫独立耦合；若用某种配对规则产生端点，必须确保各自边缘仍是噪声和数据分布。耦合改变的是“哪些起点与哪些终点组成训练对”，它影响中间分布和速度回归难度。

式~\eqref{eq:m-flow-loss} 没有要求训练阶段解 ODE，因此常称“无模拟训练”。这表示造训练样本和标签时不用模拟生成轨迹，不表示生成阶段不用积分，也不表示整个训练没有随机计算。

\begin{examplebox}{直线 Flow Matching 的训练与生成}
\textbf{训练：}抽 $(Z,X)$ 与 $\tau$，计算 $Y_\tau$，回归 $X-Z$，更新参数。

\textbf{生成：}取 $y_0\sim\pnoise$，选网格
$0=\tau_0<\cdots<\tau_N=1$，例如用 Euler：
\[
 y_{i+1}=y_i+(\tau_{i+1}-\tau_i)u_\theta(y_i,\tau_i).
\]
训练标签的 $X$ 只在训练阶段可用；生成阶段的运动完全由当前 $y_i$、时刻和条件决定。
\end{examplebox}

一般路径可以写为
$Y_\tau=\gamma_0(\tau)Z+\gamma_1(\tau)X$，满足
$\gamma_0(0)=1,\gamma_1(0)=0,\gamma_0(1)=0,\gamma_1(1)=1$，相应标签为
$\dot\gamma_0Z+\dot\gamma_1X$。这说明 Flow Matching 首先是一套回归速度的训练原则，不是“所有方法都必须使用直线”的名称。

\subsection{直线标签为什么可能学出弯曲的生成轨迹}

每个训练对的路径确实是一条直线，但网络观察的是当前位置，不是端点对。最优网络为
\begin{equation}
u^*_\tau(y)=\E[X-Z\mid Y_\tau=y].
\label{eq:m-flow-mean}
\end{equation}
在同一位置交汇的不同训练对可能有不同方向，网络给出条件平均方向。沿着这个边缘场走的 ODE 轨迹，通常不等于任何一个训练端点对的直线。它会随位置变化而弯曲，但在适当正则性与精确学习下仍运输所需边缘分布。

\begin{figure}[htbp]
\centering
\begin{tikzpicture}[x=0.95cm,y=0.85cm,font=\small,>=stealth]
\draw[->] (-0.4,0)--(5.1,0) node[right]{$\tau$};
\node[left] at (0,2.8){端点配对};
\draw[black!30,dashed] (0,0.5)--(4.5,2.5);
\draw[black!30,dashed] (0,2.5)--(4.5,0.5);
\draw[black!30,dashed] (0,1.1)--(4.5,2.0);
\draw[black!30,dashed] (0,2.0)--(4.5,1.1);
\fill (0,0.5) circle(2pt);\fill (0,2.5) circle(2pt);
\fill (4.5,2.5) circle(2pt);\fill (4.5,0.5) circle(2pt);
\node[below] at (0,0){噪声};
\node[below] at (4.5,0){数据};
\node at (2.2,3.1){每对标签路径是直线};
\begin{scope}[xshift=7.0cm]
\draw[->] (-0.4,0)--(5.1,0) node[right]{$\tau$};
\draw[thick,->] (0,0.5) .. controls (1.4,0.15) and (3.0,0.9) .. (4.5,0.5);
\draw[thick,->] (0,2.5) .. controls (1.1,2.9) and (3.3,2.1) .. (4.5,2.5);
\draw[thick,->] (0,1.1) .. controls (1.8,1.0) and (2.7,1.8) .. (4.5,2.0);
\node[below] at (0,0){噪声};
\node[below] at (4.5,0){数据};
\node at (2.2,3.1){回归边缘场后的轨迹可弯曲};
\end{scope}
\end{tikzpicture}
\caption{直线条件路径与边缘 ODE 轨迹的区别。此图为概念示意，曲线用于强调二者不应逐条等同，不代表特定分布的数值解。}
\label{fig:m-flow-lines}
\end{figure}

用平方损失的分解还能看到耦合的重要性：
\begin{align}
&\E\norm{u_\theta(Y_\tau,\tau)-V}^2\\
&\quad=\E\norm{u_\theta(Y_\tau,\tau)-u^*_\tau(Y_\tau)}^2
 +\E\norm{V-u^*_\tau(Y_\tau)}^2.
\end{align}
第二项是当前位置无法分辨端点对造成的条件方差，固定耦合后它与网络参数无关。改善配对可减少冲突方向，也可能让运输更短；这是一条改变学习任务几何的路线，不是简单提高网络求解阶数。

在噪声端与数据端，插值路径可能出现尖锐尺度变化；真实图像分布又往往接近低维结构。因此仍需关注端点截断、时间采样分布、潜空间缩放和网络正则性。抽样 $\tau$ 的密度偏向某区域而不补偿时，相当于给该区域增加训练权重。一个流模型可以使用很简单的 Euler，也可以使用高阶 ODE 求解器；直线标签不免除数值检验。

\subsection{与扩散 ODE 的共同点和不同点}

二者都可以生成一个随时间变化的确定性速度场，都可用神经网络和 ODE 求解器。扩散路径先给出带噪边缘，score 再转换为概率流速度；Flow Matching 直接回归条件速度。若选择相同高斯路径，正确转换并调整权重，二者可以对应同一边缘场。若选择不同路径或不同耦合，生成几何已经不同，不能称为只把“预测噪声”重命名成“预测速度”。

尤其要记住 $\tau=1-t$ 只是时钟反向。对于某一固定路径，它把速度乘以负号；如果同时把 VP 三角系数换成 $(1-\tau,\tau)$，就连路径与信噪比日程都变了。把差异拆成“路径变化、时钟变化、输出变化、损失变化”，往往比争论哪个方法属于哪个家族更有用。

\section{Rectified Flow：端点配对和 reflow 为什么有用}
\label{sec:m-rectified}

Rectified Flow 从连接两端样本的直线出发，回归速度并用所得 ODE 产生新的确定性端点耦合；重复这个过程称为 reflow。\cite{m_rectified} 它与直线条件 Flow Matching 的目标密切相连，强调的另一个核心是端点耦合怎样被逐次改变。

第一轮可使用独立端点对 $(Z,X)$。训练得到 $u^{(1)}$ 后，从 $Z$ 积分生成
$\widehat X=\Phi^{(1)}_{0\to1}(Z)$。这样产生新的训练对
$(Z,\widehat X)$，两端有确定性关系。第二轮对新配对重新做直线插值和速度回归，得到 $u^{(2)}$。第一轮中的目标 $X$ 来自真实数据，第二轮的目标 $\widehat X$ 则来自模型生成；这一变化必须清楚记录。

\begin{examplebox}{一轮 reflow 的操作顺序}
\begin{enumerate}[leftmargin=*,itemsep=2pt]
\item 冻结当前速度网络，抽取许多 $Z$，用足够精细的求解器生成对应端点 $\widehat X$。
\item 保存配对关系 $(Z,\widehat X)$，条件生成时同时保存条件 $c$。
\item 对这些新端点对随机抽 $\tau$，构造
$(1-\tau)Z+\tau\widehat X$，回归标签 $\widehat X-Z$。
\item 用新网络评估轨迹弯曲程度、低 NFE 质量和分布覆盖，再决定是否继续一轮。
\end{enumerate}
重新训练的是新耦合上的速度，不是直接把旧网络采样步数改成 1。
\end{examplebox}

为什么它可能缩短运输？考虑理想化情形：第一轮 ODE 的边缘与插值边缘一致，并满足使轨迹存在和唯一的条件。记其轨迹为 $Q_\tau$，则
\[
 Q_1-Q_0=\int_0^1u^*_\tau(Q_\tau)\dd\tau.
\]
由积分的 Jensen 不等式以及条件期望的 Jensen 不等式，平方成本满足
\begin{align}
\E\norm{Q_1-Q_0}^2
&\leq \int_0^1\E\norm{u^*_\tau(Q_\tau)}^2\dd\tau\\*
&=\int_0^1\E\norm{\E[X-Z\mid Y_\tau]}^2\dd\tau\\*
&\leq \E\norm{X-Z}^2.
\label{eq:m-reflow-cost}
\end{align}
中间等式用了 $Q_\tau$ 与 $Y_\tau$ 的边缘一致性。更一般的凸位移成本也可用同类推理。这是“重新配对不会增加某些理想运输成本”的含义，不能直接替换为“有限网络一定达到全局最优运输”。

更直的轨迹为什么帮助少步采样？若一条轨迹上的速度基本不变，
Euler 的冻结斜率就更准确；若速度沿途大幅改变，一步就会走偏。
但端点直线距离较短、轨迹几何较直和速度随时间变化较小，是相关而不完全相同的性质。非均匀时间重参数化可以让同一条直轨迹的速度大小变化剧烈，所以少步误差仍要检查时钟。

实际 reflow 还有边界：教师端点已经含有生成误差，再训练只能保留或改变教师所覆盖的分布；数值生成配对也耗费计算；若训练端点的多样性不足，少步学生可能继承这种不足。reflow 是改变耦合、让少步运动更容易的一条路线，蒸馏则更直接地学习有限步映射。二者可以结合，却不是同一个操作。

\section{从求场到求结果：蒸馏与一致性模型}
\label{sec:m-distillation}

\subsection{渐进蒸馏：让学生的一步匹配教师的两步}

\term{蒸馏}{gl:distillation}是用已训练模型的行为监督另一个模型。扩散蒸馏的特别之处在于，可以监督一段运动的终点，而非只要求学生在每个时刻预测同一个局部场。渐进蒸馏用确定性教师的两步结果作为学生一步目标，再多轮减半采样步数。\cite{r_vpred}

设教师一步映射为 $T_{k\to j}$，下一步为 $T_{j\to i}$，学生为
$S^\theta_{k\to i}$，目标可表示为
\begin{equation}
L_{\mathrm{dist}}
=\E\norm{S^\theta_{k\to i}(x_k)
-T_{j\to i}(T_{k\to j}(x_k))}^2.
\label{eq:m-distill}
\end{equation}
对于 DDIM 式的学生一步映射，可将教师终点反解成所需的数据或其他预测标签，让学生仍使用熟悉的输出接口。学生学习的是在这个跨度上能正确落点的输出，未必等于原模型的无穷小局部噪声预测。

这一区别很重要。数值求解器保持原网络，努力更准确地积分原场；蒸馏改变网络，让粗步规则本身在训练分布上更好用。一个只为四步网格训练的学生，如果被套上原教师的 1000 步祖先采样器，未必正确。学生的适用步数、时刻区间、条件输入和引导尺度都是模型接口的一部分。

教师可以生成很多便宜的监督对，但这些对反映教师质量。若学生容量足够、优化充分，能逼近教师的有限步生成器；若学生容量有限或只用极少步，误差仍会积累。减少推理计算通常会增加训练或配对生成计算，适合反复部署同一个模型，而不是无需成本的数学变换。

\subsection{一致性模型：同一条轨迹的各点映到同一端点}

一致性模型换了输出问题。不是问“这里的方向是什么”，而是问“沿这条概率流轨迹继续走，最后会到哪里”。它学习从不同噪声级别直接映到低噪声端点的函数，并支持一步或少步生成；原工作给出了蒸馏训练和独立训练两类路径。\cite{m_consistency}

将端点设为一个很小的正扩散时刻 $t_{\min}$，以避开可能奇异的精确数据端，理想端点映射为
\[
 C(x_t,t)=\Phi_{t\to t_{\min}}(x_t).
\]
同一条 ODE 轨迹上的两点满足
\begin{equation}
C(x_t,t)=C(x_s,s),\qquad
C(x,t_{\min})=x.
\label{eq:m-consistency}
\end{equation}
由链式法则，它还满足
\[
 \partial_tC(x,t)+J_xC(x,t)w(x,t)=0,
\]
其中 $w$ 是概率流场，$J_xC$ 是输入\term{雅可比矩阵}{gl:jacobian}。它并不是零速度方程，而是沿着轨迹端点标签不改变的运输方程。

蒸馏时，从较高噪声点 $x_{t_{i+1}}$ 出发，用教师 ODE 求解器走到
$\widetilde x_{t_i}$。然后让学生预测
$C_\theta(x_{t_{i+1}},t_{i+1})$ 接近一个停止梯度的目标
$C_{\theta^-}(\widetilde x_{t_i},t_i)$。目标网络可采用学生参数的指数滑动平均，用于减缓目标随优化的变化。若只要求两边相同而没有端点条件，恒定函数也能满足一致性，因此边界不是装饰。

常见结构
\[
 C_\theta(x,t)=c_{\mathrm{skip}}(t)x+
 c_{\mathrm{out}}(t)F_\theta(x,t)
\]
令 $c_{\mathrm{skip}}(t_{\min})=1$、
$c_{\mathrm{out}}(t_{\min})=0$，就把端点边界嵌入模型。独立一致性训练则从带噪数据构造局部配对，不依赖已经训练好的扩散教师；其配对、离散尺度和优化方法有专门设计，不能把任意两个共享噪声的时刻预测相同当作完整算法。

一次生成可从高噪声 $x_{t_{\max}}$ 直接计算 $C_\theta$。少步生成还可先得到低噪声预测，再按选定路径重新注入适量噪声、重复映射。这个迭代不是普通 Euler 对局部速度积分，而是在直接端点映射之间安排修正。精确端点映射能一步生成，是因为它把整个积分结果编码到了函数中；普通去噪回归的条件均值没有这个保证。

\subsection{LCM：把端点映射放进潜空间并吸收引导}

LCM 将一致性蒸馏用于潜空间扩散，并把引导后的生成过程作为需要快速近似的目标。\cite{m_lcm} 潜空间的作用是降低生成状态的维度和每次网络调用的成本，数学上仍需要定义噪声路径、场和端点映射。潜变量必须遵循编码器训练时的尺度；最终图像还要经过解码器。

在式~\eqref{eq:cfg} 的约定下，引导场以条件预测和无条件预测线性组合。若学生要吸收一个范围的引导尺度，可以把尺度也作为输入，并在该范围内造教师目标。这样推理时减少重复计算具有明确含义：学生学会了某类引导后的映射。训练范围之外的尺度、不同的负面条件或者修改后的噪声日程，可能不再属于同一个目标接口。

这里还有一个概率概念需要谨慎：使用引导后的近似场，不能笼统声称它仍精确生成原无条件数据分布。引导主动改变了运动方向；它常提高条件服从度，但也会改变多样性和数值稳定性。蒸馏得到的学生首先逼近教师的引导行为，是否满足其他概率性质需要另外分析。

\section{把谱系用于判断：面对一个新方法看什么}
\label{sec:m-reading}

\begin{table}[htbp]
\centering
\caption{本章各路线主要改变哪一环。年代按首个公开版本，会议年份可能在下一年。}
\label{tab:m-lineage}
\small
\begin{tabularx}{\linewidth}{p{0.16\linewidth}p{0.11\linewidth}XX}
\toprule
路线 & 公开年份 & 核心操作 & 理解时应保留的边界\\
\midrule
DDPM & 2020 & 固定前向链，回归噪声，学习反向高斯 & 反向有限步条件一般只近似高斯\\
DDIM & 2020 & 保留条件边缘，改变联合构造和更新 & 确定性有限步更新仍有误差\\
Score SDE & 2020 & 将 score 与反向随机过程、ODE 联系 & 同边缘需真实场和正确积分\\
Improved DDPM & 2021 & 日程、方差及训练目标改进 & 似然目标与图像质量目标有不同权重\\
EDM & 2022 & 预条件、训练噪声与数值采样整理 & 设计部件可以分开使用\\
DPM-Solver & 2022 & 利用半线性结构的高阶积分 & 求解器不能消除场的学习误差\\
Flow Matching & 2022 & 条件速度回归，学习生成流 & 路径与耦合需要明确指定\\
Rectified Flow & 2022 & 直线回归与重新端点配对 & 直线标签不保证边缘轨迹直\\
渐进蒸馏 & 2022 & 学生一步拟合教师多步结果 & 依赖教师与训练跨度接口\\
一致性、LCM & 2023 & 不同噪声点直接映到轨迹端点 & 端点约束与目标配对不可省略\\
\bottomrule
\end{tabularx}
\end{table}

表~\ref{tab:m-lineage} 不是“后者淘汰前者”的时间线。例如，潜空间 DiT 可以使用 Flow Matching 训练，也可以蒸馏为少步学生；噪声预测 U-Net 可以换成 DPM-Solver 采样；EDM 的尺度思想可以帮助其他去噪网络。一个系统通常同时处于几条设计线上。

读新论文时，可先把它还原成以下操作：输入 $y$ 和时间，网络输出 $q_\theta$；目标标签是什么；损失在哪些时刻取期望；输出如何转换成生成场或映射；采样器计算什么。然后再问它改进的证据属于哪个维度：固定网络比求解器，固定计算量比质量，固定教师比学生，还是扩大训练数据与参数之后比最终系统。

例如，“只需四步”至少有四种可能解释：用普通网络粗积分，设计更好的有限步求解器，训练更容易积分的路径，或者蒸馏了四步映射。四种方式在推理接口上都可能出现四次循环，但数学意义、训练成本以及适用边界不同。若实验只有总步数，没有 NFE、教师成本、引导计算和质量指标，就还不能完整评估效率。

生成质量也不是一个数字。FID 使用特征统计衡量两组样本的差异，参见术语附录中的\term{FID}{gl:fid}；它不能逐项验证文字、空间关系、长尾概念和模式覆盖。少步方法可能在某个数据集指标上表现很好，却在复杂组合条件中误差明显。教材中最值得形成的习惯，是先写清比较问题，再选择证据。

\subsection{一个可复现的二维思想实验}

在真实图像之前，可以用二维双峰数据检验自己的理解。设数据为
\[
 X\sim\tfrac12\N((-2,0),0.2^2I)
       +\tfrac12\N((2,0),0.2^2I),
\qquad Z\sim\N(0,I).
\]
对于高斯扰动路径，每个分量的含噪均值和协方差都能解析计算：第 $i$ 个分量均值为 $a(t)m_i$，协方差为
$[a(t)^2\,0.2^2+b(t)^2]I$。将两个高斯密度按权重相加，求其对数梯度，就得到精确 score。

这时无需训练网络也能比较采样器，先隔离数值误差。用精确场分别运行 ODE Euler、Heun 和合适坐标中的 DDIM，固定 NFE 逐次细化网格，观察两峰是否都出现、峰宽和权重是否正确。再把场换成拟合的小网络，就能看到学习误差叠加后的效果。二维实验不能预测高维图像的最终质量，但能检查公式、方向和尺度是否写错。

对直线 Flow Matching，也可抽大量端点对在网格附近估计条件速度均值，画出场与轨迹。对比原始独立配对和教师生成的新配对，观察“训练直线”和“生成曲线”是否不同。重要的不是把图画得像论文，而是把理论中不可见的条件平均变成可观察对象。

\begin{keybox}{三种误差，三种不同的补救}
\textbf{路径或端点不匹配：}检查训练与采样的噪声定义、终点先验、潜变量尺度。

\textbf{场的学习误差：}改善网络、数据、目标权重与训练覆盖。

\textbf{积分或有限步映射误差：}调整节点、求解器、路径几何，或用蒸馏学习粗步结果。

增加采样步数主要针对第三项中的积分误差；无法替代前两项的检查。
\end{keybox}

\section{练习与解答思路}
\label{sec:m-exercises}

\begin{exercisebox}{练习一：一次预测为什么不能恢复分布}
设 $X$ 等概率取 $-1$ 和 $1$，$x_t=aX+b\eps$，$\eps\sim\N(0,1)$。
求 $\E[X\mid x_t=x]$，并说明从很高噪声一次输出该条件均值，为什么不会产生等概率的两个端点。它与一致性模型一步生成有什么不同？
\end{exercisebox}

\textbf{解答思路。}两类后验概率之比为
$\exp(2ax/b^2)$，因此
$\E[X\mid x_t=x]=\tanh(ax/b^2)$。当 $a\to0$，它趋近 0；
条件平均把不确定端点平均掉。端点映射模型则学习确定性运输轨迹的落点，可以把不同噪声起点送到不同端点；其监督不是该数据去噪条件均值。精确离散端点会涉及终端奇异性，实践中可先以有限低噪声分布理解这种运输，再讨论极限。

\begin{exercisebox}{练习二：推导 DDIM 的一阶坐标}
从 $\eta=0$ 的式~\eqref{eq:m-ddim-step} 出发，令
$y_k=x_k/a_k$、$r_k=b_k/a_k$，推导
$y_j=y_k+(r_j-r_k)e_k$。它是否等于原始时间 $t$ 中的 Euler 更新？为什么一步跳到 $j=0$ 的反例不违背前向边缘一致性？
\end{exercisebox}

\textbf{解答思路。}把 $\widehat x_0=(x_k-b_ke_k)/a_k$ 代入并除以
$a_j$ 即得。Euler 的形式依赖坐标和时间变量，变换前后有限步规则一般不同。前向边缘证明在已知真实 $x_0$ 条件下成立，生成时替换成回归估计并做粗步更新，不具有相同保证。

\begin{exercisebox}{练习三：转换目标等于转换权重}
在 VP 路径中，使用彼此兼容的估计器，证明
$\norm{\delta\eps}^2=(a^2/b^2)\norm{\delta X}^2$，
以及 $\norm{\delta v}^2=\norm{\delta X}^2/b^2$。
若想让三种参数化都优化无权重数据均方误差，各自应乘什么权重？
\end{exercisebox}

\textbf{解答思路。}固定输入 $x_t$，分别对
$X=(x_t-b\eps)/a$ 和 $X=ax_t-bv$ 作差即可。噪声损失乘
$b^2/a^2$，$v$ 损失乘 $b^2$，数据损失乘 1。在精确端点某些权重或转换奇异，因此实际应指定有限区间，并认识到有限网络仍可能因优化尺度产生差异。

\begin{exercisebox}{练习四：时间变换会改变直线速度}
对 $Y_\tau=(1-\tau)Z+\tau X$，改用时钟
$\tau=r(s)=s^2$，其中 $s\in[0,1]$。求关于 $s$ 的速度。
轨迹在空间中仍是直线，为什么用均匀 $s$ 的一步 Euler 可能失败？
\end{exercisebox}

\textbf{解答思路。}速度为 $2s(X-Z)$，起点速度为 0，因此一步
Euler 输出仍是 $Z$，尽管精确终点是 $X$。空间轨迹的直与时间速度的恒定是不同性质；选时钟、取节点和分析少步精度都需要考虑后者。

\begin{exercisebox}{练习五：一步学生和一个高阶求解器如何公平比较}
系统 A 用四步蒸馏网络，每步一次评估；系统 B 用两步 Heun，
每步两次评估原网络。提出至少四项应记录的量，并解释只报告“步数”可能得出什么错误结论。
\end{exercisebox}

\textbf{解答思路。}两者都约四次网络评估，但网络大小和引导实现可能不同。
至少记录 NFE、网络实际耗时、峰值显存、条件与引导设置、同一组样本的质量和多样性；若评价总成本，还记录学生训练和教师配对计算。固定质量比延迟与固定预算比质量都是合理问题，应分别报告。学生针对特定粗步接口训练，Heun 则近似原网络的局部 ODE，二者的可迁移性也需要独立测试。

\chaptertakeaway{算法名称背后有四条可组合的设计线：概率路径、监督输出、采样规则与网络结构。DDPM 提供容易训练的去噪任务，DDIM 使用相同单时刻信息安排另一种运动，Score SDE 连接随机过程与概率流，EDM 和 DPM-Solver 整理数值尺度与积分结构，Flow Matching 和 Rectified Flow 改变速度学习及端点耦合，蒸馏与一致性模型则把昂贵轨迹压缩为有限步结果。真正的理解，是能够写出它的输入、标签、场和更新，并指出每个等价结论成立的条件。下一章将这些部件放回真实生成系统，观察表示空间和网络结构如何扩展应用。}
}{}
\IfFileExists{chapters/04-applications.tex}{\chapter{从图像生成到多模态世界}
\label{ch:applications}

\readingguide{前两章告诉我们怎样学习一个把噪声变成数据的场，第三章告诉我们怎样训练和积分这个场。本章要问：这个场由什么网络表示，又怎样接收人的意图？我们先拆开图像生成系统，再沿着 U-Net、隐空间、DiT、视频和统一多模态模型的路线前进。阅读时始终区分概率路径、网络结构、条件接口与数据，它们共同决定系统能力，却不是同一件事。文献与产品信息核查截至 2026 年 10 月 9 日；2026 年的研究例子明确标注为预印本，产品公开演示不作为未披露架构的证据。}

\section{先拆开一套生成系统}
\label{sec:application-system}

看到一张清晰的生成图，最容易产生的印象是：模型内部一定藏着一个直接画图的巨大函数。实际系统通常更像一条分工明确的生产线。文字先被编码成条件；随机变量提供尚未确定的细节；一个神经网络反复预测噪声、干净样本、score 或速度；采样器把预测变成状态更新；最后，解码器可能把隐变量变成像素。用户看到的是最终图像，训练者需要同时照顾这条链上每一个环节。

可以把一套系统写成一个配置组
\[
\mathcal S=(\text{表示空间},\text{概率路径},\text{预测目标},
\text{网络},\text{条件},\text{采样器},\text{数据}).
\]
这个写法有一个实际好处：它阻止我们把性能变化全部归因于论文标题里最醒目的那个名词。把 U-Net 换成 Transformer，改变的是函数的参数化；把噪声预测换成直线路径速度回归，改变的是训练目标和概率路径；把像素换成 VAE 隐变量，改变的是网络处理的随机变量；把短标题换成详细描述，改变的是监督信号。它们能够组合，但没有哪一项独自保证生成质量。

图像生成网络的输入和输出往往形状相同。例如，输入是一个含噪的 $h\times w\times C$ 张量，输出也是同样大小的噪声或速度张量。这个外观容易让人以为生成就是图像到图像的回归。关键区别在于，输入的噪声水平也作为条件，而且同一个网络要覆盖整条概率路径。它学到的是许多分布之间的局部迁移规律，采样器才把这些局部规律连接成一张完整图。可以回看第\ref{ch:math}章的式\eqref{eq:reverse-sde}、\eqref{eq:pfode}和\eqref{eq:cfm}，把其中抽象的场理解成下文讨论的网络输出。

\begin{intuitionbox}{网络负责方向，采样器负责行走}
U-Net 和 DiT 回答的是“怎样用有限参数表示这个高维场”；DDPM、DDIM 和概率流 ODE 回答的是“沿着什么随机或确定性规则更新状态”。一张图是否清晰，既取决于方向是否准确，也取决于步长、条件和表示空间是否合适。
\end{intuitionbox}

这也解释了为什么软件中的一个模型名称，未必对应一种固定数学定义。一个产品系列可以在不同代际中更换主干网络、训练路径和文本编码器。阅读论文时，应当落实到具体版本、具体训练设置。我们下面使用历史工作作为概念的锚点，不把某个版本的实现推广为整个品牌永远不变的性质。

\section{U-Net 为什么适合学习图像中的场}
\label{sec:application-unet}

\subsection{从局部纹理到整体布局}

一张图像同时包含多个尺度的信息。头发的细纹只占几个像素，人脸的比例需要几十个像素，人物和背景的布局则跨越整张画面。如果每次去噪只看很小的邻域，模型可能把局部画得精致，却无法协调左右手、物体数量和远处背景。反过来，如果网络一开始就把画面压缩得很粗，细节的位置又可能丢失。

U-Net 以一条下降路径和一条上升路径回应这个矛盾。下降路径逐渐减小空间分辨率、增加通道数，把局部证据汇总成较大的语义结构；上升路径逐渐恢复空间分辨率，把整体决定落实到局部。对应尺度之间的跳跃连接直接传递浅层特征，因此深层网络不必凭空重建所有精细位置信息。这里的“U”描述拓扑形状，既不是某一种噪声日程，也不指定采样随机性。

图\ref{fig:application-unet}给出的是原创教学示意，不是某篇论文的逐层复刻。含噪输入经过不同尺度，最终得到一个同形状的预测。中间的跳跃连接传递的是特征，而不是自动保留了正确答案。尤其在高噪声阶段，浅层特征本身充满噪声，模型仍需根据时间条件决定如何使用它们。

\begin{figure}[htbp]
\centering
\begin{tikzpicture}[x=1cm,y=1cm, every node/.style={font=\small},
  block/.style={draw,rounded corners,align=center,minimum width=2.5cm,minimum height=0.8cm},
  arr/.style={->,thick}]
\node[block] (a) at (0,0) {高分辨率特征\\含噪输入};
\node[block] (b) at (0,-1.5) {中分辨率特征};
\node[block] (c) at (0,-3) {低分辨率特征};
\node[block] (d) at (6,-3) {低分辨率融合};
\node[block] (e) at (6,-1.5) {中分辨率融合};
\node[block] (f) at (6,0) {高分辨率预测\\噪声或速度};
\draw[arr] (a)--node[left]{下采样}(b);
\draw[arr] (b)--(c);
\draw[arr] (c)--node[above]{大范围交互}(d);
\draw[arr] (d)--(e);
\draw[arr] (e)--node[right]{上采样}(f);
\draw[arr,dashed] (a)--node[above]{跳跃连接}(f);
\draw[arr,dashed] (b)--node[above]{跳跃连接}(e);
\node[align=center] at (3,-4.1) {时间嵌入告诉每一层当前噪声水平；\\文字或类别条件告诉每一层要生成什么。};
\end{tikzpicture}
\caption{U-Net 的尺度分工。虚线传递对应尺度的特征；图中省略了残差块和具体注意力位置。}
\label{fig:application-unet}
\end{figure}

卷积提供了一个很有用的归纳偏置：同一种局部规则可以在画面的不同位置重复使用。一个边缘检测器不必为左上角和右下角分别学习。它特别适合纹理与局部几何，但纯卷积的长距离交互要经过多层传播。于是扩散 U-Net 通常还引入注意力，让相距较远的特征直接交换信息。注意力不是专门为语言设计的，它是一种按内容建立连接的方法，术语解释见\hyperref[gl:attention]{注意力条目}。

\subsection{时间如何进入网络}

把一幅几乎干净的图和一幅接近纯噪声的图交给同一个网络，却不说明噪声水平，会使回归目标变得含混。网络看到一个弱边缘时，不知道它是应该保留的图像结构，还是应当消除的噪声偶然性。因此，时间 $t$ 或离散步数 $k$ 通常先变成一个向量，再被送到许多残差块。

一个常见的构造是先用正弦、余弦形成不同频率的编码，再经过小型全连接网络。这样的编码把一个标量变成多个平滑变化的坐标，有助于网络区分相近的噪声水平。时间信息可以直接加到特征上，也可以产生特征归一化后的缩放和偏移。若 $h$ 是某层特征，后一种方式可以示意为
\[
\widetilde h=\alpha(t,c)\odot\operatorname{Norm}(h)+\beta(t,c).
\]
这里 $\odot$ 是逐元素乘法，$\alpha$、$\beta$ 由条件网络产生。不要把本式的 $\alpha(t,c)$ 与 DDPM 的 $\alpha_k=1-\beta_k$ 混为一谈：前者是神经网络的特征调制参数，后者是概率路径的系数。它们只是在不同文献中恰好使用了同一个希腊字母。

可以把特征调制理解为网络内部的一组工作模式。同样一层卷积，在大噪声时侧重粗略结构，在小噪声时侧重边缘和材质。这个比喻并不意味着每个通道都能得到清晰的人类语义解释；它说明的是条件可以在许多层持续改变计算，而不只是作为开头的一个标签出现。

\subsection{类别条件走向文字条件}

类别条件只需说明“这张图属于哪一类”。文字条件却包含属性、关系、数量和动作。一条提示词“左侧是一只红杯子，右侧是一只蓝杯子”，如果只压缩成一个整体向量，红、蓝与左右的绑定就可能不稳定。保留文字的 token 序列，允许不同图像位置查询不同文字片段，是更细致的接口。

令图像特征产生查询 $Q=H W_Q$，文字特征产生键和值 $K=CW_K$、$V=CW_V$。交叉注意力写为
\begin{equation}
\operatorname{CrossAttn}(H,C)
=\operatorname{softmax}\!\left(\frac{QK^\top}{\sqrt{d_k}}\right)V.
\label{eq:application-crossattention}
\end{equation}
矩阵的一行对应某个图像位置或图像块。该行通过文字键的相似度，对文字值进行加权汇总。$d_k$ 是单个注意力头中键的维数，除以其平方根是为了避免内积幅度随维数增长过快。softmax 是按行计算的归一化指数函数，因此每行权重和为一。

这个机制并不保证“红色”一定绑定到正确的杯子。训练只通过最终回归误差监督注意力，数据描述也可能只说“桌上的杯子”，没有提供所有关系。注意力提供了表示关系的通道，关系是否被学会还取决于数据、编码器和优化。这个区别非常重要：一种结构能表达某个能力，与训练后可靠拥有这个能力，是两种不同结论。

\section{文字到图像：条件比人们想象的更重要}
\label{sec:application-textimage}

\subsection{ADM：先建立一个高质量图像先验}

Dhariwal 与 Nichol 的 ADM 工作通过 U-Net 结构实验和分类器引导，说明扩散模型在当时的图像基准上可以与强 GAN 系统竞争。论文的关键教学价值，是把高质量图像先验与引导机制分开研究，而不是把优势归结为“多做几步去噪”。\cite{a_adm}

用第\ref{sec:guidance}节的语言说，类别引导在无条件 score 上增加 $\grad_x\log p(c\mid x_t)$。梯度告诉我们：在当前噪声水平附近，朝哪个方向移动更容易被识别成所需类别。引导系数增大时，样本可能更典型，但也更集中。图像“更像一只狗”不等于狗的全部姿态、品种和背景都被更好地覆盖。第\ref{ch:concepts}章中学习分布的目标，到这里变成了一个很现实的质量与覆盖率权衡。

这里也可看出分类器的额外要求。它不能只识别干净图像，还必须在各种噪声水平下给出有意义的梯度。如果分类器在高噪声处不稳定，错误方向会进入整个采样过程。这推动了不另训噪声分类器的条件方法发展。

\subsection{GLIDE：让文本参与每次去噪}

GLIDE 把文字条件引入扩散图像生成，比较了 CLIP 引导与\term{无分类器引导}{gl:cfg}，并展示文字引导的图像补全。其研究结果支持了把有条件与无条件预测组合起来的实用路线；公开小模型与论文大模型的规模和训练条件应分别阅读。\cite{a_glide}

无分类器引导之所以容易移植，源于它只要求网络在训练时有时看见条件、有时看不见条件。推理时，两次预测的差给出一个与条件相关的方向。这个做法可以对噪声、score 或恰当参数化的速度进行线性组合，但必须保持训练路径和参数换算一致，具体公式见式\eqref{eq:cfg}。它把“更听话”变成可调节的采样操作，并不自动补足训练数据中缺失的概念。

如果模型从未学会某种空间关系，把引导强度继续增大，可能只是强化一种错误解释。实践中的饱和、过强对比度和形态扭曲，提醒我们：引导是改变已有场，而不是向网络注入全新的知识。

\subsection{Imagen：语言理解也是图像生成的瓶颈}

Imagen 使用预训练大型文本编码器和级联扩散模型，并通过实验强调文本表征对画面质量与图文一致性的作用。论文还引入包含复杂提示的 DrawBench，帮助评估数量、属性和组合关系。\cite{a_imagen}

一个冻结的文本编码器为什么仍然重要？因为去噪网络并不直接看到自然语言字符串，它只能看到编码器输出的向量。若两条本应不同的提示在表征中几乎重合，后面的生成器再强，也难以稳定区分。反过来，语言编码器保留了语义差异，图像网络才有机会把差异落实到空间布局。

级联生成体现了另一种分工。基础模型先生成低分辨率图像；超分辨率模型在其基础上生成更细的图像。数学上，这可以示意为
\[
p(x^{\mathrm{high}}\mid c)
=\int p(x^{\mathrm{high}}\mid x^{\mathrm{low}},c)
        p(x^{\mathrm{low}}\mid c)\,\dd x^{\mathrm{low}}.
\]
低分辨率结果已经承担了构图的很大部分，因此高分辨率模型主要解决在给定结构下补足细节的问题。它能减轻一次生成整张大图的困难，却也可能继承上一级的计数或布局错误。把三只杯子放大，通常不会自动变成正确的四只杯子。

\subsection{DALL-E 2：语义隐变量也是一种生成对象}

DALL-E 2 的论文把图像生成分成文字到 CLIP 图像表征的先验，以及从该表征到图像的解码器。先验研究了自回归和扩散方案，解码阶段使用扩散生成。它展示了一种与空间 VAE 隐变量不同的语义中间层。\cite{a_dalle2}

CLIP 指对比式语言图像预训练：文字与图像分别经过编码器，配对样本的表征被拉近，不配对样本的表征被区分。它因此提供一个适合比较图文语义的空间。\cite{a_clip}CLIP 图像向量不是一张缩小后的图片，也不是像素的无损编码。

为了理解这一区别，考虑两个问题：“画面总体表达什么？”与“每一个像素在哪里？”语义表征更适合前一个问题，保留整体概念和图文对齐；空间隐变量更接近后一个问题，保留可解码的局部结构。两者都叫 latent，却服务于不同接口。把“用了隐空间”直接理解成“用了 Stable Diffusion 那样的压缩图像”，会错过模型的真正设计。

中间层越抽象，语义控制可能越便利，但对精确文字、细小位置和对象数量的保留也可能越困难。这个取舍会反复出现：生成系统需要既懂意义又记得细节，而一种紧凑表征很难免费满足二者。

\section{隐空间扩散：把算力花在更合适的变量上}
\label{sec:application-latent}

\subsection{从像素分布到压缩表征分布}

高分辨率像素数量很大，其中相当一部分变化来自局部纹理、采样噪声和细微颜色。一个生成器若要在每个采样步都处理所有像素，计算量会很快增加。隐空间扩散模型（Latent Diffusion Model，LDM）的重要选择，是先学习一个图像自编码器，再在它的空间隐变量上学习扩散先验，并用交叉注意力接入条件。\cite{c_rombach2022ldm}

记编码器为 $E_\phi$，解码器为 $D_\psi$，压缩后的干净隐变量为 $\ell_0$。实际编码器可以输出后验分布的参数，而不是唯一的确定向量；为了突出系统结构，暂时用 $\ell_0=E_\phi(x_0)$ 记一次编码结果。隐空间训练对 $\ell_0$ 加噪：
\[
\ell_t=a(t)\ell_0+b(t)\eps.
\]
第二章的去噪推导仍然适用，只是随机变量已经从像素换成隐变量。生成时先从噪声得到 $\widehat\ell_0$，再得到图像 $D_\psi(\widehat\ell_0)$。在随机解码器的写法下，总分布为
\begin{equation}
p_{\theta,\psi}(x\mid c)
=\int p_\psi(x\mid\ell)\,p_\theta(\ell\mid c)\,\dd\ell.
\label{eq:application-latent-model}
\end{equation}
若解码器是确定性的，就把 $p_\psi(x\mid\ell)$ 理解为集中在 $D_\psi(\ell)$ 处的点质量。该积分描述的是解码器对隐变量分布的映射，术语可对照\hyperref[gl:latent]{隐空间条目}和\hyperref[gl:vae]{VAE 条目}。

注意，VAE 中的先验正则不意味着真实编码结果严格服从标准高斯。图像经编码后形成的聚合分布仍然复杂；扩散模型正是来拟合这个分布。否则，若标准高斯直接经解码器已经能生成足够好的图像，就没有再训练昂贵隐空间先验的必要。

\subsection{为什么压缩并非越强越好}

设原图是 $H\times W\times 3$，编码器把空间边长缩小 $f$ 倍，隐通道数为 $C_\ell$。元素数之比为
\[
\frac{(H/f)(W/f)C_\ell}{3HW}
=\frac{C_\ell}{3f^2}.
\]
这个式子只是存储元素数的比较，并不直接等于训练速度提升：通道数、网络宽度、注意力和解码成本都参与实际计算。但它清楚显示了空间压缩为何有吸引力。例如一个教学设定 $f=8$、$C_\ell=4$，隐变量元素数是像素元素数的 $1/48$。这里是计算示例，不是对所有模型配置的断言。

代价是压缩会丢信息。若一串小字号文字在编码和重建之后已经模糊，隐空间生成器再精确，也不能从这个接口稳定恢复每个字形。它可能学到一个看似合理的纹理补偿，却没有受监督保留原来的文字。这是表示瓶颈，而不是采样器做得不够好。

\begin{examplebox}{先判断错误来自哪里}
假设一个模型总把商品包装上的小字画错。先把真实包装图送进编码器再解码。如果重建图已经错了，优先检查压缩率和自编码器；如果重建正确而生成错误，再检查文字条件、数据和生成先验。把所有错误都交给“更多去噪步数”处理，常常找错了层级。
\end{examplebox}

自编码器的目标也未必只是像素均方误差。它可以结合重建、\term{感知损失}{gl:perceptual}和适当的正则，使重建图保留人们在视觉上在意的结构。感知损失通常在某个特征空间比较图像，而不是逐像素比较。它有助于保持视觉锐度，但优化的“相似”取决于该特征空间，因此仍需按具体任务验证。

\subsection{Stable Diffusion 是一套可组合的接口}

以公开的 Stable Diffusion v1 系统为例，它把文本编码器、隐空间去噪 U-Net、自编码器和采样器连接起来。\cite{a_sdcard}理解这个组合比记住某一代模型的层数更有长期价值。文字编码器告诉模型意图，去噪主干学习条件隐变量分布，自编码器负责像素接口，采样器负责数值更新。只要接口保持相容，其中某个组件就可以改进而不必重训所有组件。

这解释了生态里为什么会出现许多不同采样器、条件插件与个性化微调。但“可连接”不等于任意组合都有效。更换文本编码器会改变特征维数和语义分布；更换 VAE 会改变隐变量尺度和几何；更换预测目标会改变采样器读取输出的方式。兼容性既是形状问题，也是训练分布问题。

SDXL 在隐空间路线中进一步研究更大的去噪网络、文本条件和分辨率相关条件，并包含基础模型与细化模型的安排。它特别适合作为一个提醒：高分辨率表现来自网络、数据、训练任务和条件的联合工程，而不是单独增加画布尺寸。\cite{a_sdxl}

我们可以亲自理解分辨率条件为什么有用。同一张照片，裁成正方形和保留原宽高比，画面中对象的位置会不同。如果训练时只告诉网络文字，不告诉它原尺寸和裁剪位置，这些不同结果就被混在同一个条件分布里。尺寸、宽高比和裁剪信息作为附加条件，可以让模型解释一部分原先看似随机的布局差异。

\section{DiT：把场学习变成 token 之间的协作}
\label{sec:application-dit}

\subsection{图像块为什么能成为 token}

Transformer 不要求输入必须是文字。它需要的是一个向量序列，以及识别这些向量位置或关系的办法。把隐变量切成互不重叠的小块，将每块展开并线性投影，就得到一串\term{token}{gl:token}。这里 token 是计算中的一个序列元素，它可以对应词片段，也可以对应\term{图像块}{gl:patch}，不能一律理解为文字词语。

若隐变量为 $h\times w\times C_\ell$，块边长为 $p$，并且 $h,w$ 可被 $p$ 整除，则图像 token 数量为
\begin{equation}
N=\frac{hw}{p^2}.
\label{eq:application-token-count}
\end{equation}
每个块经投影变成宽度 $D$ 的向量。位置编码告诉网络它来自哪里，否则一个向量序列在置换之后，网络难以知道“左边”和“右边”。Transformer 块让这些 token 通过自注意力交换信息，然后通过逐 token 的前馈网络做特征变换。最后，预测头把向量还原成每个块中的噪声或速度值，并重新拼成张量。

原始 DiT（Diffusion Transformer，扩散 Transformer）论文在隐空间中用 Transformer 替代扩散 U-Net，系统比较条件注入方法和模型计算规模。其典型设计使用自适应层归一化及零初始化残差调制；主要实验是类别条件生成，不应直接说原始 DiT 就是一套完整文字到图像系统。\cite{a_dit}

\subsection{条件调制与稳定优化}

Transformer 的残差更新可以示意为
\begin{align}
\widetilde H&=(1+\alpha_c)\odot\operatorname{LN}(H)+\beta_c,\\
H'&=H+\gamma_c\odot\operatorname{Attn}(\widetilde H).
\label{eq:application-adaln}
\end{align}
这里 $\operatorname{LN}$ 是层归一化，条件网络从时间和条件产生 $\alpha_c,\beta_c,\gamma_c$，公式省略了多头拆分和后续前馈子层。缩放、偏移决定注意力看到的特征，门 $\gamma_c$ 决定残差更新的强度。零初始化某些调制输出，可让训练开始时的残差分支接近零，深层网络先接近容易优化的恒等映射，再逐步学习复杂更新。

这个设计连接了数学与优化：理论上我们需要一个准确的时间相关向量场，实际上还需要训练从随机参数出发能稳定到达它。损失函数正确不代表优化会顺利，网络中的归一化和初始化属于这条桥梁的一部分。

\subsection{为什么“大一点”可能更好，也可能更贵}

一层标准全注意力中，注意力矩阵有 $N^2$ 个位置对，主要相关运算规模约为 $N^2D$；线性投影和前馈网络的主要规模约为 $ND^2$。忽略常数后，一次网络调用的计算可以写成
\[
\operatorname{Cost}\approx L\bigl(c_1N^2D+c_2ND^2\bigr),
\]
其中 $L$ 是层数，$c_1,c_2$ 是由具体结构决定的常数。实际高效实现可以减少存储开销、利用硬件并行，但不会凭空消除所有 token 交互成本。这个估算用来理解尺度，不用于预测某张显卡的精确运行时间。

减小块大小 $p$ 会增加 $N$。若从 $p=4$ 改成 $p=2$，token 数变为四倍，$N^2$ 部分变为十六倍。与此同时，空间描述更精细，模型可以在更细的粒度上交换信息。增加深度、宽度和 token 数，并非三种完全等价的扩容：它们提供不同的计算能力，也改变不同的瓶颈。

\begin{keybox}{规模规律是一种受条件约束的经验规律}
同一训练体系中，更多计算可能带来更低验证误差和更好样本；这不能推出任意数据、任意网络、任意预算下，参数更多都更好。比较时至少应报告数据、训练计算、采样次数和评价设置。
\end{keybox}

U-Net 的多尺度结构和卷积归纳偏置使局部视觉规律容易学习；DiT 的统一序列接口使长距离交互和多模态扩展自然。两者的区别不该简化成“旧网络”和“新网络”。在预算有限、任务较窄或局部结构特别重要时，卷积式设计仍有合理性；在大规模预训练、复杂条件和统一接口中，Transformer 展现了另一组优势。

\subsection{MM-DiT：文字不只是一份静态查询资料}

SD3 对应的研究结合直线路径的流训练、噪声水平采样研究和多模态 Transformer。MM-DiT 为图像和文字保留各自的参数变换，在联合注意力中交换信息，使文字表征也能随图像状态更新。这与固定文字表征供图像交叉查询的接口不同。\cite{a_sd3}

用简化的矩阵写法，先分别计算
\[
(Q_I,K_I,V_I)=\mathcal P_I(H_I),\qquad
(Q_T,K_T,V_T)=\mathcal P_T(H_T),
\]
再在注意力操作中拼接键、查询和值。例如图像侧更新使用
\[
\operatorname{softmax}\!\left(
 \frac{Q_I[K_I;K_T]^\top}{\sqrt{d_k}}\right)[V_I;V_T].
\]
文字侧也可以读取图像侧信息。分号表示沿序列维度拼接。两种模态不用强迫所有变换共享一套参数，但能够在每层协调。这是“多模态”一词在这里的具体含义：两类 token 共同参与场的计算；它不自动意味着模型能回答任意视觉问题。

在直线路径的\hyperref[sec:flowmatching]{流匹配}设定下，按本书时钟写为
\[
\ell_\tau=(1-\tau)Z+\tau\ell_0,
\qquad u^{\mathrm{target}}=\ell_0-Z.
\]
网络回归这个条件速度的最优平均，再由 ODE 生成隐变量。这里的速度与第\ref{sec:parameterization}节的 $v$-prediction 不是同一个定义，不能只因为两者都在代码里叫 \texttt{v} 就相互替换。原论文可能使用从数据到噪声的时钟，读者还需检查符号方向。

图\ref{fig:application-interfaces}把三种接口放在一起。所有箭头表示信息通道，不代表所有组件都端到端一起训练。能认清这张图，阅读新的生成模型时就不必被名字牵着走：先找它的表示空间，随后找模态之间在哪里交流，再找它究竟预测什么。

\begin{figure}[htbp]
\centering
\begin{tikzpicture}[x=1cm,y=1cm,every node/.style={font=\small},
 b/.style={draw,rounded corners,align=center,minimum height=0.65cm},
 ar/.style={->,thick}]
\node[anchor=east] at (-0.1,0) {交叉注意力};
\node[b,minimum width=2.1cm] (ta) at (1.7,0) {固定文字特征};
\node[b,minimum width=2.5cm] (ia) at (5.6,0) {含噪图像主干};
\draw[ar] (ta)--(ia);
\node[anchor=east] at (-0.1,-1.4) {联合注意力};
\node[b,minimum width=2.1cm] (tb) at (1.7,-1.4) {文字状态};
\node[b,minimum width=2.5cm] (ib) at (5.6,-1.4) {图像状态};
\draw[<->,thick] (tb)--node[above]{每层交流}(ib);
\node[anchor=east] at (-0.1,-2.8) {统一任务};
\node[b,minimum width=2.1cm] (tc) at (1.7,-2.8) {文字与视觉上下文};
\node[b,minimum width=2.5cm] (ic) at (5.6,-2.8) {共享或分专家主干};
\draw[ar] (tc)--(ic);
\node[b,minimum width=2.5cm] (out) at (5.6,-4.2) {文字概率／视觉速度};
\draw[ar] (ic)--(out);
\end{tikzpicture}
\caption{条件接口逐步走向统一任务的三种概念结构。双向交流只是结构能力，仍需任务和数据让它形成有用协作。}
\label{fig:application-interfaces}
\end{figure}

\section{编辑与控制：在可能的画面中选择符合要求的一张}
\label{sec:application-control}

\subsection{为什么一句文字不足以规定一幅图}

“一个人在街上骑自行车”可以对应无数姿态、构图和视角。文字擅长说出语义，却很难高效指定每条轮廓。条件图像、深度图、边缘图和姿态骨架能够补充空间约束。数学上，任务从 $p(x\mid c_{\mathrm{text}})$ 变成
\[
p(x\mid c_{\mathrm{text}},c_{\mathrm{spatial}}).
\]
新增条件使允许的画面集合缩小，但通常仍然保留大量纹理、光照和背景自由度。控制生成的目的，就是在约束下保留合理变化，而不是把所有自由度消灭。

补全是一个直观例子。令掩码 $M$ 在已知区域取一，在需要生成的区域取零，已知像素为 $x_{\mathrm{known}}$。一个理想约束是
\[
M\odot x_0=M\odot x_{\mathrm{known}}.
\]
生成器要根据可见区域推断不可见区域的条件分布。采样过程若持续使用已知区域的相应噪声状态，能够加强一致性；具体方法也可以训练专门接收掩码和条件图的网络。掩码边界附近仍需协调纹理、光照和几何，简单粘贴并不能代替条件建模。

\subsection{ControlNet：给预训练先验增加空间输入}

ControlNet 冻结预训练生成主干，用可训练分支处理附加空间条件，并通过零初始化卷积将分支信号送回主干。它让较小的任务数据能够利用已有图像先验，支持边缘、深度和姿态等条件。\cite{a_controlnet}

零初始化在这里有一个清楚的作用。训练刚开始，附加分支对主干的注入接近零，因此系统行为接近原来的生成器；随着训练进行，条件分支逐渐学会对不同尺度提供修正。这是一种有用的工程策略：不要让一个尚未学会含义的新接口立即扰乱成熟的表征。

条件图和生成图并非像素对像素的硬复制。骨架固定大致关节位置，仍然允许不同衣服和人物；深度固定相对空间结构，仍然允许不同材质；边缘图本身有歧义，同一轮廓可能对应不同对象。控制效果要评价约束是否被满足，也要评价保留的自由度是否合理。

\subsection{DreamBooth 与指令编辑：记住主体和理解修改}

DreamBooth 通过少量主体照片微调文字到图像模型，把一个特别的标识与主体关联，并用类别先验保留思想缓解微调后泛化受损的问题。它研究的是“这个具体对象如何进入已有生成先验”。\cite{a_dreambooth}

主体个性化的难点是区分身份与偶然环境。一组照片若都在同一背景、同一光照下拍摄，模型可能把背景当成身份的一部分。换一个地点生成时，便可能连背景也一起带走。增加视角和环境变化，并适度限制微调幅度，才有机会学到更稳定的主体特征。这个问题本质上是条件下的可辨识性：训练数据是否足以告诉模型哪些变化属于身份，哪些只是拍摄情境。

InstructPix2Pix 则学习原图、编辑指令和目标图之间的关系，利用合成的编辑配对训练条件扩散模型，推理时可直接依据指令修改图像。它强调编辑行为需要专门的监督，而不是仅凭文字生成能力便会自然出现。\cite{a_instructpix}

对编辑任务，单纯评价最终图像漂亮与否是不够的。若要求“把杯子变成蓝色”，模型同时改变桌面、人物和构图，就没有正确完成局部指令。应当至少检查三件事：要求修改的属性是否改变，主体身份是否保留，不相关区域是否稳定。生成质量和编辑忠实度并不总是同向变化。

\begin{cautionbox}{编辑强度与保持程度的矛盾}
从原图加更多噪声再生成，通常会给模型更大改动自由，但也削弱原图约束。改得少不一定更正确，改得多也不一定更听话。理想编辑需要理解“应该改变什么”和“必须保持什么”，而不只是一个统一的强度旋钮。
\end{cautionbox}

\section{视频生成：时间一致性不是逐帧画图}
\label{sec:application-video}

\subsection{从边缘分布到联合分布}

设视频有 $F$ 帧，单独为每帧生成一张好图，只能保证每个边缘分布看起来合理。视频需要的是联合分布
\[
p(x^{(1)},\ldots,x^{(F)}\mid c).
\]
同一个人物跨帧应保持身份，物体运动应连续，遮挡后重现的东西应保留先前属性。即使每一帧单独看都毫无问题，它们的组合仍可能出现衣服闪烁、物体漂移和人数跳变。这与概率论里“边缘分布不能唯一确定联合分布”完全一致，相关概念见附录\ref{app:probability}。

早期 Video Diffusion Models 用时空网络研究视频的联合去噪，并探索图像和视频共同训练。它是理解从二维生成到时序生成的重要入口：时间轴进入网络，学习目标也随之涵盖跨帧相关性。\cite{a_videodiff}

可以把图像 U-Net 中的空间算子扩展成时空算子，也可以分别做空间注意力和时间注意力。后一种分解中，同一帧内的 token 先交换信息，相同空间位置附近的 token 再跨帧交流。它降低直接在全部时空 token 上交互的压力，但分解的顺序和范围也影响模型对复杂运动的表达。

\subsection{时空压缩和视频 Transformer}

视频张量比图像多了时间维度，内存与计算问题更加突出。若直接把 $F$ 帧都切成 token，数量约为 $N=Fhw/p^2$，全注意力中的 $N^2$ 项会随帧数平方增长。时空自编码器因此不仅压缩空间，也压缩时间，把多帧映射成较短的隐变量序列。

Wan 的 2025 年技术报告公开了视频 DiT、时空 VAE、数据处理与训练策略。其 Wan-VAE 使用因果时间结构；报告说明普通后续帧可按时间四倍、空间八倍压缩，第一帧单独保留时间处理方式。这些是具体公开实现，应与产品系列其他版本分开。\cite{a_wan}

“因果”在这里指一个时间位置的计算不使用未来帧信息，它支持分块编码、解码和缓存设计；它不意味着整个视频生成主干只能按帧自回归。自编码器的因果性、Transformer 的注意力范围与采样的时间顺序，属于三种不同层面的选择。

一个重要取舍是时间压缩率。压缩过强时，短促动作和嘴唇运动可能在接口中丢失；压缩较弱时，长视频计算昂贵。把时空压缩只当作无损的加速器，会低估它对运动质量的影响。对于视频也应先做重建实验：真实视频经过编解码后，快速运动和细小文字是否仍然稳定？

\subsection{视频生成与世界模型之间还隔着什么}

OpenAI 2024 年 Sora 技术披露说明了压缩视觉数据、时空 patches 和扩散 Transformer 的路线。披露同时展示了一些持续性与模拟能力，也报告了物理和长时间一致性的局限。它适合作为公开研究思路的案例，无法据此还原全部训练细节。\cite{a_sora2024}

一段看起来真实的物体运动，不等于模型已经掌握可供规划的真实动力学。视频数据告诉模型哪些画面经常相邻，却可能没有告诉它哪些变化由动作造成。要把模型用于决策，还需要回答反事实问题：如果施加另一个动作，后续状态如何变化？需要区分相机移动与物体移动，需要在新的环境中保持预测校准，还需要为不确定性留下空间。

例如，一只玻璃杯在许多视频里被完整拿起。模型据此生成“完整拿起”的画面很合理；但若机器人用过大的力抓杯子，真实结果可能是破裂。仅从常见画面学习到的统计规律，不一定覆盖关键而罕见的动作后果。一个能够辅助规划的世界模型，需要在行动条件下学习动态，而不仅是生成视觉上顺畅的故事。

Sora 2 的 2025 年官方材料披露同步音视频生成能力，但没有给出足以独立重建系统的完整网络和训练目标。本书将它作为产品能力的历史例子，不把其演示推断为特定未公开架构。按官方后续公告，Sora 网页与应用在 2026 年 4 月 26 日停止，API 公告停止日期为 2026 年 9 月 24 日；这个状态变化也说明电子书应区分技术贡献和服务可用性。\cite{a_sora2,a_soraend}

\section{统一多模态：理解与生成怎样共享计算}
\label{sec:application-unified}

\subsection{“能看图”和“能画图”为什么不自动统一}

视觉理解需要从像素中提取适合回答问题的语义。例如，回答“桌上有几个杯子”时，背景纹理大多不重要，物体数量却非常重要。视觉生成需要恢复可解码的细节，包括理解任务愿意忽略的纹理、颜色和局部位置。一种理解表征可能对光照变化保持不变，这对识别有帮助；但生成器如果需要改变光照，便不能把所有光照信息都抹去。

因此，统一多模态不只是把所有东西转成 token，然后宣布问题解决。它需要考虑至少三个接口：用于读取图像的表征，用于生成图像的表征，以及二者与文字上下文之间的交流方式。一个系统可以为理解和生成使用不同视觉编码器，同时在主干中共享上下文；也可以学习共同隐变量，再用不同解码器得到不同模态。

文本的离散概率目标与图像的连续回归目标也不同。一个常见混合训练的概念形式是
\begin{align}
\mathcal L={}&\lambda_T\,
\E[-\log p_\theta(w_j\mid w_{<j},\text{上下文})]\\
&+\lambda_I\,
\E[\norm{u_\theta(\ell_\tau,\tau,\text{上下文})
-(\ell_0-Z)}^2].
\label{eq:application-hybrid-loss}
\end{align}
本式是教学抽象，不是某一模型完整训练配方。$\lambda_T,\lambda_I$ 调节任务权重；$w_j$ 是第 $j$ 个文字 token。两项对共享参数提出不同要求，任务比例、数据质量和参数分工都会影响协作还是干扰。

\subsection{Transfusion：共享主干，保留不同生成机制}

Transfusion 把文字的下一 token 预测与图像的扩散目标放入同一个多模态模型，使用离散文字与连续图像表征，在图像块内部允许双向交互。这说明统一主干不要求所有模态都使用同一种输出分布。\cite{a_transfusion}

为什么图像内部的双向交流合理？图像的整个含噪状态在一次去噪调用时已经可见，我们要联合预测它的更新，没有必要人为规定“左上角必须先于右下角”。文字生成却通常需要遵守因果顺序，只让当前位置读取已经确定的前文。统一模型因而需要按模态调整注意力可见范围。

这是一条容易遗漏的训练原则。若文字目标能看到未来正确答案，训练误差会虚假降低，推理却不能复现；若图像被迫只有单向局部上下文，又会丢失扩散并行协同的一部分优势。注意力掩码必须与生成时实际可见的信息一致。

\subsection{BAGEL：两种视觉表征与专家分工}

BAGEL 的 2025 年研究采用文字自回归与视觉 Rectified Flow 的组合，使用理解用的视觉编码器和生成用的 VAE 表征。其 MoT 结构对理解与生成 token 使用不同参数专家，在统一注意力中交流，并用大规模交错多模态数据训练。\cite{a_bagel}

这里的“专家”不是几个会互相聊天的小模型。它指计算图中不同 token 类型使用不同的可训练参数。以这种方式，理解任务可保留语言与语义处理的能力，生成任务可拥有专门处理连续视觉状态的变换，而上下文仍能跨模态传播。总参数量、每个 token 激活的参数量与一次完整生成的计算量，是不同概念，不能混着比较。

交错数据指一条样本中可能出现文字、图像、再文字、再图像。例如，先给出人物照片，再说明要更换衣服，然后给出修改结果，随后继续要求改变背景。它与孤立的“图片加标题”配对相比，明确教模型如何把之前的视觉内容当成后续任务的上下文。系统想获得多轮编辑、跨图身份一致性和图文交替表达，就需要相应的数据组织，而不只是更多无关图片。

读者可以用一个小例子检验自己的理解：第一轮让模型画一只戴红围巾的猫，第二轮要求“让它坐在窗边，保留围巾”。第二轮条件不应只有一句孤立文字，还应包含第一轮的视觉信息。否则“它”指的是哪只猫，以及围巾的具体样子，都缺少稳定依据。多模态上下文提供记忆接口，模型是否有效利用仍需实际评估。

\subsection{2026 年的两个研究窗口}

MUNI 是 2026 年 6 月的预印本，研究共享随机隐变量下的任意模态到任意模态生成。它把模态编码器、解码器与共享流先验联合训练，并讨论跨模态一致性、子集条件的预测充分性和隐变量简洁性之间的关系。其理论结论依赖可实现性和优化等假设。\cite{a_muni}

这项研究引出的一个基础问题值得深入思考。假设我们要同时生成“雨中的街景图”和“雨声”，分别生成一张好街景、一个好雨声，只保证各自边缘分布。若图中雨很小，声音却是猛烈暴雨，它们没有表达同一个情境。共享隐变量希望承载共同场景，再由模态解码器表达各自细节。可共享信息太少，会出现不一致；共享信息太多，又可能把声音里的随机相位、图片里的随机纹理都强塞进同一个复杂变量。这是一个关于充分性与压缩的概率建模问题。

WorldBagel 是 2026 年 7 月的预印本，把 BAGEL 式统一模型扩展到视觉、语言、动作和未来观察，在机器人任务上研究统一训练的作用。它提出动作表征和交错任务组织，并报告若干指定环境中的实验；这些实验不能当作通用物理模拟已经实现的证据。\cite{a_worldbagel}

这里最有启发的是任务关系的改变。过去，文字到图像问“给定描述，什么画面合理”；行动世界建模问“给定当前观察和动作，什么后果合理”。它必须把动作放进条件，考虑多步误差与决策反馈。生成器面对的对象也从静态图像分布扩展到环境转移分布。两个预印本作为正在发展的研究方向阅读，比把它们当成已定型标准更合适。

\subsection{商业多模态系统：明确公开知识的边界}

OpenAI 在 2025 年发布的 4o 图像生成材料强调多模态上下文、文字渲染和编辑能力；2026 年 9 月 8 日的 ChatGPT Images 2.5 发布材料继续强调细节、精确编辑与多轮保持。这里我们引用的是官方产品披露，不能从表现推导出完整主干网络、噪声路径或训练损失。\cite{a_gptimage,a_images25}

同一个能力可以由不同组合实现：语言模型先形成语义计划再交给图像生成器，共享 Transformer 同时处理文字和连续视觉状态，或者通过其他尚未公开的混合设计。界面里的“原生多模态”说明的是系统层面的能力和使用方式，不是一条足够精确的数学定义。

评价这样的系统，可以研究可观察行为：多轮修改是否保留身份，文本信息是否准确，复杂关系是否满足，延迟是否符合任务需求。解释底层原理则应依赖论文、技术报告、公开代码和可复现实验。把这两类证据分开，既能欣赏系统实际进展，也能保持技术判断清楚。

\section{从声音到分子：扩散学习的不一定是像素}
\label{sec:application-beyond}

\subsection{音频：同一个条件，多种合理声音}

音频生成可在波形、时频谱或压缩音频隐变量中进行。AudioLDM 研究文本到音频的隐空间扩散，用音频表征建立生成条件，并把压缩表征解码为声音。它说明“压缩表示加条件生成”的思想可以从图像移植到别的感知模态。\cite{a_audioldm}

一条提示“雨点落在金属屋顶上”，不会规定每个雨点的准确落下时刻。因此，输出仍然是条件分布中的一个样本。音频的挑战包括持续时间、事件先后、节奏和声源一致性。若任务是语音，音素内容还必须精确；若任务是音乐，长时间结构与和声关系很重要。只让短窗口内的声音真实，并不能保证整段作品的组织合理。

时频谱把声音分解为时间和频率附近的能量信息，但从谱回到波形还需处理相位或依赖专门解码器。因而“把音频当成一张图片来生成”只是一种局部直觉，不能忽略声音特有的物理和时间结构。

\subsection{三维与分子：先选对几何空间}

三维对象可以表示为点云、网格、体素、隐式场或可渲染场景。表示一换，网络要遵守的结构也换了。点云对点的排列通常不敏感，三维旋转不应改变对象本身；网格还要考虑连接关系。第\ref{ch:math}章中直接向所有欧氏坐标加高斯噪声的例子，是一个重要起点，但不一定是所有几何任务的合适终点。

DiffDock 把分子对接写成配体位姿的生成问题，在平移、旋转和扭转自由度构成的空间上建立扩散过程，并配合置信度估计。关键变化是：生成变量不是一幅分子图片，而是实际几何自由度。\cite{a_diffdock}

旋转矩阵受到正交性和行列式约束，不能把矩阵的每个元素任意扰动后，仍当作合法旋转。周期性的扭转角也不能简单按无边界实数处理。这样的例子提醒我们，score、噪声和速度必须尊重变量所在空间的几何；术语可参考\hyperref[gl:manifold]{流形}与\hyperref[gl:equivariance]{等变性}条目。

RFdiffusion 将蛋白质结构设计表述为逐步去噪的生成任务，并研究基序支架、结合体等条件设计。论文包含实验验证的设计案例，因此展示了从生成结构到物理实验验证的完整研究链条。\cite{a_rfdiffusion}

这里的质量标准明显不同于图像。一个结构外观像蛋白质，并不意味着它会折叠成预期形态，更不意味着它具有所需功能。计算预测、结构筛选和实验测量分别提供不同证据。对于科学应用，生成器提出候选，验证流程决定哪些候选可信。扩散模型扩展了可以探索的设计空间，却没有替代领域中的约束与检验。

\section{怎样评价一种进展}
\label{sec:application-evaluation}

\subsection{一个分数不能承担所有结论}

\term{FID}{gl:fid}在固定特征空间里比较生成与真实样本的统计量，它对分布偏差有用，却不直接回答一张图是否符合某条复杂指令。图文相似度可以检查总体语义，却可能漏掉数量、左右关系和文字拼写。人类偏好评测更贴近视觉体验，但结果会受到提示集合、评价问题、样本选择和比较对手影响。

表\ref{tab:application-evaluation}按任务而非模型列出评价问题。它是一份原创检查框架，不是论文里的性能排行榜。一个实际实验应先选定任务，再确定哪些错误最有代价，然后选择能显露这些错误的指标。

\begin{table}[htbp]
\centering
\small
\begin{tabularx}{\linewidth}{@{}p{2.1cm}XX@{}}
\toprule
任务 & 最该检查的问题 & 单一视觉质量分数容易漏掉什么\\
\midrule
文字到图像 & 对象、数量、属性和关系是否同时满足 & 正确构图中的错误数量或绑定\\
局部编辑 & 指令是否完成，其他部分是否保持 & 漂亮但更换了主体和背景\\
空间控制 & 生成是否遵守骨架、深度或边缘 & 纹理逼真但几何偏离\\
视频 & 身份、运动、遮挡和长期状态是否连续 & 每帧清晰但跨帧闪烁\\
多模态生成 & 各模态是否表达同一事件与时间结构 & 图片与声音各自合理却不匹配\\
科学设计 & 约束、可实现性与实验功能是否成立 & 看似合理的结构不能实际工作\\
\bottomrule
\end{tabularx}
\caption{从任务出发选择评价问题。评价应观察完整条件任务，而不只是输出是否好看。}
\label{tab:application-evaluation}
\end{table}

一个公平的生成比较，还需要报告采样预算。模型一次生成一张图，与先生成二十张再筛选最好的一张，体现的是不同计算预算。采样步数、随机种子数、提示词重写、引导强度和后处理都应清楚说明。论文展示一组很好的图，只能证明系统能产生这些图；稳定成功率需要在有代表性的提示集合上验证。

\subsection{用消融实验定位贡献}

消融实验是有控制地改变一个组件，观察结果的实验。若新系统同时更换文字编码器、扩大数据、增加模型宽度和改变损失，最终变好，不能直接说全部收益来自新损失。一个有说服力的结论需要尽可能固定其他条件，或者明确承认组件间耦合。

读论文时，可沿着一条具体问题链阅读：作者先测量了什么失败？新方法改变哪个组件？对照是否只有该组件不同？收益在哪些任务成立？代价是更多训练还是更多推理？是否公开了失败样例？这些问题比“模型名字里有没有最新术语”更能帮助判断一个进展。

\section{一条可走通的论文阅读路线}
\label{sec:application-reading}

下面选择二十篇具有明确教学价值的论文或技术报告。它们并非二十个竞争榜单位置，而是二十个不同的问题窗口。第二、三章的 DDPM、DDIM、score-SDE 和 Flow Matching 奠定数学前提；读到本章时，优先看问题设置、方法图和消融表，再回到推导。每条入口直接链接一手来源。

\begin{enumerate}[leftmargin=*,itemsep=0.65em]
\item \textbf{ADM，2021。}问题：扩散的图像先验怎样与强生成基线竞争？关注结构消融和引导的质量、覆盖率取舍；局限是类别条件并不等同于开放语言控制。\href{https://arxiv.org/abs/2105.05233}{阅读入口}。\cite{a_adm}
\item \textbf{GLIDE，2021。}问题：怎样在去噪过程中遵循文字？关注条件表示、引导比较和补全设置；不要把公开的小模型表现当成全部论文实验。\href{https://arxiv.org/abs/2112.10741}{阅读入口}。\cite{a_glide}
\item \textbf{Imagen，2022。}问题：图像模型之外，还有什么限制图文一致性？优先读文本编码器与扩容实验，再看级联生成；复杂关系仍需要专门评估。\href{https://arxiv.org/abs/2205.11487}{阅读入口}。\cite{a_imagen}
\item \textbf{DALL-E 2，2022。}问题：语义表征怎样成为中间生成变量？先画清先验和解码器的分工；语义接口会压缩部分细节。\href{https://arxiv.org/abs/2204.06125}{阅读入口}。\cite{a_dalle2}
\item \textbf{LDM，2021／2022。}问题：高分辨率扩散能否在更紧凑空间中进行？先检查自编码器重建，再读条件机制；压缩损失构成能力边界。\href{https://arxiv.org/abs/2112.10752}{阅读入口}。\cite{c_rombach2022ldm}
\item \textbf{DreamBooth，2022。}问题：怎样让已有先验记住少量照片中的主体？关注身份与类别泛化；数据少时容易把背景学成身份。\href{https://arxiv.org/abs/2208.12242}{阅读入口}。\cite{a_dreambooth}
\item \textbf{InstructPix2Pix，2022。}问题：怎样学习“照着指令改图”？关注训练配对从何而来以及保持能力；合成配对的偏差会影响可学编辑。\href{https://arxiv.org/abs/2211.09800}{阅读入口}。\cite{a_instructpix}
\item \textbf{DiT，2022／2023。}问题：Transformer 能否成为可扩展的扩散主干？重点读条件注入和计算规模实验；原始类别实验不能直接推出所有文图任务。\href{https://arxiv.org/abs/2212.09748}{阅读入口}。\cite{a_dit}
\item \textbf{ControlNet，2023。}问题：怎样利用成熟先验满足空间条件？追踪可训练分支怎样进入主干；约束服从与视觉质量需要分别评估。\href{https://arxiv.org/abs/2302.05543}{阅读入口}。\cite{a_controlnet}
\item \textbf{AudioLDM，2023。}问题：隐空间条件生成怎样迁移到声音？关注条件表征与声音解码接口；长时间事件组织不能由短片段质量替代。\href{https://arxiv.org/abs/2301.12503}{阅读入口}。\cite{a_audioldm}
\item \textbf{SDXL，2023。}问题：如何把隐空间系统推进到更高分辨率？看尺寸条件和组件分工；多项改动需要消融定位收益。\href{https://arxiv.org/abs/2307.01952}{阅读入口}。\cite{a_sdxl}
\item \textbf{SD3／MM-DiT，2024。}问题：如何让文字和图像在流模型里共同计算？区分路径实验与结构实验；比较损失大小时注意目标和权重不同。\href{https://arxiv.org/abs/2403.03206}{阅读入口}。\cite{a_sd3}
\item \textbf{Video Diffusion Models，2022。}问题：从图像边缘分布走向视频联合分布，需要改变什么？关注时空交互与联合训练；短时一致性不等于长期动力学。\href{https://arxiv.org/abs/2204.03458}{阅读入口}。\cite{a_videodiff}
\item \textbf{Wan，2025。}问题：开放的视频基础模型怎样组织编解码、数据和训练？报告较长，可先读模型设计；具体版本与评测条件必须一起记录。\href{https://arxiv.org/abs/2503.20314}{阅读入口}。\cite{a_wan}
\item \textbf{Transfusion，2024。}问题：自回归文字与连续视觉生成怎样共享主干？先看注意力掩码和两个目标；统一参数仍面对任务平衡。\href{https://arxiv.org/abs/2408.11039}{阅读入口}。\cite{a_transfusion}
\item \textbf{BAGEL，2025。}问题：大规模交错数据怎样支持理解与生成？关注专家、两种视觉表征和上下文；任务能力应与数据组织一同讨论。\href{https://arxiv.org/abs/2505.14683}{阅读入口}。\cite{a_bagel}
\item \textbf{MUNI，2026，预印本。}问题：多模态共同生成怎样保持一致？先读潜变量标准，再检查理论假设；基准结果不能替代所有模态上的证明。\href{https://arxiv.org/abs/2606.16408}{阅读入口}。\cite{a_muni}
\item \textbf{WorldBagel，2026，预印本。}问题：统一模型怎样加入动作和未来状态？关注转移任务与动作表征；指定机器人环境中的收益有清楚适用范围。\href{https://arxiv.org/abs/2607.03461}{阅读入口}。\cite{a_worldbagel}
\item \textbf{DiffDock，2022／2023。}问题：存在多个合理配体位姿时，怎样生成并筛选？关注几何空间和置信度；结构误差与实际药物功能属于不同证据。\href{https://arxiv.org/abs/2210.01776}{阅读入口}。\cite{a_diffdock}
\item \textbf{RFdiffusion，2023。}问题：怎样在条件下生成可验证的蛋白质结构？先看设计任务及实验链条；生成候选仍需结构与功能检验。\href{https://www.nature.com/articles/s41586-023-06415-8}{阅读入口}。\cite{a_rfdiffusion}
\end{enumerate}

读完其中任意一篇，可以尝试用五句话复述：它生成什么随机变量，如何构造训练样本，网络如何接收条件，推理如何更新状态，实验排除了哪一种替代解释。若五句话还说不清，继续查结构图和任务设置，通常比再看十张示例图更有帮助。

\section{未来方向：从已知瓶颈提出判断}
\label{sec:application-future}

以下是基于本章讨论的研究判断，不是关于下一代产品的已知事实，也不是保证会实现的时间表。

第一，\textbf{更准确的表示接口}可能与更大的主干同样重要。小字、快速动作、几何细节若在编码器中丢失，后面的场学习无法稳定补救。未来的表示可能按内容分配不同粒度，使平坦背景用较少 token，文字和关键动作获得更细的表征。但自适应接口还必须保持可训练、可解码和高效批处理，不能只在概念上节省 token。

第二，\textbf{生成计算可能更依赖任务难度}。简单背景也许不需要与复杂空间关系同样多的推理预算。第三章的蒸馏、少步生成和求解器改进提供了加速基础；另一条路线是让模型检查约束，发现失败后继续修改。难点在于检查器也会犯错，而且多次生成再选择的额外成本应当诚实计算。

第三，\textbf{多轮编辑需要可维护的视觉状态}。一幅图连续修改十次，哪些对象身份必须保持，哪些属性已经被用户改变，不能仅靠一个不断增长的文字摘要猜测。场景结构、对象引用和连续视觉表征之间的配合，可能成为更可靠的编辑接口。但显式结构并非越多越好：过硬的接口也可能妨碍风格化和创作自由。

第四，\textbf{多模态统一将继续面对共享与专门化的取舍}。所有模态共用一套参数，结构简洁却可能产生任务干扰；为每个模态完全独立建模，又浪费共同知识。共享注意力、专门参数、不同输出分布和统一上下文的组合，是一个开放设计空间。真正有意义的进展应当证明新增协作改善了任务，而不是仅仅让结构图看起来统一。

第五，\textbf{世界建模需要行动和可检验的因果约束}。视频生成带来丰富视觉先验，但用于机器人和科学预测时，误差不再只是画面瑕疵。未来系统需要识别分布外状态，表达多种后果，在闭环中接受真实反馈。把动态模型当成一个允许检验和修正的预测工具，会比把漂亮视频直接等同于真实世界理解更有成效。

这些方向最终又回到本书的第一句话：生成是在学习分布。应用越丰富，分布的对象、条件和评价越具体。扩散与流模型提供了一个强大框架，真正完成任务还需要正确的表示、可学习的条件、合适的数据，以及诚实的验证。

\section{练习与解答思路}
\label{sec:application-exercises}

\begin{exercisebox}{练习 1：定位小字错误}
某文字到图像系统在 $1024\times1024$ 图像中总把商品标签的小字画成乱码。设计三个实验，分别判断错误是否来自自编码器、文字条件或生成先验。
\end{exercisebox}

\begin{exercisebox}{练习 2：计算 token 的代价}
隐变量尺寸为 $64\times64\times16$。将 patch 边长从 $4$ 改成 $2$，token 数、注意力矩阵元素数各改变多少倍？如果图像空间边长再翻倍且压缩率不变，这两项又改变多少倍？说明为什么这些倍率不等于完整推理时间倍率。
\end{exercisebox}

\begin{exercisebox}{练习 3：联合分布与逐帧质量}
构造一个最小例子：每帧都有同样的合理边缘分布，但两种视频生成器的时间一致性完全不同。解释为什么只计算单帧指标不能区分它们。
\end{exercisebox}

\begin{exercisebox}{练习 4：为统一模型设计掩码}
一条训练样本依次包含文字说明、输入图像、编辑指令和目标图像。目标图像用于速度回归，后续文字用于自回归回答。说明哪些 token 应看见哪些信息，并指出一种会导致训练时信息泄漏的设计。
\end{exercisebox}

\begin{exercisebox}{练习 5：把未来方向写成实验}
选择“多轮身份保持”“图像和声音一致”或“动作条件世界预测”之一，提出一个有控制的实验。给出训练条件、比较基线、评价指标和能够推翻你假设的结果。
\end{exercisebox}

\paragraph{练习 1 的思路。}
先对真实标签图做编解码重建，检查字形是否保留；再用内容只有一个大字、一个短词和一条明确关系的提示检验条件理解；最后固定编码器与文字编码器，比较数据或生成主干变体。若重建就失败，先验实验无法单独解决接口缺失。评价文字可用转录准确率与人工核验，不只用整体图像相似度。

\paragraph{练习 2 的思路。}
token 数从 $256$ 变成 $1024$，增加四倍；注意力矩阵元素数增加十六倍。固定 patch 后空间边长翻倍，token 数又增加四倍，注意力矩阵又增加十六倍。线性投影、前馈层、自编码器、采样调用次数和硬件利用率也影响时间，因此完整延迟需要测量。

\paragraph{练习 3 的思路。}
令每帧只有红球或蓝球，概率各为 $1/2$。系统 A 每帧独立选色；系统 B 在第一帧选色后一直保持。所有单帧边缘分布相同，但 A 频繁闪烁，B 始终一致。可用相邻帧颜色改变率区分，说明联合指标不可由单帧指标替代。

\paragraph{练习 4 的思路。}
目标含噪图像可以读取先前说明、输入图像和编辑指令，并在目标图像内部双向交流。文字预测应遵守其实际推理顺序。若让当前文字读取尚未生成的正确未来文字，便发生直接泄漏；若后续任务在训练中读取不可能在推理时获得的干净目标表示，也会形成条件分布不一致。需先明确推理顺序，再设计掩码。

\paragraph{练习 5 的思路。}
以多轮编辑为例，固定主干、训练预算和提示集合，比较有交错数据与仅孤立配对数据的系统；记录每轮指令成功率、主体保持和无关区域变化。若加入交错数据后只让画面更好看，却没有提高保持或任务完成率，原假设未获支持。报告全部轮次和失败，而不只选最佳结果。

\chaptertakeaway{应用发展的主线不是名字不断变新，而是接口不断扩展：U-Net 协调不同图像尺度，隐空间降低处理维度，DiT 把视觉场写成 token 交互，空间条件使生成可控，视频要求联合时间建模，统一多模态让理解与生成共享上下文。无论生成像素、声音、位姿还是蛋白质，都应先选对随机变量与几何空间，再决定条件、网络和验证方法。}
}{}
\appendix
\IfFileExists{appendices/01-foundations.tex}{\chapter{概率工具：读懂分布与回归}\label{app:probability}
这个附录不是概率论的完整替代，而是正文使用的工具箱。每个概念都围绕扩散模型里的具体用法展开。读到某个式子卡住时，可以先查对应工具，再返回原文。

\section{随机变量、分布与密度}
随机变量是从随机结果到数值或向量的映射。例如抽取一张训练图像，把它展平成 $d$ 维向量，便得到随机向量 $X$。分布告诉我们 $X$ 落在一个集合 $A$ 中的概率，记为 $P(X\in A)$。分布可以存在而没有相对于普通体积的密度；例如只落在一条曲线上的向量就可能如此。

若分布对 Lebesgue 体积绝对连续，则存在非负函数 $p$ 满足
\begin{equation}
 P(X\in A)=\int_Ap(x)\dd x,\qquad\int_{\R^d}p(x)\dd x=1.
\end{equation}
$p(x)$ 不是单点概率。连续变量满足 $P(X=x)=0$，但其附近小区域的概率可以为正。在同一维度和同一坐标尺度下，密度越大，一小块等体积区域的概率通常越大。换坐标会改变体积单位，因而改变密度数值。

实际数字图像有有限精度，原本是离散分布。为应用连续模型，可以归一化像素并作去量化，例如给整数像素加入区间均匀噪声。另一种处理是用连续模型的密度在像素量化区间上积分，得到离散概率。直接将一个连续密度值当成量化图像的概率质量，会把两种对象混在一起。严格报告 bits/dim 时还须明确去量化方式与尺度修正。

\section{联合、条件、边缘与贝叶斯公式}
两个随机变量的联合密度 $p(x,y)$ 同时描述二者。边缘密度通过积分获得：$p(x)=\int p(x,y)\dd y$。当 $p(y)>0$ 时，条件密度是 $p(x\mid y)=p(x,y)/p(y)$。于是
\begin{equation}
 p(x\mid y)=\frac{p(y\mid x)p(x)}{p(y)}.
\end{equation}
这是贝叶斯公式。扩散中 $p(y\mid x)$ 是已知的加噪核，$p(x)$ 是数据分布，$p(x\mid y)$ 则描述哪些原图可能导致当前含噪观测。

这里容易出现一种误解：训练时我们知道“这个 $y$ 是由哪张 $x$ 生成的”，不代表仅给定 $y$ 就能在数学上唯一识别 $x$。训练配对提供一个条件后验的随机样本；网络只能根据输入可用的信息作预测。

条件期望 $\E[V\mid U]$ 本身是 $U$ 的函数。它满足塔式法则
\begin{equation}
 \E[V]=\E\bigl[\E[V\mid U]\bigr].
\end{equation}
更一般地，对任何有界的 $U$ 的函数 $h$，或使下列相关乘积可积的 $h$，
\begin{equation}
 \E[h(U)V]=\E[h(U)\E[V\mid U]].
\end{equation}
正文多次使用的“把条件标签平均成边缘对象”正是这个法则的向量版本。需要平方回归时还要求相关变量二阶可积。

\section{平方回归为什么得到条件均值}
令 $m(U)=\E[V\mid U]$，展开平方：
\begin{align}
 \norm{h(U)-V}^2=&\norm{h(U)-m(U)}^2+\norm{m(U)-V}^2\\
 &+2\ip{h(U)-m(U)}{m(U)-V}.
\end{align}
给定 $U$ 时，交叉项的条件期望是零。所以平方风险分成“预测距离最优函数的误差”和“观察不到的随机性所造成的不可消除方差”。这不仅证明最优函数为条件均值，也解释了为什么最佳损失一般不等于零。

如果监督是图像 $X$，条件平均可能把不同合法图像叠加成模糊图像。如果监督是局部速度或 score，平均并不会立即成为最终输出；它成为概率动力学里的场。任务的语义与损失的数学形式需要一起看，不能笼统地说“均方误差只能生成模糊图像”。

\section{高斯向量与条件公式}
若 $X\sim\N(m,C)$，其中 $C$ 正定，则
\begin{equation}
 p(x)=\frac{\exp[-\tfrac12(x-m)^\top C^{-1}(x-m)]}{(2\pi)^{d/2}\det(C)^{1/2}}.
\end{equation}
仿射变换 $AX+b$ 仍是高斯，均值为 $Am+b$，协方差为 $ACA^\top$。两个独立高斯之和仍是高斯，均值和协方差分别相加。仅仅“不相关”一般不足以保证独立；在联合高斯情况下，不相关等价于独立。

对联合高斯
\begin{equation}
 \begin{bmatrix}X\\Y\end{bmatrix}\sim\N\left(
 \begin{bmatrix}m_X\\m_Y\end{bmatrix},
 \begin{bmatrix}C_{XX}&C_{XY}\\C_{YX}&C_{YY}\end{bmatrix}\right),
\end{equation}
若 $C_{YY}$ 可逆，则
\begin{align}
 \E[X\mid Y=y]&=m_X+C_{XY}C_{YY}^{-1}(y-m_Y),\\
 \Cov(X\mid Y)&=C_{XX}-C_{XY}C_{YY}^{-1}C_{YX}.
\end{align}
一种证明方式是令 $R=X-m_X-C_{XY}C_{YY}^{-1}(Y-m_Y)$。直接计算 $\Cov(R,Y)=0$；联合高斯性使 $R$ 与 $Y$ 独立，因此上面给出条件均值与残差方差。DDPM 后验也可用这一公式求得，正文选择配方是为了让系数来源更清楚。

\section{KL、似然与变分下界}
KL 散度定义为
\begin{equation}
 \KL(q\Vert p)=\E_q\log\frac{q(X)}{p(X)}.
\end{equation}
它要求 $q$ 对 $p$ 绝对连续；若 $q$ 给某处正概率而 $p$ 给零概率，KL 可以是无穷。KL 非负，但通常不对称，也不满足三角不等式，不能称为一般意义上的距离。

由 $-\log$ 的凸性，
\begin{equation}
 \KL(q\Vert p)=-\E_q\log\frac pq
 \ge-\log\E_q\frac pq\ge0,
\end{equation}
在同支撑、规范化等常见条件下最后的期望为 $1$。等号只在两个分布几乎处处相等时成立。

两个高斯的 KL 可以把对数密度相减并用二次型期望计算：
\begin{align}
 \KL(\N(m_q,C_q)\Vert\N(m_p,C_p))
 =\tfrac12\bigl[&\tr(C_p^{-1}C_q)
 +(m_p-m_q)^\top C_p^{-1}(m_p-m_q)\\
 &-d+\log\tfrac{\det C_p}{\det C_q}\bigr].
\end{align}
这里的基本恒等式是 $\E[(X-m)^\top A(X-m)]=\tr(AC)$。它解释了正文中的均值误差系数与方差项。

似然把参数看作可变对象，样本看作已观测数据：$\theta\mapsto p_\theta(x)$。最大似然训练最大化训练数据的平均对数似然。对真实分布而言，
\begin{equation}
 \E_{\pdata}[-\log p_\theta(X)]
 =H(\pdata)+\KL(\pdata\Vert p_\theta),
\end{equation}
在这些量存在时，最小负对数似然等价于最小这个方向的 KL。

若模型有隐藏变量 $Z$，令任意辅助条件密度 $q(z\mid x)$ 满足相关支撑条件，定义
\begin{equation}
 \mathrm{ELBO}(x)=\E_{q(z\mid x)}[\log p_\theta(x,z)-\log q(z\mid x)].
\end{equation}
利用 $p_\theta(x,z)=p_\theta(z\mid x)p_\theta(x)$，直接得到
\begin{equation}
 \log p_\theta(x)-\mathrm{ELBO}(x)
 =\KL(q(z\mid x)\Vert p_\theta(z\mid x))\ge0.
\end{equation}
这展示了下界的松紧来源。DDPM 的隐藏变量是整条含噪轨迹，辅助分布是前向链；并不是在每个时刻单独创建一个无关的 VAE。

\section{推前、耦合与最优输运}
给定映射 $T$，令 $Y=T(Z)$。它的分布记为 $T_\#p_Z$，意思是
\begin{equation}
 P(Y\in A)=P(Z\in T^{-1}(A)).
\end{equation}
这叫推前分布。它不要求 $T$ 可逆。若 $T$ 可逆且光滑，则变化量公式为
\begin{equation}
 p_Y(y)=p_Z(T^{-1}(y))\left|\det\frac{\partial T^{-1}}{\partial y}\right|.
\end{equation}
因此 GAN 的确定性生成器也能推前随机噪声；其难点不是“确定性就不能随机”，而是训练与密度计算方式。

耦合 $\pi(z,x)$ 是具有指定两个边缘的联合分布。相同的起点、终点分布可以存在许多耦合：随机配对、按距离配对、模型流产生的配对都可能不同。最优输运在这些耦合里最小化一个代价，例如
\begin{equation}
 \inf_{\pi\in\Pi(p_Z,p_X)}\E_\pi\norm{X-Z}^2.
\end{equation}
这个最小值是平方 Wasserstein-2 距离（要求有限二阶矩）。选择直线插值只是给某个耦合选路径；它本身不负责把耦合优化到全局最优。小批量最优配对也不自动等于整个数据分布的最优耦合。

\section{熵、归一化与有限样本}
离散熵是 $-\sum_xp(x)\log p(x)$；连续微分熵是 $-\int p(x)\log p(x)\dd x$。微分熵会随坐标缩放变化，甚至可以为负。扩散中“加噪增加不确定性”的直觉不意味着任意带漂移过程的微分熵都单调增加；漂移可能压缩空间。

真实数据分布不是训练集的经验分布本身。经验分布是各样本位置上的点质量混合，常写 $\widehat p=\frac1n\sum_i\delta_{x_i}$。高斯平滑后它变成可微核密度；只要模型容量与优化允许，网络也可能记住这些样本的局部结构。去噪训练并没有在数学上禁止记忆。泛化还依赖数据、归纳偏置、正则化及训练规模。

\chapter{连续动力学：ODE 与随机微积分入门}\label{app:stochastic}
\section{ODE 是随时间变化的向量场}
普通微分方程 $\dot x_t=v(x_t,t)$ 表示粒子当前位置的速度由场决定。初始条件 $x_0=z$ 选定一条轨迹。如果 $v$ 连续并在状态变量上局部 Lipschitz，就有局部存在唯一性；适当的增长控制可保证有限时间内不爆炸。更强的光滑条件使流映射对初始状态可微。

Euler 方法 $x_{t+h}=x_t+h v(x_t,t)$ 用当前速度近似整段移动，局部截断误差通常为 $O(h^2)$、全局误差为 $O(h)$，前提是足够光滑并稳定。Heun 方法先预测下一点，再平均两个端点的速度，满足条件时有二阶全局精度。阶数描述渐近误差随步长变化的规律，不代表在任意粗步长或任意网络场下都更好。

刚性粗略指不同方向或区域具有很不相同的变化时间尺度，迫使某些显式算法使用很小步长来保持稳定。生成模型低噪声端的场可能变得陡峭，强引导也可能放大场幅度或敏感性。不能仅凭“轨迹看起来平滑”排除数值困难。

\section{布朗运动与 SDE 的含义}
标准布朗运动 $W_t$ 从零开始，路径连续，独立增量满足 $W_t-W_s\sim\N(0,(t-s)I)$。路径连续不意味着可微；它几乎必然处处不能按普通方式求导。形式上的“白噪声” $\dot W$ 应理解为广义随机对象。

SDE
$\dd X_t=f(X_t,t)\dd t+G(X_t,t)\dd W_t$
是积分方程的简写：
\begin{equation}
 X_t=X_0+\int_0^t f(X_r,r)\dd r+\int_0^t G(X_r,r)\dd W_r.
\end{equation}
正文使用 It\^o 随机积分，离散近似在每段左端取系数。对适应于过去信息、平方可积的系数，随机积分满足均值零与等距性质：
\begin{equation}
 \E\norm{\int_0^tG_r\dd W_r}^2
 =\E\int_0^t\norm{G_r}_{F}^{2}\dd r.
\end{equation}
这里 Frobenius 范数的平方是矩阵所有元素平方和。It\^o 与 Stratonovich 是两种不同积分约定；若扩散系数依赖状态，漂移转换会出现额外修正。正文的 $g(t)I$ 不依赖状态，避免了该处歧义。

\section{It\^o 二阶项从哪里来}
对一维小增量 $\Delta X=f\Delta t+g\Delta W$，Taylor 展开为
\begin{equation}
 \Delta\varphi=\varphi'\Delta X+\tfrac12\varphi''(\Delta X)^2+\cdots.
\end{equation}
由于 $\Delta W$ 的标准差是 $\sqrt{\Delta t}$，$(\Delta W)^2$ 在累积意义上留下 $\Delta t$ 的贡献；$(\Delta t)^2$ 与 $\Delta t\Delta W$ 则是更高阶。于是得到 $\tfrac12g^2\varphi''\dd t$。多维情况用 Hessian 与 $GG^\top$ 收缩：
\begin{align}
 \dd\varphi(X_t,t)=&\left(\partial_t\varphi+f\cdot\grad\varphi
 +\tfrac12\tr(GG^\top\nabla^2\varphi)\right)\dd t\\
 &+(\grad\varphi)^\top G\dd W_t.
\end{align}
这是正文由 SDE 得到密度方程的基础。形式记忆法 $\dd W_i\dd W_j=\delta_{ij}\dd t$ 是对二次变差的简写，不能当成普通代数里每次随机增量平方恰等于时间的断言。

\section{弱形式为何适合概率路径}
如果密度或场不够光滑，逐点微分方程未必成立。我们改用测试函数积分来定义方程。对于连续性方程，弱形式为
\begin{equation}
 \frac{\dd}{\dd t}\E[\varphi(X_t)]
 =\E[\grad\varphi(X_t)\cdot v_t(X_t)].
\end{equation}
它表达所有可观測平滑量的期望如何变化。流匹配对固定端点条件下的路径可能只是一个移动的点质量，但对测试函数作求导仍然有意义。再把端点混合，就得到边缘输运方程。这也是正文选择测试函数证明而不强行给每条直线写普通密度的原因。

\section{哪些地方应保留小的正噪声端点}
当真实数据严格位于低维集合上时，零噪声端的普通 score 可能不存在。即使数据有密度，$b(t)\to0$ 时的 $\epsilon/b(t)$ 也可能具有高方差，数值积分容易不稳定。实践中常只积分到 $t_{\min}>0$，再作去噪估计或解码。理论端点、训练抽样端点与采样端点应分别说明。

}{}
\IfFileExists{appendices/02-proofs.tex}{\chapter{补充推导与公式核验}\label{app:proofs}
本附录补全正文没有展开到底的步骤，也提供几种独立的符号检查方式。它与第二章使用同一时间约定。

\section{ELBO 的望远镜重排}
直接由负下界得到
\begin{align}
 L=&\E_q\left[-\log p(x_K)-\sum_{k=1}^K\log p_\theta(x_{k-1}\mid x_k)
 +\sum_{k=1}^K\log q(x_k\mid x_{k-1})\right].
\end{align}
对 $k\ge2$，贝叶斯公式给出
\begin{equation}
 \log q(x_k\mid x_{k-1})=
 \log q(x_{k-1}\mid x_k,x_0)+\log q(x_k\mid x_0)-\log q(x_{k-1}\mid x_0).
\end{equation}
最后两个对数对 $k$ 求和会相互抵消，留下 $\log q(x_K\mid x_0)-\log q(x_1\mid x_0)$。再与 $k=1$ 的 $\log q(x_1\mid x_0)$ 抵消，便得到正文式\eqref{eq:elbo}的三个部分。每个 KL 的条件变量和期望分布必须一并保留；只写一个抽象的“KL 求和”容易漏掉重建项。

\section{经典 score matching 的积分分部}
在密度 $p$ 足够光滑、边界项消失的情况下，考虑
\begin{equation}
 J(\theta)=\tfrac12\E_p\norm{s_\theta(X)-\grad\log p(X)}^2.
\end{equation}
展开并略去与参数无关的部分，
\begin{equation}
 J=\tfrac12\int p\norm{s_\theta}^2\dd x
 -\int s_\theta\cdot\grad p\dd x+C.
\end{equation}
逐坐标积分分部得到
\begin{equation}
 -\int s_\theta\cdot\grad p\dd x=\int p\divg s_\theta\dd x.
\end{equation}
所以
\begin{equation}
 J(\theta)=\E_p\left[\tfrac12\norm{s_\theta(X)}^2+\divg s_\theta(X)\right]+C.
\end{equation}
未知的真实 score 从可优化项中消失。这是经典 score matching 的关键结构\cite{r_scorematching}。但高维网络需要计算散度，未加噪数据也可能集中在低维流形。去噪 score matching 改用已知平滑核的梯度，不需要直接计算这项散度，并在多个噪声尺度学习平滑后的场。

\section{反向漂移的局部贝叶斯计算}
先假设 $g$ 仅依赖时间且局部看作常数，$f$ 和密度足够光滑。令 $h>0$，给定当前状态 $X_t=x$，前一状态写作 $y=x+\delta$。小步转移密度近似为
\begin{equation}
 q_h(x\mid y)\propto\exp\left[-\frac{\norm{x-y-hf(y,t-h)}^2}{2hg^2}\right].
\end{equation}
逆条件密度正比于 $p_{t-h}(x+\delta)q_h(x\mid x+\delta)$。主导的 $\delta$ 尺度为 $\sqrt h$。对对数先验作展开：
\begin{equation}
 \log p_{t-h}(x+\delta)
 =\log p_t(x)+s_t(x)\cdot\delta+O(h)+O(\norm\delta^2).
\end{equation}
保留决定均值一阶项的部分，转移指数为
\begin{equation}
 -\frac{\norm{\delta+hf(x,t)}^2}{2hg^2}
 +s_t(x)\cdot\delta.
\end{equation}
配方后，均值为 $-hf(x,t)+hg^2s_t(x)$，协方差主导项为 $hg^2I$。漂移导数和密度二阶项对均值、协方差给出更高阶修正；使该展开严格需控制余项和尾部。该计算展示了反向 drift 的来源，不代替一般反向扩散定理的完整条件证明。

\section{矩阵扩散且依赖状态时的修正}
对于
\begin{equation}
 \dd X_t=f(X_t,t)\dd t+G(X_t,t)\dd W_t,
 \qquad A(x,t)=G(x,t)G(x,t)^\top,
\end{equation}
Fokker--Planck 方程为
\begin{equation}
 \partial_tp=-\sum_i\partial_i(f_ip)
 +\tfrac12\sum_{i,j}\partial_i\partial_j(A_{ij}p).
\end{equation}
定义向量 $(\divg A)_i=\sum_j\partial_jA_{ij}$。将二阶项展开为概率流，得到概率流速度
\begin{equation}
 v^{\mathrm{PF}}=f-\tfrac12\divg A-\tfrac12A\grad\log p.
\end{equation}
使用 $t$ 向后走的反向 SDE 漂移则为
\begin{equation}
 f^{\mathrm{reverse}}=f-\divg A-A\grad\log p.
\end{equation}
相应的逆向正时钟漂移要整体取负号。正文 $A=g(t)^2I$ 不依赖状态，$\divg A=0$，才化成熟悉的标量公式。把状态依赖扩散硬套进 $f-g^2s$ 会漏掉散度修正。这一点可对照 Score SDE 原文附录 A\cite{r_sde}。

\section{通过流映射的 Jacobian 检查密度符号}
设 $T_t(z)$ 是 ODE 流，$J_t=\partial_zT_t$。对初始状态求导：
\begin{equation}
 \dot J_t=(\partial_xv_t)(T_t)J_t.
\end{equation}
利用行列式的导数公式，
\begin{equation}
 \frac{\dd}{\dd t}\log|\det J_t|
 =\tr(J_t^{-1}\dot J_t)=\divg v_t(T_t).
\end{equation}
变化量公式 $p_t(T_t(z))|\det J_t(z)|=p_0(z)$ 因而给出
\begin{equation}
 \frac{\dd}{\dd t}\log p_t(T_t(z))=-\divg v_t(T_t(z)).
\end{equation}
空间体积膨胀时密度降低，这与正文的连续性方程推导一致。这个证明要求可微可逆流；奇异端点则需要单独处理，不能直接使用可逆变化量公式。

\section{一维高斯概率流的完整解析轨迹}
设边缘路径为 $p_t=\N(m_t,h_t)$，$h_t>0$。选择速度
\begin{equation}
 v_t(x)=m_t'+\frac{h_t'}{2h_t}(x-m_t).
\end{equation}
令 $x_t=m_t+\sqrt{h_t}\,z$，直接求导即验证它满足此 ODE。若 $z\sim\N(0,1)$，边缘恰为 $p_t$。因此给定任意可微高斯均值和正方差路径，生成速度由“移动中心”与“缩放偏离中心的位置”两项组成。

对于 VP 前向过程，$m_t=a(t)m_0$、$h_t=a(t)^2h_0+b(t)^2$。取 $s_t=-(x-m_t)/h_t$，代入 $f-g^2s/2$，利用 $m_t'=\gamma m_t$ 与 $h_t'=2\gamma h_t+g^2$，得到同一速度。这个独立检查能同时核验 score 符号和概率流的半系数。

\section{CFM 等梯度结论的精确范围}
记 $V$ 为条件速度标签，$U=(Y_\tau,\tau)$，$m(U)=\E[V\mid U]$。由投影恒等式，
\begin{equation}
 L_{\mathrm{CFM}}(\theta)-L_{\mathrm{FM}}(\theta)
 =\E\norm{V-m(U)}^2.
\end{equation}
如果配对分布与概率路径固定，不随 $\theta$ 改变，右边为常数；在可交换求导与积分的条件下，两者梯度相同。若用当前模型生成配对，再让配对的采样过程也参与反向传播，右边一般会随参数变化。reflow 等过程常把新配对视为固定训练数据，再训练下一轮模型；不能忽略“路径是否固定”这个条件。

类似地，可更换时间抽样或正权重，但必须保证比较的两个损失使用相同权重。模型容量不足时，“每个位置的理论最优相同”也不意味着两种有限网络训练具有完全相同的优化轨迹。

\section{热方程与熵：加噪的一个背景直觉}
纯布朗加噪 $X_t=X_0+W_t$ 的密度满足 $\partial_tp_t=\Delta p_t/2$。在密度、尾部和导数足以支持以下运算时，微分熵 $H(p_t)=-\int p_t\log p_t$ 的变化为
\begin{align}
 \frac{\dd H}{\dd t}
 &=-\tfrac12\int(1+\log p_t)\Delta p_t\dd x\\
 &=\tfrac12\int\frac{\norm{\grad p_t}^2}{p_t}\dd x
 =\tfrac12\E_{p_t}\norm{s_t(X_t)}^2\ge0.
\end{align}
右边称为与平移有关的 Fisher 信息的一种形式。纯加性噪声使分布变得更不集中，而 score 的平方平均描述此熵增长速率。对于带漂移的 VP 过程，熵导数还有漂移散度贡献；不能将这条单调性无条件移植到所有扩散过程。

}{}
\IfFileExists{appendices/03-glossary.tex}{\chapter{术语词典与回读入口}\label{app:glossary}
术语按理解依赖组织，中文名称后给出常见英文写法。正文中的青绿色术语会跳转到这里。每组末尾提供回到相关正文的链接；PDF 阅读器的“上一视图”也可直接返回原位置。

\section{分布、几何与统计}
\subsection*{分布（distribution）}\label{gl:distribution}
随机变量落在各集合中的概率规律。对图像而言，分布既约束颜色、纹理，也约束形状、空间关系及它们的组合。学习分布不是把每张训练图都编码成一个文件；生成只要求新随机样本具有相应统计规律。回读：第\ref{ch:concepts}章、附录\ref{app:probability}。

\subsection*{密度（probability density）}\label{gl:density}
将区域概率写为体积积分时的被积函数。密度值不是单点概率，也可能大于 $1$。有些分布，例如严格集中在低维集合上的分布，没有相对于全空间体积的密度。回读：附录\ref{app:probability}。

\subsection*{条件与条件分布（conditioning）}\label{gl:conditional}
在已知某些信息之后，对其余随机对象的概率规律进行更新。例如 $p(x\mid c)$ 表示给定文本或类别后的图像分布。这里的“条件”不一定是训练代码里拼接的一个向量；如何编码和注入它，是网络设计问题。回读：第\ref{sec:guidance}节。

\subsection*{流形（manifold）}\label{gl:manifold}
在局部可以用较少坐标描述的几何空间。自然图像位于“数据流形附近”是解释高维结构的有用假设，而不是已经证明所有真实图像都严格落在某个光滑低维流形上。噪声会使原本集中的分布铺开，从而得到普通密度与可学习的 score。回读：第\ref{ch:concepts}章。

\subsection*{隐变量与隐空间（latent variable / latent space）}\label{gl:latent}
概率模型中的隐变量是未被直接观测、却参与生成或推断的变量；随机生成起点、DDPM 的辅助含噪轨迹、编码器得到的压缩表示，是不同的具体用法。隐空间指这些变量所在的表示空间，在隐空间扩散中通常特指自编码器的内部表示。它可以比像素空间维度低，但低维不自动意味着每一维都可解释为某个语义。隐空间扩散学习编码后分布，最后用解码器恢复可见数据。回读：第\ref{ch:concepts}章、第\ref{ch:applications}章。

\subsection*{似然（likelihood）}\label{gl:likelihood}
已观测样本作为固定输入，密度或概率质量作为参数的函数。它衡量模型给数据赋予的统计支持，不直接等于审美质量或提示词遵循。连续图像的似然报告还要明确去量化、归一化和维度。回读：第\ref{sec:ddpm-elbo}节、式\eqref{eq:cnf-likelihood}。

\subsection*{KL 散度（Kullback--Leibler divergence）}\label{gl:kl}
$\KL(q\Vert p)=\E_q\log(q/p)$，比较两个概率规律。它非负、不对称，支撑不匹配时可能无穷。优化哪个方向的 KL 会影响对漏掉模式和分配错误概率的惩罚。回读：附录\ref{app:probability}。

\subsection*{ELBO（evidence lower bound）}\label{gl:elbo}
对隐藏变量模型的对数似然给出的可优化下界。下界差距为辅助后验与模型后验的 KL。DDPM 用完整前向含噪轨迹作为辅助后验；简单噪声损失是对原变分目标重新加权后的训练选择。回读：第\ref{sec:ddpm-elbo}节。

\subsection*{熵（entropy）}\label{gl:entropy}
度量分布不确定性或平均信息量的一种统计量。连续微分熵依赖坐标尺度，不同于离散熵。纯加性高斯平滑在适当条件下使微分熵增大；带漂移的扩散不能据此无条件断言熵单调增加。回读：附录\ref{app:proofs}。

\subsection*{信噪比（signal-to-noise ratio, SNR）}\label{gl:snr}
在 $x_t=ax_0+b\epsilon$ 的日程表示里常指 $a^2/b^2$。它表示信号与噪声的系数尺度比；数据方差不均匀时，不同方向的实际信息保留量仍不同。log-SNR 常用于时间参数化与损失权重分析。回读：第\ref{sec:parameterization}节。

\subsection*{重参数化（reparameterization）}\label{gl:reparameterization}
将复杂随机变量写成简单噪声的可微确定性变换。例如直接构造 $x_k=\sqrt{\bar\alpha_k}x_0+\sqrt{1-\bar\alpha_k}\epsilon$。它使训练能随机抽取时刻，而不必模拟完整前向链。回读：式\eqref{eq:forward-marginal}。

\section{动力学、输运与训练}
\subsection*{Markov 性（Markov property）}\label{gl:markov}
给定当前状态后，未来条件分布不再依赖更早历史。DDPM 前向链具有这一性质；DDIM 的辅助联合构造可以改变多时刻依赖，而保留训练所需的单时刻边缘。回读：第\ref{ch:methods}章。

\subsection*{Score（状态的对数密度梯度）}\label{gl:score}
$s_t(x)=\nabla_x\log p_t(x)$。它是状态空间的向量场，区别于对参数求导的统计 score。高斯条件核的 score 可计算，条件均值则给出未知边缘 score。它不等于图像质量函数。回读：第\ref{sec:score}节。

\subsection*{ODE（ordinary differential equation）}\label{gl:ode}
普通微分方程，用速度场和初始位置确定连续轨迹。给定噪声起点后，生成轨迹可以确定；随机起点仍使最终输出具有多样性。回读：第\ref{sec:pfode}节、附录\ref{app:stochastic}。

\subsection*{SDE（stochastic differential equation）}\label{gl:sde}
随机微分方程，在漂移之外还由布朗增量驱动。它是积分方程的简写，不能把噪声项视为普通可微函数。扩散生成使用正确的时间反向 SDE，逆漂移包含边缘 score。回读：第\ref{sec:sde}节。

\subsection*{布朗运动（Brownian motion）}\label{gl:brownian}
连续的随机过程，独立增量的方差与时间间隔成正比。它的小步典型尺度是时间间隔的平方根，路径几乎必然不可按普通方式求导。回读：附录\ref{app:stochastic}。

\subsection*{It\^o 公式（It\^o formula）}\label{gl:ito}
随机过程上的链式法则。与普通链式法则相比，多了由布朗二次变差产生的二阶项。正是这个项让 SDE 的密度方程包含扩散与 $1/2$ 系数。回读：第\ref{sec:sde}节、附录\ref{app:stochastic}。

\subsection*{Fokker--Planck 方程}\label{gl:fp}
描述 SDE 所诱导密度随时间变化的偏微分方程。对于状态无关标量扩散，$\partial_tp=-\nabla\cdot(fp)+g^2\Delta p/2$。它把粒子的随机运动转成整体密度演化。回读：式\eqref{eq:fokker-planck}。

\subsection*{连续性方程（continuity equation）}\label{gl:continuity}
表达输运中的概率质量守恒：$\partial_tp+\nabla\cdot(pu)=0$。局部概率减少对应净流出，不是概率消失。概率流 ODE 与流匹配都使用这一关系。回读：第\ref{sec:pfode}、\ref{sec:flowmatching}节。

\subsection*{散度（divergence）}\label{gl:divergence}
向量场各分量对对应坐标偏导之和，$\nabla\cdot v=\sum_i\partial_iv_i$。它描述局部体积膨胀率；沿 ODE 轨迹的对数密度变化为负散度。概率流密度的散度 $\nabla\cdot(pv)$ 与速度散度不同。回读：式\eqref{eq:cnf-likelihood}。

\subsection*{Jacobian（雅可比矩阵）}\label{gl:jacobian}
向量函数的一阶偏导矩阵。其行列式描述局部体积缩放，迹在同维速度场中给出散度。连续归一化流将离散变换的行列式计算转化为沿时间积分散度。回读：附录\ref{app:proofs}。

\subsection*{Lipschitz 条件}\label{gl:lipschitz}
用 $\norm{v(x)-v(y)}\le L\norm{x-y}$ 控制函数对输入变化的敏感性。局部 Lipschitz 有助于 ODE 解唯一；与增长条件配合可防止有限时间爆炸。神经场的实际 Lipschitz 行为也影响误差放大。回读：附录\ref{app:stochastic}。

\subsection*{刚性（stiffness）}\label{gl:stiffness}
动力学中不同变化尺度相差很大，使某些显式数值方法为了稳定而不得不用很小步长的现象。高阶不自动解决刚性，也不自动消除网络误差。回读：附录\ref{app:stochastic}、第\ref{ch:methods}章。

\subsection*{耦合（coupling）}\label{gl:coupling}
具有给定两个边缘分布的联合分布。耦合决定噪声与数据如何配对；改变配对可能改变流匹配的中间路径和学习难度，而端点边缘保持相同。回读：第\ref{sec:flowmatching}节。

\subsection*{概率输运与推前（transport / pushforward）}\label{gl:transport}
用映射或动力学把一份概率分布搬成另一份。推前记作 $T_\#p$；最优输运另有指定代价的最小化要求，不能把任意直线插值都称为全局最优。回读：附录\ref{app:probability}。

\subsection*{桥（bridge）}\label{gl:bridge}
带有端点或端点分布约束的路径建模。不同论文的桥可能指条件随机桥、Schr\"odinger bridge 或其他输运构造，不能只因“连接两端”就视为同一种算法。阅读时应检查参考过程、耦合、路径随机性与优化目标。回读：第\ref{ch:methods}章。

\subsection*{CFG（classifier-free guidance）}\label{gl:cfg}
组合条件预测和空条件预测来增强条件影响。训练常用条件丢弃获得两种预测。引导改变生成场，并可能改变多样性、稳定性和条件遵循；引导尺度的数值必须连同公式约定比较。回读：第\ref{sec:guidance}节。

\subsection*{蒸馏（distillation）}\label{gl:distillation}
利用较强或较慢教师的输出、轨迹或预测训练更便宜的学生。少步生成可以通过蒸馏得到，但“减少采样步数”也可由求解器改进实现；两者的训练成本与误差来源不同。回读：第\ref{ch:methods}章。

\section{生成模型与网络组件}
\subsection*{VAE（variational autoencoder）}\label{gl:vae}
用隐变量、编码器近似后验和解码器构成生成模型，并优化变分下界。隐空间扩散的自编码器是表示工具，扩散先验则建模编码后的分布；两部分职责不同。回读：第\ref{ch:concepts}章。

\subsection*{GAN（generative adversarial network）}\label{gl:gan}
用生成器与判别器的对抗训练学习分布匹配。通常采样很快，直接密度求值却不容易；稳定性与模式覆盖依赖具体目标和训练方法。它同样从随机噪声生成图像，不能把“用噪声作起点”视为扩散独有性质。回读：第\ref{ch:concepts}章。

\subsection*{归一化流（normalizing flow）}\label{gl:flow}
利用可逆变换及变化量公式构造可计算密度的模型。连续流以 ODE 表示变换，沿轨迹积分散度计算密度。Flow Matching 是训练连续流的一类方式，而非与所有 flow 模型无关的名称。回读：第\ref{ch:concepts}章、第\ref{sec:flowmatching}节。

\subsection*{自回归（autoregressive, AR）}\label{gl:ar}
把联合概率分解为有序条件概率乘积，例如 $p(x)=\prod_ip(x_i\mid x_{<i})$。可以用于像素、离散视觉 token 或文本。自回归规定因子分解，Transformer 是网络结构；两者不是同义词。回读：第\ref{ch:concepts}章。

\subsection*{注意力（attention）}\label{gl:attention}
通过 query 与 key 的相关性计算权重，再对 value 汇总信息。自注意力让同一组位置交换信息；交叉注意力让一组位置读取另一组信息，例如图像读取文本。注意力的作用来自信息路由，不等于模型已经学会真实物理因果关系。回读：第\ref{ch:applications}章。

\subsection*{Token（处理单元）}\label{gl:token}
网络输入序列中的一个单元。文本 token 常为词片段或离散符号；视觉 token 可以是连续 patch 向量，并不一定是离散词表索引。多模态系统把不同模态组织成 token 后，仍要设计它们的空间位置、时间位置及注意力约束。回读：第\ref{ch:applications}章。

\subsection*{Patch（局部块）}\label{gl:patch}
把图像或隐空间网格切成局部小块，再映射成序列向量。patch 大小影响 token 数、计算成本和细节表达。DiT 可以在连续隐空间 patch 上处理扩散或流速度预测，不能仅因 patch 化就称其为离散扩散。回读：第\ref{ch:applications}章。

\subsection*{感知损失（perceptual loss）}\label{gl:perceptual}
在预训练特征空间里比较图像，而不只比较像素。它能更关注某些可见结构，但会继承特征提取器的偏好；低感知损失不代表所有信息无损保存。隐空间自编码器的压缩质量决定生成系统能够恢复的细节上限。回读：第\ref{ch:applications}章。

\subsection*{等变性（equivariance）}\label{gl:equivariance}
输入经过某种变换后，输出按相应方式变化，例如旋转分子后预测坐标也随之旋转。它与不变性不同：不变输出完全不变。科学与三维生成常需要这样的结构约束来提高数据效率与物理一致性。回读：第\ref{ch:applications}章。

\subsection*{FID（Fr\'echet Inception Distance）}\label{gl:fid}
在某个特征空间中，以均值与协方差的高斯近似比较真实和生成样本集。它是集合统计指标，受特征提取器、样本量与评估协议影响；不能直接衡量某张图的质量，也不能单独衡量文本准确性、物理合理性或安全性。回读：第\ref{ch:applications}章的评价讨论。
}{}
\IfFileExists{appendices/04-reading.tex}{\chapter{阅读路线与可复现实验}\label{app:exercises}
\section{三种阅读节奏}
\textbf{先建立直觉。}读第一章和第二章的双峰例子、条件投影框、标准高斯的 SDE/ODE 对照，再读第三章的四轴分类，最后进入第四章的网络和应用。这条路线先获得各对象的定位，随后补推导。

\textbf{以数学为主。}按第二章顺序阅读，做五道章末练习，再补附录 A--C。推导中遇到障碍时，先检查条件和时间方向；不要先试图背最终公式。最后用第三章检查自己能否把公式变成可运行算法。

\textbf{以系统为主。}第一章后进入第四章的 U-Net、隐空间、DiT 与多模态信息流，再回第三章辨认实际系统的路径、目标、求解器和网络，最后回第二章核验公式。仅看到一个网络结构图，不足以确定模型使用哪一种生成动力学。

\section{六个无需大模型的实验}
这些实验是进一步学习建议，并不是本书声称已经训练完成的模型结果。它们可以用一维或二维数据完成，让你把数学关系看成数值事实。

\begin{enumerate}
\item \textbf{DDPM 边缘核验。}对固定 $x_0$，一组样本逐步加噪，另一组直接用式\eqref{eq:forward-marginal}生成。比较单时刻均值、方差与直方图，再比较跨时刻协方差。你会看见边缘一致而联合构造可不同。
\item \textbf{双峰 score。}使用第二章解析混合分布，在不同噪声尺度上绘制密度和 score。找出 score 零点，区分密度极大值与极小值。随后用精确场生成，避免把网络训练误差和动力学原理混在一起。
\item \textbf{SDE/ODE 同边缘。}从标准高斯出发模拟 VP SDE，并比较零速度概率流 ODE。画逐时刻直方图与单条轨迹：前者接近相同，后者明显不同。
\item \textbf{高斯 Flow Matching。}用练习 2.4 的解析速度求解 ODE，比较 Euler、Heun 与解析解。对步数作误差曲线，检查阶数成立的区间。
\item \textbf{参数化权重。}在相同批次与日程上计算噪声损失、原图损失及转换后的权重。逐时刻核验代数关系，观察未经权重修正的损失如何关注不同尺度。
\item \textbf{条件投影。}构造两组在中间位置相交的二维插值，统计相同小邻域里的速度标签和均值。把个别条件直线与边缘速度积分轨迹一起画出，检查它们为何不同。
\end{enumerate}

\section{阅读论文时用同一张问题清单}
每读一篇新论文，先记录：数据表示是什么；端点与概率路径是什么；网络预测哪个目标；损失权重怎样设计；条件怎样注入；采样如何积分；是否另有蒸馏阶段；作者的实验控制了哪些因素。随后再读结果表。

当两篇论文给出不同的结论时，先检查数据、表示、训练预算、引导、评估协议与推理成本是否相同。不要把更大的模型或更多数据带来的收益全归因于一个新公式。第四章的精选阅读路线提供了入口；参考文献中的链接指向论文或官方披露来源。

\section{使用与编译这份书稿}
本书主文件是 \path{main.tex}，正文与附录分别存于独立文件，模板样式在 \path{book-style.tex}。推荐 XeLaTeX 与 Biber。完整顺序为 XeLaTeX、Biber、XeLaTeX、XeLaTeX；最后两遍建立引用和目录。Windows 本地可运行附带的 \path{scripts/build.ps1}。

项目保留 ElegantBook 原始类和许可证，私有依赖只是补齐本地环境所缺的公开组件。在完整 TeX Live 或 Overleaf 环境中通常无需这些额外依赖。具体编译与质量检查记录保存在项目说明中。修改章节后应重新构建整本书，检查未定义引用、长公式和页面溢出，避免只更新源码而仍阅读旧 PDF。
}{}
\backmatter
\printbibliography[heading=bibintoc,title={参考文献与原始论文}]
\end{document}
